\chapter{Reading the paper: a guided tour}\label{ch:paper}
\chaptermark{Reading the paper}

This chapter reads the research note ``Local factors of the adjoint $L$-function and the $\ell$-adic valuations of the modular degree'' section by
section, in the paper's own order. It is written last, from the finished book. For each section and each numbered result of the paper the tour does three
things: it states the result (in the paper's words when they are quoted, in the book's notation otherwise); it says in which chapter of the book each
ingredient is proved or explained; and it expands the steps that the paper compresses, with short guidance on the hard ones. The chapter ends with a map
of the paper (dependencies, hypotheses, imported statements) and a table of worked verifications.

\begin{sloppypar}
\section*{Core results at a glance}

The paper's results are listed in the paper's order, with the paper's labels. Throughout, $E/\Q$ has conductor $N$, $f=f_E$, $\ell$ is an odd prime,
$T=T_\ell E$, $V=V_\ell E$, $A_f=\ad V$, $W=\ad T\otimes\Qlmod$, and $\Frob_p^{-1}$ is the geometric Frobenius.

\begin{enumerate}[label=\textbf{\thechapter.\arabic*}]
\item \textbf{Conventions.} Standing notation; $\Frob_p$ arithmetic, $L_p(M,s)=\det(1-\Frob_p^{-1}p^{-s}\mid M^{I_p})^{-1}$, $L_p(\Sym f,s)=L_p(A_f,s-1)$, $L_q(\Sym f,s)=P_q(q^{-s})^{-1}$; duals
$M^\vee$, $X^*$, $M^D$, and $T^D\simeq A(1)$ for $\ell$ odd. \emph{(convention)} \Cref{conv:paper:standing,conv:paper:frobenius,conv:paper:duals}.
\item \textbf{Definition def:terms.} $t_\ell(E,q)=v_\ell(-v_q(j))$ if $v_q(j)<0$, else $0$ (every bad $q$); $e_\ell(E,q)=v_\ell\bigl(P_q(q^{-2})/(1-a_q^2q^{-2})\bigr)$ ($q\ne\ell$), a nonnegative integer, $0$ at
multiplicative $q$; $t_\ell$ is twist invariant, $e_\ell(E,q)$ invariant under twists unramified at $q$ only (\curve{11a1}, \curve{121d2}). \emph{(definition; properties proved here.)} \Cref{def:paper:terms}.
\item \textbf{Tamagawa factor} (definition). $\Tam_p(M)=\Fitt_{\Zl}((H^1(I_p,M)_{\mathrm{tor}})^{G_{\Q_p}})$; $H^1_f(\Q,\ad T\otimes\Qlmod)$ the Bloch-Kato Selmer group. \emph{(definition.)} \Cref{def:paper:tamagawa}.
\item \textbf{Theorem thm:main.} $E$ optimal, $\ell\nmid2N$ odd, $\ell\notin S_f$: $v_\ell(\degphi)=\len H^1_f(\Q,W)+\sum_{q\mid N}v_\ell(\Tam_q(\ad T))+\sum_{q\mid N}e_\ell(E,q)$, hence
$v_\ell(\degphi)\ge\sum_q(e_\ell(E,q)+t_\ell(E,q))$. \emph{(proved here, given the theorems of Diamond, Flach and Guo, \cite[Theorem 2.2]{ARS12} and Proposition prop:integral, which is proved here.)}
\Cref{thm:paper:main}.
\item \textbf{Conjecture conj:euler.} For $E$ optimal, non-CM, $E[\ell]$ irreducible: $v_\ell(\degphi)\ge\sum_{q\ne\ell}e_\ell+\sum_qt_\ell+[\ell=3]w_3(E)-h_\ell(E)$; $h_\ell=0$ for $\ell\ge5$, $h_3\in\{0,1\}$ with $h_3=1$ iff
image $\Ns$. \emph{(stated; the facts on $h_\ell$ proved here.)} \Cref{conj:paper:euler}.
\item \textbf{Component groups and weights} (definition): $\#\Phig_q$ by type; $w_\ell(E,q)$; (eq:main), (eq:twist); $w_\ell=\max(v_\ell\#\Phig_q(E),v_\ell\#\Phig_q(E^{(q)}))$ for odd $q$. \Cref{def:paper:weights}.
\item \textbf{Proposition prop:forms.} Under the hypothesis of Theorem thm:main, (eq:main) and (eq:twist) hold; likewise under (eq:rule) with $h_\ell=0$. \emph{(proved here,
given \Cref{thm:paper:main,lem:paper:IV}.)} \Cref{prop:paper:forms}.
\item \textbf{Conjecture conj:eis.} Eisenstein deficit $D\le v_\ell(\#E(\Q)_{\mathrm{tors}})+[E[3]\simeq\Z/3\oplus\mu_3]$ and $\le$ the number of curves of the class with a rational $\ell$-torsion point;
$D\le2$ follows from the first bound. \emph{(stated.)} \Cref{conj:paper:eis}.
\item \textbf{Lemma lem:tamagawa-exp.} At multiplicative $q\ne\ell$: $\rho\bmod\ell^t$ unramified at $q$ iff $\ell^t\mid n_q$. \emph{(proved here, given the Tate curve, \Cref{ch:reduction}.)}
\Cref{lem:paper:tamagawa-exp}.
\item \textbf{Inertia invariants.} At additive $q\ne\ell$: $E[\ell^k]^{I_q}\simeq\Phig_q[\ell^k]$. \emph{(proved here, given the N\'eron filtration, \Cref{ch:reduction}.)} \Cref{lem:paper:inertia-invariants}.
\item \textbf{Lemma lem:IV.} At $q\ne3$ of type $\kod{IV}$, $\kod{IV}^*$: potentially good, $\rho(I_q)$ cyclic of order $3$, $E[3]^{I_q}$ a line, Serre exponent $1$, $f_q=2$, $P_q=1-\chi_{-3}(q)qX$, $e_\ell=v_\ell(q-\chi_{-3}(q))$; at
$q=2$, $P_2=1+2X$, $e_3=1$. \emph{(proved here, given \cite{SerreTate68} and \Cref{ch:ant,ch:local}.)} \Cref{lem:paper:IV}.
\item \textbf{Lemma lem:twist.} (i) Same type, $\Phig_p$, $f_p$ at $p\ne q$; (ii) $\kod{I}_n\leftrightarrow\kod{I}_n^*$ ($n\ge0$), $\kod{II}\leftrightarrow\kod{IV}^*$, $\kod{III}\leftrightarrow\kod{III}^*$, $\kod{IV}\leftrightarrow\kod{II}^*$ at every odd $q$, with $f_3$
unchanged by the last three at $q=3$; (iii) Atkin-Li. \emph{((i), (ii) proved here, given Tate's algorithm and, at $q=3$, the Ogg-Saito formula; (iii) stated, reference \cite{AtkinLi78}.)}
\Cref{lem:paper:twist}.
\item \textbf{Lemma lem:euler.} The table: $L_q(A_f,s)^{-1}$ and $e_\ell$ for $\kod{I}_n$ (every $q\ne\ell$), $\kod{I}_n^*$, $\kod{I}_0^*$ (odd $q$), $\kod{II},\kod{II}^*,\kod{IV},\kod{IV}^*$ and $\kod{III},\kod{III}^*$ ($q\ge5$). \emph{(proved here,
given the inertia images of \Cref{ch:galrep}.)} \Cref{lem:paper:euler}.
\item \textbf{Proposition prop:shapes.} At every bad $q\ne\ell$, $P_q$ is $1-X$, the $\kod{I}_0^*$ cubic $(1-qX)(1-(a^2-2q)X+q^2X^2)$, $1-cqX$ ($c=\pm1$), or $1$ (nonabelian inertia, $q=2,3$ only).
\emph{(proved here, given \Cref{ch:galrep,ch:lfunctions}.)} \Cref{prop:paper:shapes}.
\item \textbf{Lemma lem:tamfactor.} At potentially multiplicative $q\ne\ell$ ($q=2$ included): $v_\ell(\Tam_q(\ad T))=v_\ell(n)+\min(v_\ell(n),v_\ell(q-1),v_\ell(c(u)))\ge t_\ell(E,q)$. \emph{(proved here, given
\Cref{ch:cohomology,ch:galcoh}.)} \Cref{lem:paper:tamfactor}.
\item \textbf{Lemma lem:tamadd.} At potentially good additive $q\ne\ell$: $H^1(I_q,\ad T)_{\mathrm{tor}}=0$, so $v_\ell(\Tam_q(\ad T))=0$. \emph{(proved here, given \Cref{ch:cohomology,ch:ant}.)}
\Cref{lem:paper:tamadd}.
\item \textbf{Congruence number} (definition) and Agashe-Ribet-Stein: $\degphi\mid r_E$, equal $p$-parts for $p^2\nmid N$. \emph{(the theorem stated; reference \cite[Theorem 2.2]{ARS12}.)}
\Cref{def:paper:congruence-number}.
\item \textbf{Lemma lem:crit.} $\T/(\ker\phi_f+\Ann(f^\perp))\simeq\Z/r_E$; a congruent eigenform orthogonal to $f$ gives $\ell\mid r_E$. \emph{(proved in \Cref{ch:modgalrep}; walked through here.)}
\Cref{lem:paper:crit}.
\item \textbf{Lemma lem:old.} A congruent newform $h\ne f$ of level $M\mid N$, with $p\nmid M$ at the $p\,\|\,N$ where $\bar\rho$ is unramified, gives $\ell\mid r_E$. \emph{(proved in \Cref{ch:modgalrep}; walked through
here.)} \Cref{lem:paper:old}.
\item \textbf{Optimal-level theorem} (Ribet, Carayol, Diamond; \cite[Theorem 3.15]{DDT}). \emph{(stated; reference.)} \Cref{thm:paper:optimal-level}.
\item \textbf{Theorem thm:qual.} $\ell\nmid N$ odd, $\bar\rho$ irreducible (absolutely over $\Q(\sqrt{-3})$ if $\ell=3$), and a prime of type $\kod{IV},\kod{IV}^*$ ($q\ge5$ or $q=2$, $\ell=3$), $\kod{II},\kod{II}^*$ ($q\ge5$, $\ell=3$) or
$\kod{I}_n^*$ ($n\ge1$, $\ell\mid n$, $q$ odd): then $\ell\mid\deg\varphi_{E_0}$. \emph{(proved here, given \Cref{thm:paper:optimal-level} and \cite{ARS12}.)} \Cref{thm:paper:qual}.
\item \textbf{Objects of Diamond, Flach and Guo} (definition): $M^\Sigma$, $f^\Sigma$, the integral structure, $M_{f^\Sigma,\lambda}$, $\hat\delta_!$, $\eta^\Sigma_{f,\lambda}$. \Cref{def:paper:dfg-objects}.
\item \textbf{Proposition prop:integral.} For every prime $\ell$, pull-back identifies $H^1(X_0(M^\Sigma),\Z)$ with the Betti lattice $M(M^\Sigma,1)_{!,B}$, Hecke equivariantly, $U_p$ included; so
$M(M^\Sigma,1)_{!,\lambda}=H^1(X_0(M^\Sigma),\Zl)$, with no localisation. \emph{(proved here; expanded proof in \Cref{ch:congruence}.)} \Cref{prop:paper:integral}.
\item \textbf{Theorem thm:dfg14} (Diamond, Flach and Guo, Proposition 1.4(c)). For $\ell\nmid2M$, $\ell\notin\Sigma$, $\bar\rho_f$ irreducible: $\eta^\Sigma=\eta^\emptyset\prod_{p\in\Sigma}L^{\nv}_p(A_f,1)^{-1}$. \emph{(stated; reference.)}
\Cref{thm:paper:dfg14}.
\item \textbf{Proposition prop:pullback.} For optimal $E$: $\varphi^*H^1(E,\Zl)=H^1(X_0(N),\Zl)[I_f]$, $\langle\varphi^*a,w_N^*\varphi^*b\rangle=\varepsilon_N\degphi\langle a,b\rangle$; for $\ell\nmid2N$,
$v_\lambda(\eta^\emptyset_{f,\lambda})=v_\ell(\degphi)$. \emph{(proved here, given \Cref{ch:moddeg} and Proposition prop:integral.)} \Cref{prop:paper:pullback}.
\item \textbf{Lemma lem:ddt417.} $d\,v(\eta)\ge v(\Delta)$, with equality if $L$ is free over $\T$. \emph{(proved here, given \Cref{ch:modules}.)} \Cref{lem:paper:ddt417}.
\item \textbf{Lemma lem:hecke.} $v_\ell(\degphi)=v_\ell(r_E)=v_\lambda(\eta^\emptyset_{f,\lambda})$ for $E$ optimal, $\ell\nmid2N$, $\bar\rho$ irreducible; no freeness. \emph{(proved here, given \cite{ARS12} and
\Cref{prop:paper:pullback}.)} \Cref{lem:paper:hecke}.
\item \textbf{Selmer set-up} (definition): $W$, $T^D$, $V^D$, $W^D$, $H^1_\Sigma$, $S_f$. \Cref{def:paper:selmer-setup}.
\item \textbf{Proposition prop:Sf.} For $\ell\nmid N$ odd and $E[\ell]$ irreducible: $\ell\notin S_f$ if $\ell\ge5$; $3\in S_f$ iff image $N_s(\F_3)$ (then ordinary at $3$); $\ell\notin S_f$ iff $h_\ell=0$.
\emph{(proved here, given \cite[Section 1]{Serre72} and \Cref{ch:reps}.)} \Cref{prop:paper:Sf}.
\item \textbf{Theorems thm:dfg37 and thm:dfg27} (Diamond, Flach and Guo, Theorems 3.7 and 2.7). For $\ell\notin S_f$: $\len H^1_\Sigma(\Q,W)=v_\lambda(\eta^\Sigma)$; $H^0(\Q,V^D)=H^1_f(\Q,V^D)=0$.
\emph{(stated; reference; they cover $\ell=3$.)} \Cref{thm:paper:dfg37,thm:paper:dfg27}.
\item \textbf{Lemma lem:dfg21.} For $\ell\notin\Sigma$: $0\to H^1_f\to H^1_\Sigma\to\bigoplus H^1_f(\Q_p,T^D)^*\to H^0(\Q,W^D)^*\to0$; $H^0(\Q,V^D)=0$ automatic; $H^1_f(\Q_p,T^D)$ the torsion of $H^1(\Q_p,T^D)$.
\emph{(stated; reference; the paragraph after it proved here.)} \Cref{lem:paper:dfg21}.
\item \textbf{Lemma lem:fpr.} $\Fitt H^1_f(\Q_p,T^D)=L_p(A_f,1)^{-1}\Tam_p(T^D)$ and $\Tam_p(T^D)=\Tam_p(A)$. \emph{(first identity proved here; second stated, reference \cite[(57)]{DFG04}, and checked for the
elliptic adjoint.)} \Cref{lem:paper:fpr}.
\item \textbf{Lemma lem:h0.} $H^0(\Q,W^D)$ finite cyclic of length $h_\ell(E)$; $h_\ell=0$ for $\ell\ge5$; $h_3\le1$, with equality iff $\bar\rho\simeq\bar\rho\otimes\chi_{-3}$ iff image $N_s(\F_3)$; CM allowed.
\emph{(proved here, given \Cref{ch:reps,ch:ant}.)} \Cref{lem:paper:h0}.
\item \textbf{Corollary cor:explicit.} $v_\ell(\degphi)=\len H^1_f+\sum e_\ell+\sum_{v_q(j)<0}(t_\ell+\min(t_\ell,v_\ell(q-1),v_\ell(c(u_q))))$. \emph{(proved here, given \Cref{thm:paper:main}.)}
\Cref{cor:paper:explicit}.
\item \textbf{The degree formula} (Remark rem:degree). $\degphi=c_E^2N\,L^{\nv}(\ad f,1)/(2\pi\vol(\Lambda_E))=B_E\prod_{q^2\mid N}L_q(A_f,1)^{-1}$; $v_\ell(\degphi)=v_\ell(B_E)+\sum e_\ell$ for $\ell\nmid N$.
\emph{(proved here, given the formula of Zagier and Watkins and the Petersson formula.)} \Cref{prop:paper:degree}.
\item \textbf{Theorem thm:twist.} $r_{E_S}\mid r_E$ if $q^2\mid N$ ($q\in S$) and $\lcm(N_S,D^2)\mid N$; equality if also $\lcm(N,D^2)\mid N_S$. \emph{(proved here, given \Cref{ch:hecke}.)}
\Cref{thm:paper:twist}.
\item \textbf{Lemma lem:step1.} For every $E$, with $S$ its odd $\kod{I}_0^*$, $\kod{I}_n^*$ primes: $v_\ell(r_E)\le v_\ell(\eta^{\mathrm{Hecke}}(g^S))$ at every prime $\ell$, with equality at every $\ell\notin S$.
\emph{(proved here and in \Cref{ch:congruence}.)} \Cref{lem:paper:step1}.
\item \textbf{Theorem thm:quant.} For $\ell\nmid2N$, $E[\ell]$ irreducible, $S$ the odd $\kod{I}_0^*$, $\kod{I}_n^*$ primes: $v_\ell(\degphi)\ge v_\ell(r_{E_S})+\sum_Se_\ell$, with equality if $\ell\notin S_f$.
\emph{(proved here, given \Cref{thm:paper:dfg14,thm:paper:main}; also in \Cref{ch:congruence}.)} \Cref{thm:paper:quant}.
\item \textbf{Proposition prop:ns.} $E[3]$ irreducible: image $\GL_2(\F_3)$, $N_{ns}(\F_3)$ or $N_s(\F_3)$; absolutely irreducible over $\Q(\sqrt{-3})$ unless $N_s$; then $\bar\rho\otimes\F_9=\Ind\psi$, $\bar\rho\simeq\bar\rho\otimes\chi_{-3}$,
$h_3=1$; $\psi_3$ is two quadratics or irreducible. \emph{(proved here, given \Cref{ch:groups,ch:reps}.)} \Cref{prop:paper:ns}.
\item \textbf{Proposition prop:ns-structure.} For image $N_s(\F_3)$: (i) $\bar\rho(I_p)\subseteq C_4$ ($p\ne3$), so $t_3\ge1$ at potentially multiplicative $p\ne3$; (ii) the order-four primes are the $p\ne3$ ramified in
$K_1$, a nonempty set, with $e_3=v_3(p-\chi_{-3}(p))\ge1$ there; (iii) at $p\ge5$ they are the primes of type $\kod{III}$, $\kod{III}^*$, and $p\equiv\pm1\pmod{12}$; (iv) biquadratic splitting field; (v)
$3\mid a_p$ for $p\equiv2\pmod3$, $p\nmid3N$; (vi) ordinary at $3$ if $3\nmid N$; (vii) $h_3=1$. \emph{(proved here, given Kraus's theorem and \Cref{prop:paper:Sf}.)} \Cref{prop:paper:ns-structure}.
\item \textbf{Remark rem:ns} (conditional proposition). For image $\Ns$, if $H^1_f(\Q,V^D)=0$: $\len H^1_\Sigma=\len H^1_f+\sum_p(v(L_p^{-1})+v(\Tam_p))-1$; with $v(\eta^\Sigma)=\len H^1_\Sigma$ this gives (eq:identity)
minus $h_3$. \emph{(proved here, conditionally, given \Cref{lem:paper:dfg21}.)} \Cref{prop:paper:ns-conditional}.
\item \textbf{Proposition prop:no9.} A non-CM $E/\Q$ with $E[3]\simeq\Z/3\oplus\mu_3$ has no rational point of order $9$; so $D\le2$ in Conjecture conj:eis follows from its first bound. \emph{(proved here,
given \cite{Ligozat75} and \cite{Mazur77}.)} \Cref{prop:paper:no9}.
\item \textbf{The map and the verifications.} Dependencies, hypotheses and imported statements (Section~\ref{sec:paper:map}); recomputed curve data (Section~\ref{sec:paper:verifications}).
\end{enumerate}

\section*{Where the paper uses this chapter}

All of it: this chapter is the paper, read with the book open beside it. Each numbered section below covers one section of the paper (or a part of
one) and opens by saying which. Sections~\ref{sec:paper:map} and~\ref{sec:paper:verifications} are devices of the tour: the dependency map and the table of
hypotheses (used to see which inputs each statement needs), and the worked verifications (used to check every number quoted from a curve). A reader who
wants only the proofs can read \Cref{ch:congruence} for Section 4 of the paper, \Cref{ch:galrep,ch:reduction} for Sections 2 and 5, \Cref{ch:modgalrep} for
Section 3, \Cref{ch:blochkato,ch:lfunctions} for the local terms and \Cref{ch:computation} for Section 6 and the appendix; this chapter says, result by result,
which of these to open.

\section{How to read the tour}\label{sec:paper:reading}

This section fixes the conventions the tour uses and explains how it is marked. It corresponds to the paper's Section 1.1 (conventions) and to the
opening of every section of the paper, which use the same notation.

\begin{convention}[standing notation]\label{conv:paper:standing}
$E/\Q$ is an elliptic curve of conductor $N$; unless said otherwise it is the $X_0(N)$-optimal curve of its isogeny class, $\varphi_E\colon X_0(N)\to E$ its
modular parametrisation and $\degphi$ its modular degree; $f=f_E\in S_2(\Gamma_0(N))$ is its newform and $\phi_f\colon\T\to\Z$, $T_n\mapsto a_n(f)$, the eigencharacter of
$f$ on the full Hecke algebra $\T$ (generated by all $T_n$, $n\ge1$, the $U_p$ for $p\mid N$ included). The prime $\ell$ is odd; $T=T_\ell E$, $V=V_\ell E$, $\rho=\rho_{E,\ell}$, $\bar\rho$ its
reduction on $E[\ell]$; $A_f=\ad V$ and $A=\ad T$ are the trace-zero adjoints (the macro already contains the superscript $0$); $A_f$ is never an abelian
variety. A bad prime is written $q$, any prime $p$. For $p\ne\ell$, $\delta_p=\dim V^{I_p}$ is $2,1,0$ for good, multiplicative, additive reduction. $v_\ell$ is the
$\ell$-adic valuation, extended to $\Q^\times$, to ideals of $\Zl$, and to a class $c\in\Z/\ell^m$ by $v_\ell(c)=\min(v_\ell(\tilde c),m)$ for any lift $\tilde c$; $\len$ is
$\Zl$-length. Characters: $\chi_\ell$ is the $\ell$-adic cyclotomic character and $\bar\chi_\ell$ its reduction; a numerical subscript means the quadratic
character of a quadratic field of that discriminant ($\chi_{-3}$, $\chi_{-4}$, $\chi_{-7}$, $\chi_{-11}$); for an odd prime $q$, $q^*=(-1)^{(q-1)/2}q$ and $\chi_{q^*}$ is the
quadratic character of conductor $q$; at $\ell=3$ the cyclotomic character is $\chi_3^{\mathrm{cyc}}$, and $\bar\chi_3^{\mathrm{cyc}}=\chi_{-3}$. $E^{(q)}$ is the twist of $E$ by
$\chi_{q^*}$ and $E_S=E\otimes\chi_S$ with $\chi_S=\prod_{q\in S}\chi_{q^*}$. The Tamagawa number is $c_q=\#\Phi_q(\F_q)$ and $\Phig_q=\Phi_q(\Fbar_q)$ the geometric component
group; the class of the unit part of a Tate parameter is $c(u)$; the Manin constant is $c_E>0$. Depletion is $f^\Sigma$; the naive twist of a form is
$h\otimes\chi=\sum\chi(n)a_n(h)q^n$, the newform twist $h^\chi$. These are the conventions of \Cref{ch:notation}, and the paper uses the same ones.
\end{convention}

\begin{convention}[Frobenius and $L$-factors]\label{conv:paper:frobenius}
$\Frob_p$ is an arithmetic Frobenius: it induces $x\mapsto x^p$ on the residue field and $\chi_\ell(\Frob_p)=p$. On the tame quotient of inertia,
$\Frob\, t\,\Frob^{-1}=t^p$, equivalently $\Frob^{-1}t\,\Frob=t^{p^{-1}}$. The geometric Frobenius is $\Frob_p^{-1}$ (the letter $\Phi$ is not used for it in this
chapter, since $\Phig_q$ is a component group). For $p\ne\ell$ and a $\Ql$-representation $M$ of $G_{\Q_p}$,
\[
L_p(M,s)=\det\bigl(1-\Frob_p^{-1}p^{-s}\mid M^{I_p}\bigr)^{-1},
\]
the convention of \cite[Sections 1.1.1, 1.7.2]{DFG04}. On $V$ at a good prime the arithmetic Frobenius has the roots $\alpha,\beta$ of $X^2-a_pX+p$; on $\ad V$ both
Frobenii have the eigenvalues $\alpha/\beta,1,\beta/\alpha$; at a multiplicative prime the arithmetic Frobenius acts on $(\ad V)^{I_q}$ by $q$ and the geometric one by $q^{-1}$,
and only the latter gives $L_q(A_f,s)^{-1}=1-q^{-1-s}$. Further, $L_p(\Sym f,s)=L_p(A_f,s-1)$, $L_q(\Sym f,s)=P_q(q^{-s})^{-1}$ defines $P_q\in\Z[X]$, and
$L_p(M(1),s)=L_p(M,s+1)$. These facts are proved in \Cref{ch:local,ch:lfunctions}; the paper states the same convention in Section 1.1 and says that the
matrices in the proof of Lemma lem:tamfactor are those of the arithmetic Frobenius.
\end{convention}

\begin{convention}[duals]\label{conv:paper:duals}
$M^\vee=\Hom(M,\Zl)$ is the linear dual, $X^*=\Hom_{\mathrm{cont}}(X,\Qlmod)$ the Pontryagin dual (the paper's ${}^*$ in Lemma lem:dfg21), $M^D=\Hom(M,\Zl(1))$ the Kummer
dual; a star after the name of a map is a pullback ($\varphi_E^*$). For $A=\ad T$: $W=A\otimes\Qlmod$, $T^D=\Hom(A,\Zl(1))$, $V^D=T^D\otimes\Q_\ell$, $W^D=V^D/T^D$. For $\ell$
odd the trace form is perfect on $\ad T$ (\Cref{ch:reps}), so $T^D\simeq A(1)$ and $W^D\simeq\ad T(1)\otimes\Qlmod$. The adjoint does not see duals or twists:
$\ad(T^\vee)=\ad T=\ad(T(n))=\ad(T\otimes\chi)$.
\end{convention}

\paragraph{How the tour is marked.} Every result of the paper is restated in a numbered environment whose title carries the paper's label (the
paper's labels are quoted as text, for instance Lemma lem:tamfactor). After the statement comes a line saying where the book proves each
ingredient, and then the expansion of the paper's argument, with guidance where a step is hard. Where the book proves more than the paper needs, a one-line
remark says so. Sentences of the paper are quoted verbatim, between quotation marks, with the paper's own symbols; everything else is paraphrase.

\paragraph{What the tour does not claim.} The paper's large-range counts (conductor up to $500000$ and the curves outside the tables) are not
computed by the book's own scripts; the paper's pipeline was rerun in full in the book's repository and reproduced them (\Cref{ch:computation}, in its
section on reproducibility), and the tour reads the paper's tables from their printed rows and quotes their entries as printed. The book's own computations cover conductor below $10{,}000$ ($38{,}042$ optimal curves,
Cremona's files \cite{ecdata}, PARI/GP \cite{PARI} version 2.17.2 through the Python interface \texttt{cypari2}); they are always labelled as the book's, with their
range, and they were made with the paper in hand, so they are checks, not blind tests. A theorem number of another author's paper is said to be \emph{verified}
only when it was compared with the published text; otherwise it is given \emph{as cited} by the paper.

\paragraph{How the paper was produced.} The paper ends with a statement on the use of an AI system, which says that the work was carried out
with an AI system at every stage, under the author's direction, that several of its results and proofs were worked out with the same system while this book
was being written, that the note has not yet been refereed and ``the author has to review it before it is circulated'', and that the author takes full
responsibility for all statements. This book was itself drafted with the assistance of an AI system, under human direction (see the preface). Neither fact
changes what is proved: every argument below is meant to be checked line by line, and every count can be rerun.

\section{The abstract and the introduction}\label{sec:paper:intro}

This section reads the abstract and the first two paragraphs of the paper's Section 1. Their job is to announce the identity (eq:identity), to say
where the modular degree meets the adjoint $L$-function, and to name the inputs. Nothing is proved there; the tour says where each claim is proved.

\paragraph{The abstract, claim by claim.} The abstract makes nine claims. (1) For $E$ optimal and $\ell\nmid N$ odd with $\bar\rho$ absolutely irreducible
on $\Gal(\Qbar/\Q(\sqrt{\ell^*}))$, the identity
\[
v_\ell(\degphi)=\len_{\Zl}H^1_f(\Q,\ad T\otimes\Qlmod)+\sum_{q\mid N}v_\ell\bigl(\Tam_q(\ad T)\bigr)+\sum_{q\mid N}e_\ell(E,q)
\]
holds: this is Theorem thm:main, proved in Section~\ref{sec:paper:selmer} below and in \Cref{ch:congruence}. The parenthesis ``(for $\ell\ge5$ this only asks that
$E[\ell]$ be irreducible)'' is Proposition prop:Sf (\Cref{prop:paper:Sf}). (2) The values of $e_\ell(E,q)$ at a tame $q\ge5$ are Lemma lem:euler (\Cref{lem:paper:euler}). (3) The
Tamagawa factors are trivial at potentially good primes (Lemma lem:tamadd) and at least $t_\ell(E,q)$ at potentially multiplicative ones (Lemma lem:tamfactor). (4) The
proof is an assembly of four inputs from Diamond, Flach and Guo and Fontaine and Perrin-Riou, the local computations and the projection formula; the abstract
adds that ``it also uses an identification of their integral structure with the cohomology of $X_0(N)$, which we prove'': that is Proposition prop:integral
(\Cref{prop:paper:integral}), which holds integrally, for every prime $\ell$, with no localisation. (5) ``The Tamagawa factor raises the bound by at least one whenever a Steinberg prime
$q\equiv1\pmod\ell$ with $\ell\mid n_q$ has a Tate parameter whose unit part is an $\ell$-th power'': a Steinberg prime is a prime of multiplicative reduction (the
local representation is an unramified twist of the Steinberg representation), and the claim is the case $\min(v_\ell(n),v_\ell(q-1),v_\ell(c(u)))\ge1$ of Lemma
lem:tamfactor. (6) The $\kod{I}_0^*$ and $\kod{I}_n^*$ terms hold without the hypothesis over $\Q(\sqrt{\ell^*})$: Theorem thm:quant, which still needs $E[\ell]$ irreducible and
$\ell\nmid2N$. (7) The failure of the inequality at $\ell=3$ for image $N_s(\F_3)$, and (8) its validity with the correction term ``for every curve of conductor at most
$500000$ that the tables determine to be optimal, with $E[\ell]$ irreducible and without complex multiplication'', are data statements (Section~\ref{sec:paper:data});
the two hypotheses matter (\curve{11a1} at the Eisenstein prime $\ell=5$ violates the bound, Section~\ref{sec:paper:appendix}). (9) In the Eisenstein case the
abstract claims only ``that the bound $2$ for the deficit at $\ell=3$ follows from a conjectured bound by the rational torsion'': that is Proposition prop:no9
(\Cref{prop:paper:no9}); the torsion bound itself is Conjecture conj:eis.

\paragraph{The first paragraph of the introduction.} ``When $E[\ell]$ is irreducible, every curve of the class is linked to $E$ by a cyclic isogeny of
degree prime to $\ell$, so $v_\ell$ of the degree of its minimal parametrisation is $v_\ell(\degphi)$''. This is proved in \Cref{ch:ellcurves} (an isogeny of
degree divisible by $\ell$ between curves of the class would give a rational subgroup of order $\ell$, that is, $E[\ell]$ reducible somewhere in the class,
hence everywhere) and \Cref{ch:moddeg} (the degree of $\varphi_{E'}=\psi\circ\varphi_E$ is $\deg\psi\cdot\degphi$). The equivalence of the bound $\degphi\ll N^{2+\varepsilon}$
with a form of the $abc$ conjecture is quoted from \cite[Section 2]{ARS12} and is not used.

\paragraph{The second paragraph: naive and true $L$-functions.} Three facts are used. First, the formula of Zagier and Watkins and the special value
$\langle f,f\rangle$; the exact form is Remark rem:degree, read in Section~\ref{sec:paper:selmer} (\Cref{prop:paper:degree}). Second, Agashe, Ribet and Stein:
$\degphi\mid r_E$ with equal $p$-parts when $p^2\nmid N$ (\Cref{def:paper:congruence-number}). Third, the naive symmetric square
$L^{\nv}(\Sym f,s)=\zeta^{(N)}(2s-2)\sum_{n\ge1}a_{n^2}(f)n^{-s}$, Shimura's series \cite{Shimura75}, whose local factors are the true ones at good and multiplicative
primes and $1$ at the additive primes. A reader should not confuse it with the Rankin series $\sum_na_n(f)^2n^{-s}$, which differs: at a good prime $p$, with $X=p^{-s}$,
\[
\sum_{k\ge0}a_{p^k}^2X^k=\frac{1+pX}{(1-\alpha^2X)(1-pX)(1-\beta^2X)},\qquad\sum_{k\ge0}a_{p^{2k}}X^k=\frac{1+pX}{(1-\alpha^2X)(1-\beta^2X)},
\]
so that $\sum_na_n^2n^{-s}=\zeta^{(N)}(s-1)L^{\nv}(\Sym f,s)/\zeta^{(N)}(2s-2)$ exactly; at a multiplicative prime both local series are $1/(1-X)$, at an additive
prime both are $1$ (\Cref{ch:lfunctions}, \Cref{ex:paper:rankin}). The paragraph ends with the paper's thesis: the primes where the naive and the true adjoint
$L$-functions differ are exactly the primes where the modular degree acquires extra local divisibility, and the amount is the valuation of the omitted factor.
Everything after this is the proof of that sentence and the description of its boundary.

\section[Paper 1.1: the local terms]{Section 1.1 of the paper: the local terms}\label{sec:paper:local-terms}

This section reads the definition of the two local terms $e_\ell(E,q)$ and $t_\ell(E,q)$ and the table of omitted Euler factors. They are the
right-hand side of every inequality in the paper. The table is proved by Lemma lem:euler (Section~\ref{sec:paper:twists-euler}); here we fix the
definitions and their elementary properties.

\begin{definition}[the paper's Definition def:terms]\label{def:paper:terms}
For a bad prime $q$ put $t_\ell(E,q)=v_\ell(-v_q(j_E))$ if $v_q(j_E)<0$ and $t_\ell(E,q)=0$ otherwise; and, if $q\ne\ell$,
\[
e_\ell(E,q)=v_\ell\Bigl(\frac{L_q(A_f,1)^{-1}}{L^{\nv}_q(A_f,1)^{-1}}\Bigr)=v_\ell\Bigl(\frac{P_q(q^{-2})}{1-a_q(E)^2q^{-2}}\Bigr),
\]
where $L^{\nv}_q(A_f,s)=L_q(A_f,s)$ if $\delta_q>0$ and $L^{\nv}_q(A_f,s)=1$ if $\delta_q=0$.
\end{definition}

\noindent\emph{Where proved.} $L_q$ and $P_q$: \Cref{ch:lfunctions}; $v_q(j)$ and Kodaira types: \Cref{ch:reduction}; the Tate parameter: \Cref{ch:reduction}.

\emph{Why these two terms.} The paper compares two objects. The modular degree is computed by the Petersson norm of $f$, and the Petersson norm sees the
\emph{naive} adjoint $L$-value, in which the Euler factor at an additive prime is replaced by $1$ (Remark rem:degree, \Cref{prop:paper:degree}). The Selmer side of
the theorem of Diamond, Flach and Guo sees the \emph{true} $L$-value. The term $e_\ell(E,q)$ is the $\ell$-part, at one prime, of the ratio between the two: it records what
is lost at $q$ when the naive value replaces the true one. The term $t_\ell(E,q)$ has a different origin: when $v_q(j)<0$ the curve is, up to a twist, a Tate curve with
parameter of valuation $n=-v_q(j)$, and $\ell^{t_\ell(E,q)}$ is the $\ell$-part of $n$, the $\ell$-part of the order of the component group of that Tate curve; it enters through
the Tamagawa factor (Lemma lem:tamfactor below). A small example, with the values of Section~\ref{sec:paper:verifications}: \curve{726a1} has type $\kod{I}_6$ at $2$, $\kod{I}_3$
at $3$ and $\kod{III}$ at $11$, with $P_{11}(X)=1+11X$. At $\ell=3$: $e_3(E,2)=0$ (multiplicative), $t_3(E,2)=v_3(6)=1$; $t_3(E,3)=v_3(3)=1$, while $e_3(E,3)$ is not defined
($q=\ell$); and $e_3(E,11)=v_3(1+11\cdot11^{-2})=v_3(12/11)=1$, $t_3(E,11)=0$ ($v_{11}(j)\ge0$).

\emph{Expansion.} (i) The second expression equals the first: $L^{\nv}_q(A_f,1)^{-1}$ is $1-a_q^2q^{-2}$ in every case, since $a_q=\pm1$ at a multiplicative prime, where
the true factor is $1-q^{-2}$, and $a_q=0$ at an additive prime, where the naive factor is $1$; and $L_q(A_f,1)^{-1}=P_q(q^{-2})$ because $L_q(A_f,s)=L_q(\Sym f,s+1)$.
(ii) $e_\ell(E,q)=0$ at a multiplicative prime (naive equals true), and $e_\ell(E,q)=v_\ell(P_q(q^{-2}))$ at an additive prime. (iii) $e_\ell(E,q)$ is a finite
nonnegative integer: $P_q$ has integer coefficients and constant term $1$, $q^{-2}$ is an $\ell$-adic unit, so $P_q(q^{-2})\in\Z[1/q]$ has nonnegative valuation; it is
nonzero because each eigenvalue $\lambda$ of the geometric Frobenius on $(\ad V)^{I_q}$ has complex absolute value $1$ or $q^{-1}$ (\Cref{ch:lfunctions}), so $1-\lambda q^{-1}\ne0$. (iv)
$t_\ell(E,q)=v_\ell(n)$ for type $\kod{I}_n$ at any $q$ and for $\kod{I}_n^*$ at odd $q$ (where $v_q(j)=-n$); at $q=2$ a symbol $\kod{I}_m^*$ can have $v_2(j)\ge0$ (potentially good
reduction), and then $t_\ell(E,2)=0$, or $-v_2(j)\ne m$ (\Cref{ch:reduction}). (v) Twists: $j$ and $\ad V$ are invariant under every quadratic twist, so $t_\ell(E,q)$
and the true factor are. The naive factor at $q$ is invariant under twists unramified at $q$ (the type at $q$ does not change), and so then is $e_\ell(E,q)$. A twist
ramified at $q$ that moves the reduction between multiplicative and $\kod{I}_n^*$, or between good and $\kod{I}_0^*$, changes the naive factor and with it $e_\ell$:
\curve{11a1} has type $\kod{I}_5$ at $11$ and $e_3=0$, its twist \curve{121d2} by $\chi_{-11}$ has type $\kod{I}_5^*$ and $e_3=v_3(120)=1$; the twist \curve{539d2} of \curve{11a1} by
$\chi_{-7}$ lies in the class \curve{539d}, of type $\kod{I}_0^*$ at $7$, with $e_3(\curve{539d1},7)=v_3(6)+v_3(60)=2$, while \curve{11a1} has no term at $7$.

\emph{Reading guidance.} The paper's sentence ``The term $t_\ell$ is invariant under quadratic twist, and $e_\ell(E,q)$ under quadratic twists unramified at $q$'' is
(v); the asymmetry is the whole point of the paper, since the true factor is twist invariant and the naive one is not. Note also that $t_\ell$ is defined at every bad
$q$, $q=\ell$ included, as (eq:rule) and the appendix's function need, while $e_\ell(E,q)$ is defined only for $q\ne\ell$.

\paragraph{The table.} For a bad prime $q\ge5$, $q\ne\ell$, with $E'$ the twist of $E$ by $\Q(\sqrt{q^*})$ and $\alpha,\beta$ the roots of $X^2-a_q(E')X+q$:
\begin{center}\small
\begin{tabular}{lll}
\toprule
type at $q$ & $L_q(A_f,s)^{-1}$ & $e_\ell(E,q)$\\
\midrule
$\kod{I}_n$ & $1-q^{-1-s}$ & $0$\\
$\kod{I}_n^*$ & $1-q^{-1-s}$ & $v_\ell(q^2-1)$\\
$\kod{II},\kod{II}^*,\kod{IV},\kod{IV}^*$ & $1-\chi_{-3}(q)q^{-s}$ & $v_\ell(q-\chi_{-3}(q))$\\
$\kod{III},\kod{III}^*$ & $1-\chi_{-4}(q)q^{-s}$ & $v_\ell(q-\chi_{-4}(q))$\\
$\kod{I}_0^*$ & $(1-\frac\alpha\beta q^{-s})(1-q^{-s})(1-\frac\beta\alpha q^{-s})$ & $v_\ell(q-1)+v_\ell((q+1)^2-a_q(E')^2)$\\
\bottomrule
\end{tabular}
\end{center}
The table is the content of Lemma lem:euler (\Cref{lem:paper:euler}). The sentence after it reads: ``The row $\mathrm{I}_n$ holds at every $q\ne\ell$, $q=2$ included,
the rows $\mathrm{I}_n^*$ and $\mathrm{I}_0^*$ at every odd $q\ne\ell$, and the row $\mathrm{IV}$, $\mathrm{IV}^*$ at every $q\ne3,\ell$''. The rows $\kod{II}$, $\kod{III}$ need $q\ge5$, where
inertia acts tamely; at $q=2,3$ the factor is given by Proposition prop:shapes and computed by PARI.

The last column follows from the middle one and \Cref{def:paper:terms} by evaluating at $s=1$ and comparing with the naive factor. For $\kod{I}_n$ the naive factor is the
true one, so $e_\ell=0$. In the other rows the reduction is additive, the naive factor is $1$, and $e_\ell$ is the valuation of the true factor at $s=1$: $v_\ell(1-q^{-2})=v_\ell(q^2-1)$
for $\kod{I}_n^*$ ($q$ is an $\ell$-adic unit), $v_\ell(1-\chi q^{-1})=v_\ell(q-\chi)$ with $\chi=\chi_{-3}(q)$ or $\chi_{-4}(q)$ for the next two rows, and the $\kod{I}_0^*$ value is computed in the proof of
Lemma lem:euler. One consequence is used repeatedly: $\chi_{-3}(q)$ is the Legendre symbol $(\frac q3)$, so $\chi_{-3}(q)\equiv q\pmod3$ and $3\mid q-\chi_{-3}(q)$; every prime $q\ge5$ of type $\kod{II}$,
$\kod{II}^*$, $\kod{IV}$ or $\kod{IV}^*$ contributes at least $1$ at $\ell=3$. Likewise $\chi_{-4}(q)\equiv q\pmod4$, so $q-\chi_{-4}(q)$ is divisible by $4$, which matters only at $\ell=2$.

\section[Paper 1.2, 1.3: theorem and conjecture]{Sections 1.2 and 1.3 of the paper: the theorem and the conjecture}\label{sec:paper:main-statements}

This section reads the statement of the main theorem, the definition of the Tamagawa factor that enters it, and the conjecture that extends its
inequality beyond the hypothesis. The proofs are read in Section~\ref{sec:paper:selmer}.

\begin{definition}[Tamagawa factor and Selmer group; Section 1.2 of the paper]\label{def:paper:tamagawa}
For a prime $p\ne\ell$ and a $G_{\Q_p}$-stable $\Zl$-lattice $M$, $\Tam_p(M)=\Fitt_{\Zl}\bigl((H^1(I_p,M)_{\mathrm{tor}})^{G_{\Q_p}}\bigr)$, where $G_{\Q_p}/I_p$ acts on $H^1(I_p,M)$ by
conjugation; it is an ideal $\ell^m\Zl$ and $v_\ell(\Tam_p(M))=m$. $H^1_f(\Q,W)$, $W=\ad T\otimes\Qlmod$, is the Bloch-Kato Selmer group \cite[Section 2.1]{DFG04}.
\end{definition}

\noindent\emph{Where proved.} The definition and the identification with $c_p(E)$ for $M=T$: \Cref{ch:blochkato}; the conjugation action: \Cref{ch:cohomology}; the
Bloch-Kato conditions: \Cref{ch:blochkato,ch:galcoh}. Read the parentheses with care: the invariants are taken of the torsion submodule, and the Fitting ideal of
the resulting finite module, as in \cite[I.4.2.2]{FPR94}.

\emph{Why this definition.} At a prime $p\ne\ell$ the Bloch-Kato local condition on $M\otimes\Ql$ is the unramified one. For the lattice $M$ there are two candidate conditions:
classes that are unramified integrally, and classes whose image in $H^1(\Q_p,M\otimes\Ql)$ is unramified. Their difference is measured by the classes that become trivial on $I_p$
after inverting $\ell$, that is, by the torsion of $H^1(I_p,M)$; and the restriction map $H^1(\Q_p,M)\to H^1(I_p,M)^{G_{\Q_p}/I_p}$ is onto (inflation-restriction, $G_{\Q_p}/I_p\simeq\widehat\Z$
having cohomological dimension one), so the part of the torsion fixed by Frobenius is exactly what classes over $\Q_p$ can produce. The ideal $\Tam_p(M)$ is the size
of that discrepancy. Two cases fix the intuition. If $M$ is unramified at $p$, then $H^1(I_p,M)=\Hom(I_p,M)$ is torsion-free and $\Tam_p(M)=\Zl$. If $M=T$ and $E$ has split
multiplicative reduction of type $\kod{I}_n$ at $p$, then $\Tam_p(T)=\ell^{v_\ell(n)}\Zl$, the $\ell$-part of the Tamagawa number $c_p(E)=n$ (step (7) of the proof of
Lemma lem:tamfactor below); this is why the ideal carries the name.

\begin{theorem}[the paper's Theorem thm:main]\label{thm:paper:main}
Let $E$ be $X_0(N)$-optimal and $\ell$ an odd prime, $\ell\nmid2N$, such that $\bar\rho_{E,\ell}$ is absolutely irreducible on $\Gal(\Qbar/\Q(\sqrt{\ell^*}))$. Then
\begin{equation}\label{eq:paper:identity}
v_\ell(\degphi)=\len_{\Zl}H^1_f(\Q,\ad T\otimes\Qlmod)+\sum_{q\mid N}v_\ell\bigl(\Tam_q(\ad T)\bigr)+\sum_{q\mid N}e_\ell(E,q),
\end{equation}
and in particular
\begin{equation}\label{eq:paper:thm}
v_\ell(\degphi)\ge\sum_{q\mid N}e_\ell(E,q)+\sum_{q\mid N}t_\ell(E,q).
\end{equation}
\end{theorem}

\noindent\emph{Where proved.} The assembly is \Cref{ch:congruence}; the proof is read in Section~\ref{sec:paper:selmer}. Its inputs are the theorems of Diamond,
Flach and Guo (stated with references), \cite[Theorem 2.2]{ARS12}, and the identification of the integral structure of \cite{DFG04} with $H^1(X_0(N),\Zl)$
(Proposition prop:integral, \Cref{prop:paper:integral}). The introduction says of this identification that it is not a published statement in the form
needed, and that the paper proves it: ``which we prove; it holds integrally, with no localisation''. The book proves it as well (\Cref{ch:congruence}). The equality in (eq:identity), rather than an inequality, comes
from Proposition prop:pullback (\Cref{prop:paper:pullback}): no freeness or multiplicity-one hypothesis enters.

\emph{Reading the identity.} All three terms on the right of (eq:identity) are nonnegative integers: a length, the valuation of an ideal of $\Zl$, and the integers of
\Cref{def:paper:terms}. The global term $\len H^1_f$ is not computable from local data; the other two are, by the lemmas of Section 2 of the paper. So (eq:thm) is
obtained by dropping the Selmer length and bounding each Tamagawa term from below: $v_\ell(\Tam_q(\ad T))\ge t_\ell(E,q)$ at a potentially multiplicative $q$ (Lemma lem:tamfactor),
and $t_\ell(E,q)=0$ at the other primes. An example where every term is visible: \curve{791a1} has type $\kod{I}_3$ at $7$ and $\kod{I}_2$ at $113$, $\degphi=144$, and surjective mod $3$
image (Cremona's \texttt{galrep} lists only $2\mathrm B$), so the theorem applies at $\ell=3$. Both primes are multiplicative, so $e_3=0$ at both; at $113$, $v_3(2)=0$ and the Tamagawa
term is $0$; at $7$, $v_3(3)=1$, $v_3(7-1)=1$ and the unit part of the Tate parameter reduces to $\bar u=6=(-1)^3$, a cube in $\F_7^\times$, so Lemma lem:tamfactor gives
$v_3(\Tam_7(\ad T))=1+\min(1,1,1)=2$. Then (eq:identity) reads $v_3(144)=2=\len H^1_f+2+0$, so the Selmer group vanishes, while (eq:thm) gives only $2\ge t_3(E,7)=1$.

The paragraph after the theorem says that the hypothesis is $\ell\notin S_f$ in the notation of \cite[(31)]{DFG04} (the set $S_f$ is defined in
\Cref{def:paper:selmer-setup} below), the one under which they prove the
Tamagawa number conjecture for the adjoint motive \cite[Theorem 0.1]{DFG04} (both verified against the published text, \Cref{ch:congruence}), and that for
$\ell\nmid2N$ it means: $E[\ell]$ irreducible and, if $\ell=3$, mod $3$ image not $N_s(\F_3)$ (Proposition prop:Sf, \Cref{prop:paper:Sf}). It then lists what the
paper adds to \cite{DFG04}: the identification of the congruence ideal with the modular degree, the local terms, the twisting theorems,
the boundary of the hypothesis and the computations, each read below.

\begin{conjecture}[the paper's Conjecture conj:euler]\label{conj:paper:euler}
Let $E$ be $X_0(N)$-optimal without complex multiplication, $\ell$ an odd prime with $E[\ell]$ irreducible, and $h_\ell(E)=\len H^0(\Q,\ad T(1)\otimes\Qlmod)$. Then
\begin{equation}\label{eq:paper:rule}
v_\ell(\degphi)\ge\sum_{q\mid N,\,q\ne\ell}e_\ell(E,q)+\sum_{q\mid N}t_\ell(E,q)+[\ell=3]\,w_3(E)-h_\ell(E),
\end{equation}
where $w_3(E)=1$ for the types $\kod{II},\kod{IV},\kod{III}^*$ at $3$, $2$ for $\kod{II}^*,\kod{IV}^*$, and $0$ for $\kod{III},\kod{I}_0^*$; and $w_3(E)=0$ when $3\nmid N$ or the reduction at $3$ is potentially
multiplicative.
\end{conjecture}

\noindent\emph{What is proved about it.} (i) $h_\ell(E)=0$ for every $\ell\ge5$, and $h_3(E)\le1$ with equality exactly for image $N_s(\F_3)$ (Lemma lem:h0 and Proposition
prop:ns, \Cref{lem:paper:h0,prop:paper:ns}); this holds with or without complex multiplication. (ii) Under the hypothesis of Theorem thm:main, (eq:rule) is
(eq:thm), so it is a theorem there. (iii) For image $N_s(\F_3)$ the correction is what the sequence of Lemma lem:dfg21 produces \emph{if}
$H^1_f(\Q,V^D)=0$, which is not known for these curves (\Cref{prop:paper:ns-conditional}); the paper says so (``which is not known for these curves''). (iv) The weights $w_3$ are empirical; the paper's Remark rem:oq1, read in Section~\ref{sec:paper:data}, relates them to a twist relation at
$q=\ell$ observed on the data. In the paragraph after the conjecture, the condition for $h_\ell(E)>0$ is $\bar\rho\simeq\bar\rho\otimes\bar\chi_\ell$, a twist by the cyclotomic
character (not by the quadratic character $\chi_{\ell^*}$); it can hold only at $\ell=3$, where $\bar\chi_3^{\mathrm{cyc}}=\chi_{-3}$.

\emph{The shape of the right side.} The first two sums are the right side of (eq:thm), with the term $e_\ell(E,\ell)$ omitted because it is not defined at $q=\ell$. The weight
$w_3(E)$ stands in for the missing local term at $q=\ell=3$ (it is read off the Kodaira type at $3$, and is $0$ when $3\nmid N$). The correction $h_\ell(E)$ is the length of
$H^0(\Q,W^D)$, the last term of the four-term sequence of Lemma lem:dfg21: when that term is nonzero, the Selmer group $H^1_\Sigma$ is shorter by $h_\ell(E)$ than the sum of its
local pieces (\Cref{prop:paper:ns-conditional}), and the conjecture subtracts exactly that. A worked instance at a curve outside the theorem: \curve{726a1}, with the local
terms computed after \Cref{def:paper:terms}, has image $\Ns$ at $\ell=3$, so $h_3(E)=1$; its reduction at $3$ is multiplicative, so $w_3(E)=0$. The right side of (eq:rule) is
$e_3(E,11)+t_3(E,2)+t_3(E,3)+0-1=1+1+1-1=2$, and $v_3(\degphi)=v_3(144)=2$: equality. Without the correction the bound would be $3$, and the inequality would fail.

\section[Paper 1.4, 1.5: the other statements]{Sections 1.4 and 1.5 of the paper: component groups, history, the other statements}\label{sec:paper:forms-history}

This section reads the two component-group inequalities, Proposition prop:forms, the historical paragraph, and the statements of Theorems thm:qual
and thm:quant, Proposition prop:ns and Conjecture conj:eis, whose proofs come later in the paper.

\begin{definition}[component groups and the twist-symmetrised weight; Section 1.4]\label{def:paper:weights}
For a bad prime $q$, $\#\Phig_q$ is $n$ for $\kod{I}_n$, $1$ for $\kod{II},\kod{II}^*$, $2$ for $\kod{III},\kod{III}^*$, $3$ for $\kod{IV},\kod{IV}^*$ and $4$ for $\kod{I}_0^*$, $\kod{I}_n^*$. The weight
$w_\ell(E,q)$ is $v_\ell(n)$ if the type is $\kod{I}_n$, or $\kod{I}_n^*$ with $q$ odd; $1$ if $\ell=3$ and the type is $\kod{IV}$ or $\kod{IV}^*$, or $\kod{II}$ or $\kod{II}^*$ with $q$ odd; and $0$
otherwise. The two inequalities are
\[
\text{(eq:main)}\quad v_\ell(\degphi)\ge\sum_{q\mid N}v_\ell(\#\Phig_q),\qquad\text{(eq:twist)}\quad v_\ell(\degphi)\ge\sum_{q\mid N}w_\ell(E,q).
\]
\end{definition}

\noindent\emph{Where proved.} The orders of $\Phig_q$: \Cref{ch:reduction}. The paper's identity $w_\ell(E,q)=\max\bigl(v_\ell(\#\Phig_q(E)),v_\ell(\#\Phig_q(E^{(q)}))\bigr)$ for odd $q$
follows from the type exchanges of Lemma lem:twist(ii) by inspection. Written out, with $\ell$ odd and the orders of the component groups listed above: $\kod{I}_n\leftrightarrow\kod{I}_n^*$ gives
$\max(v_\ell(n),v_\ell(4))=v_\ell(n)$; $\kod{I}_0\leftrightarrow\kod{I}_0^*$ gives $\max(0,v_\ell(4))=0$; $\kod{II}\leftrightarrow\kod{IV}^*$ and $\kod{IV}\leftrightarrow\kod{II}^*$ give $\max(0,v_\ell(3))=v_\ell(3)$, which is $1$ for $\ell=3$ and $0$
otherwise; $\kod{III}\leftrightarrow\kod{III}^*$ gives $\max(v_\ell(2),v_\ell(2))=0$. These are the values of the definition, case by case. The point of $w_\ell$ is that it does not change under the
twist by $\chi_{q^*}$, whereas $v_\ell(\#\Phig_q)$ does (\curve{121a1}, of type $\kod{II}$ at $11$, has $\#\Phig_{11}=1$, and its twist \curve{121c2}, of type $\kod{IV}^*$, has $\#\Phig_{11}=3$). At $q=3$
the exchanges $\kod{II}\leftrightarrow\kod{IV}^*$, $\kod{III}\leftrightarrow\kod{III}^*$, $\kod{IV}\leftrightarrow\kod{II}^*$ rest on the formula of Ogg and Saito (\Cref{lem:paper:twist}, \Cref{ch:reduction}); the paper's
parenthesis ``(all five at every odd $q$; at $q=3$ the three middle exchanges keep the conductor exponent, which can exceed $2$)'' says exactly this. Since
$w_\ell\ge v_\ell(\#\Phig_q)$ termwise, (eq:twist) implies (eq:main).

\begin{proposition}[the paper's Proposition prop:forms]\label{prop:paper:forms}
Under the hypothesis of Theorem thm:main, (eq:main) and (eq:twist) hold; the same is true for every curve satisfying (eq:rule) with $h_\ell(E)=0$.
\end{proposition}

\begin{proof}
It suffices to bound $\sum_qw_\ell(E,q)$ by the right side of (eq:thm), respectively (eq:rule), prime by prime; every term of those right sides is
nonnegative. Type $\kod{I}_n$: $w_\ell=v_\ell(n)=t_\ell$. Type $\kod{I}_n^*$, $q$ odd: $w_\ell=v_\ell(n)=t_\ell$ ($v_q(j)=-n$). Type $\kod{I}_n^*$, $q=2$: $w_\ell=0$. Types $\kod{II},\kod{II}^*,\kod{IV},\kod{IV}^*$ at
$q\ge5$: for $\ell\ge5$, $w_\ell=0$; for $\ell=3$, $e_3(E,q)=v_3(q-\chi_{-3}(q))\ge1=w_3$, because $\chi_{-3}(q)\equiv q\pmod3$ ($\chi_{-3}(q)=(\frac q3)$ is $1$ for $q\equiv1$ and $-1$ for
$q\equiv2\pmod3$; values of $e_3$ from \Cref{lem:paper:euler}). The same types at $q=3$ occur only
under (eq:rule) (under Theorem thm:main $3\nmid N$ when $\ell=3$, and $w_\ell=0$ when $\ell\ge5$), and there $w_3(E)\ge1=w_3(E,3)$. Types $\kod{II},\kod{II}^*$ at $q=2$: $w_\ell=0$. Types
$\kod{IV},\kod{IV}^*$ at $q=2$, $\ell=3$: $w_3(E,2)=1=e_3(E,2)$ by Lemma lem:IV (\Cref{lem:paper:IV}). Types $\kod{III},\kod{III}^*,\kod{I}_0^*$: $w_\ell=0$.
\end{proof}

\noindent The only delicate case is the type $\kod{IV}$ or $\kod{IV}^*$ at $q=2$, where Kraus's theorem does not apply; Lemma lem:IV covers it, and the paper's sentence after the
proof checks it on data (``On the $92{,}812$ working-set curves with type $\mathrm{IV}$ or $\mathrm{IV}^*$ at $2$, PARI's factor at $2$ is $1+2X$ for every one of them'').

\paragraph{Section 1.5: the history.} Three statements of other authors are summarised. (i) Ribet and Takahashi compare, for semistable $E$, $\degphi$ with the
degree of a Shimura-curve parametrisation, exactly when $E$ is alone in its isogeny class and otherwise up to primes where $E[\ell]$ is reducible
(\cite{RibetTakahashi97,Takahashi01}; checked against the sources, \Cref{ch:moddeg}). (ii) Pollack and Weston, under their hypothesis CR (surjective $\bar\rho$,
ramified at every prime of the relevant part of the level that is $\pm1$ modulo $\ell$; \cite[Section 1 and Theorem 6.11]{PollackWeston11}, numbering of the arXiv
version). (iii) Kim and Ota \cite[Theorem 1.1]{KimOta}, as cited, under the hypotheses recalled in Remark rem:ko; the identification of their Tamagawa exponent with $t_\ell(E,q)$ is Lemma lem:tamagawa-exp (\Cref{lem:paper:tamagawa-exp}). The
final sentence, that Theorem thm:main contains $v_\ell(\degphi)\ge\sum_{q\|N}t_\ell(E,q)$ including the primes $q\equiv\pm1\pmod\ell$ and $\ell=3$, is correct under its
hypothesis.

\paragraph{The statements proved later.} Theorem thm:qual is \Cref{thm:paper:qual} (Section~\ref{sec:paper:qual}); Theorem thm:quant is
\Cref{thm:paper:quant} (Section~\ref{sec:paper:twisting}), whose proof rests on Lemma lem:step1; Proposition prop:ns is \Cref{prop:paper:ns}
(Section~\ref{sec:paper:boundary}); Conjecture conj:eis is \Cref{conj:paper:eis} below. The sentence after Theorem thm:qual says that for $\ell\ge5$
the theorem needs only the irreducibility of $\bar\rho$: for $\ell\ge5$ the condition over $\Q(\sqrt{-3})$ is not asked, and the optimal-level
theorem needs only irreducibility.

\begin{conjecture}[the paper's Conjecture conj:eis]\label{conj:paper:eis}
Let $E$ be optimal, without complex multiplication, with a rational $\ell$-isogeny, $\ell$ odd, and $D=\sum_qv_\ell(\#\Phig_q)-v_\ell(\degphi)$. Then
$D\le v_\ell(\#E(\Q)_{\mathrm{tors}})+[E[3]\simeq\Z/3\oplus\mu_3]$, and $D$ is at most the number of curves in the class with a rational point of order $\ell$.
\end{conjecture}

\noindent The consequences $D\le0$ without rational $\ell$-torsion, $D\le1$ for $\ell\ge5$ and $D\le2$ for $\ell=3$ follow from the first bound by Proposition prop:no9
(\Cref{prop:paper:no9}) and Mazur's theorem on torsion; only the first bound, and the second one, are conjectural. The section closes with a plan of the paper:
Section 2 the local facts, Section 3 Theorem thm:qual, Section 4 Theorems thm:main and thm:quant, Section 5 the boundary of the hypothesis, Section 6 the
computations, Section 7 the prime $2$; the tour follows it.

\section[Paper 2: inertia]{Section 2 of the paper, part 1: inertia at a bad prime}\label{sec:paper:inertia}

This section reads the first half of the paper's Section 2 (``Local invariants''): the Tamagawa exponent at a multiplicative prime, the inertia
invariants at an additive prime, Lemma lem:IV and Remark rem:types. They feed Theorem thm:qual (cases (a) to (c)), the table of Euler factors, and
Proposition prop:forms.

\begin{lemma}[the paper's Lemma lem:tamagawa-exp]\label{lem:paper:tamagawa-exp}
Let $q$ be a prime of multiplicative reduction, $n_q=v_q(\Delta_E)$, and $\ell\ne q$ a prime. Then $\rho_{E,\ell}\bmod\ell^t$ is unramified at $q$ if and only if $\ell^t\mid n_q$, for
split and for nonsplit reduction.
\end{lemma}

\noindent\emph{Where proved.} Tate curves and their twists: \Cref{ch:reduction}; Kummer theory over $\Q_q^{\ur}$: \Cref{ch:local}; the lemma: \Cref{ch:galrep}.

\begin{proof}[Expanded proof]
Let $K$ be the unramified quadratic extension of $\Q_q$ ($K=\Q_q$ if the reduction is split). Over $K$, $E$ is isomorphic to the Tate curve $\overline K^\times/q_E^{\Z}$ with
$v_q(q_E)=n_q$, and $I_K=I_q$ since $K/\Q_q$ is unramified. The group $E[\ell^t](\overline K)$ is generated by $\mu_{\ell^t}$ and one $\ell^t$-th root $\theta$ of $q_E$. As $\ell\ne q$,
$\mu_{\ell^t}\subset\Q_q^{\ur}$, so $I_q$ fixes $\mu_{\ell^t}$ and sends $\theta$ to $\zeta(\sigma)\theta$ with $\zeta(\sigma)\in\mu_{\ell^t}$; inertia acts trivially exactly when every $\ell^t$-th root of
$q_E$ lies in $\Q_q^{\ur}$, that is, when $q_E\in(\Q_q^{\ur\times})^{\ell^t}$ (Kummer theory over $\Q_q^{\ur}$, which contains $\mu_{\ell^t}$: one root lies in $\Q_q^{\ur}$ iff all do, since
they differ by elements of $\mu_{\ell^t}$; and the fixed field of $I_q$ in $\overline\Q_q$ is $\Q_q^{\ur}$, so $I_q$ fixes $\theta$ iff $\theta\in\Q_q^{\ur}$). Write $q_E=q^{n_q}u$ with $u$ a unit. Units of the
valuation ring of $\Q_q^{\ur}$ are $\ell^t$-th powers: their reductions are, $\Fbar_q^\times$ being divisible, and Hensel's lemma lifts roots since the derivative $\ell^tX^{\ell^t-1}$ of $X^{\ell^t}-u$ is a
unit at a root when $\ell\ne q$. So $q_E$ is an $\ell^t$-th power iff $q^{n_q}$ is. The valuation of $\Q_q^{\ur}$ takes the value $1$ on $q$ (the extension is unramified), so an $\ell^t$-th power
has valuation divisible by $\ell^t$; hence $q^{n_q}$ is not an $\ell^t$-th power unless $\ell^t\mid n_q$, and if $n_q=\ell^tm$ it is $(q^m)^{\ell^t}$. So $q_E$ is an $\ell^t$-th power iff $\ell^t\mid n_q$.
\end{proof}

\noindent The letters matter: $q_E$ is the Tate parameter and $u=q_Eq^{-n}$ its unit part, the letter Lemma lem:tamfactor uses. The proof shows why the lemma holds for
nonsplit reduction too: the unramified quadratic twist does not change the inertia group.

\begin{lemma}[inertia invariants at an additive prime; the paragraph before Lemma lem:IV]\label{lem:paper:inertia-invariants}
Let $q\ne\ell$ be a prime of additive reduction. Then $E[\ell^k]^{I_q}\simeq\Phig_q[\ell^k]$ for every $k\ge1$.
\end{lemma}

\noindent\emph{Where proved.} The N\'eron model and the filtration $E_1\subset E_0\subset E$: \Cref{ch:reduction}; the argument: \Cref{ch:galrep}.

\begin{proof}[Expanded proof]
$E[\ell^k]^{I_q}=E(\Q_q^{\ur})[\ell^k]$, since $\Q_q^{\ur}$ is the fixed field of $I_q$ in $\overline\Q_q$. Let $\mathcal O$ be the valuation ring of $\Q_q^{\ur}$ (or of its completion; torsion is the same). The N\'eron model gives exact sequences
$0\to E_0(\Q_q^{\ur})\to E(\Q_q^{\ur})\to\Phig_q\to0$ and $0\to E_1(\Q_q^{\ur})\to E_0(\Q_q^{\ur})\to\tilde E_{\mathrm{ns}}(\Fbar_q)\to0$, and for additive reduction
$\tilde E_{\mathrm{ns}}(\Fbar_q)\simeq\Fbar_q^+$. The kernel of reduction $E_1$ is the formal group, a pro-$q$ group; $\Fbar_q^+$ has exponent $q$. Both are uniquely $\ell$-divisible ($\ell\ne q$),
hence so is $E_0(\Q_q^{\ur})$ (apply the snake lemma to $\ell^k$ on the second sequence: kernels and cokernels vanish on both ends, so they vanish in the middle). Apply the snake
lemma to multiplication by $\ell^k$ on the first sequence; it gives the exact sequence
\[
0\to E_0(\Q_q^{\ur})[\ell^k]\to E(\Q_q^{\ur})[\ell^k]\to\Phig_q[\ell^k]\to E_0(\Q_q^{\ur})/\ell^k,
\]
whose two outer terms vanish, so $E(\Q_q^{\ur})[\ell^k]\to\Phig_q[\ell^k]$ is an isomorphism. (The component group over $\Q_q^{\ur}$ is the geometric one, $\Phig_q$, since the residue field
$\Fbar_q$ is algebraically closed.)
\end{proof}

\noindent This is the paragraph of the paper between Lemma lem:tamagawa-exp and Lemma lem:IV; it is what makes Lemma lem:IV work at $q=2$, where Kraus's theorem is
not available.

\begin{lemma}[the paper's Lemma lem:IV]\label{lem:paper:IV}
Let $q\ne3$ be a prime at which $E$ has reduction of type $\kod{IV}$ or $\kod{IV}^*$. Then the reduction is potentially good; $\rho_{E,\ell}(I_q)$ is cyclic of order $3$ for every
$\ell\ne q$; $\bar\rho_{E,3}(I_q)$ is a nontrivial unipotent subgroup, $E[3]^{I_q}$ is a line, $\bar\rho_{E,3}$ is tamely ramified at $q$ with Serre conductor exponent $1$, and
the conductor exponent of $E$ at $q$ is $2$. For $\ell\ne q$, $(\ad V_\ell E)^{I_q}$ is a line on which $\Frob_q$ (either one) acts by $\chi_{-3}(q)$, so $P_q(X)=1-\chi_{-3}(q)qX$ and
$e_\ell(E,q)=v_\ell(q-\chi_{-3}(q))$; at $q=2$, $P_2(X)=1+2X$ and $e_3(E,2)=1$.
\end{lemma}

\noindent\emph{Where proved.} \Cref{ch:galrep} (with the case $q=2$); the commutant of an element of order $3$: \Cref{ch:ant}; the tame relation: \Cref{ch:local}; the
conductor: \Cref{ch:reduction,ch:galrep}.

\begin{proof}[Expanded proof]
\emph{Step 1: a line of invariants.} By \Cref{lem:paper:inertia-invariants} with $\ell=3$, $k=1$: $E[3]^{I_q}\simeq\Phig_q[3]=\Z/3$, since $\#\Phig_q=3$. Inertia acts nontrivially on $E[3]$,
otherwise $\Phig_q[3]$ would have order $9$.

\emph{Step 2: potentially good.} If $v_q(j)<0$, over $\Q_q$ the curve is the twist of a Tate curve by a quadratic character $\chi$, ramified as the reduction is additive
(\Cref{ch:reduction}). On $I_q$, in a suitable basis, $\rho_{E,3}(\sigma)=\chi(\sigma)\sm1{n\theta(\sigma)}01$ with $\theta$ the $3$-adic tame character and $n=-v_q(j)$. The subgroup $H\subseteq I_q$
on which $n\theta\equiv0\pmod3$ has index $1$ or $3$ (it is the kernel of the homomorphism $n\theta\bmod3$ to $\Z/3$), so it is not contained in $\ker\chi$, of index $2$ (if $H\subseteq\ker\chi$,
then $2=[I_q:\ker\chi]$ would divide $[I_q:H]\in\{1,3\}$); an element $\sigma\in H\smallsetminus\ker\chi$ has $\bar\rho_{E,3}(\sigma)=\chi(\sigma)\sm1001=-1$, which fixes no nonzero vector of $E[3]$, contradicting
Step 1. So the reduction is potentially good and $\rho_{E,3}(I_q)$ is finite \cite[Section 2]{SerreTate68}.

\emph{Step 3: the image.} $\det\rho_{E,3}=\chi_3^{\mathrm{cyc}}$ is unramified at $q\ne3$, so $\rho_{E,3}(I_q)\subset\SL_2(\Z_3)$. The kernel of $\GL_2(\Z_3)\to\GL_2(\F_3)$ is $1+3\Mat_2(\Z_3)$,
torsion-free (if $g=1+3^aX$ with $a\ge1$ and $X\not\equiv0\pmod3$, then $g^p\equiv1+p3^aX\pmod{3^{a+1}}$ for a prime $p\ne3$, and $g^3\equiv1+3^{a+1}X\pmod{3^{a+2}}$, so
no power of $g$ of prime order is $1$), so reduction is injective on the finite group $\rho_{E,3}(I_q)$. Its image fixes a nonzero vector of $\F_3^2$ and has determinant $1$; in a basis
$(v,w)$ with $v$ fixed such an element is $\sm1b01$, so the image lies in a group of order $3$ of unipotents, and being nontrivial, equals it. Hence $\rho_{E,3}(I_q)$ is cyclic of
order $3$. By \cite[Section 2]{SerreTate68} the kernel of $\rho_{E,\ell}|_{I_q}$ does not depend on $\ell\ne q$, so $\rho_{E,\ell}(I_q)$ is cyclic of order $3$ for every $\ell\ne q$.

\emph{Step 4: tame, conductors.} $3$ is prime to $q$, so the image is a quotient of tame inertia: the ramification is tame, with no Swan term. The Serre conductor exponent
of $\bar\rho_{E,3}$ is $\dim E[3]-\dim E[3]^{I_q}=1$, and the conductor exponent of $E$ is $\dim V-\dim V^{I_q}=2$ ($V^{I_q}=0$: the generator has eigenvalues
$\zeta_3^{\pm1}\ne1$).

\emph{Step 5: the Euler factor.} Let $\sigma$ generate $\rho(I_q)$ on $V=V_\ell E$; $\sigma$ has order $3$ and determinant $1$, so its eigenvalues are $\zeta_3^{\pm1}$ and $\sigma^2+\sigma+1=0$.
An endomorphism commuting with $\sigma$ is a polynomial in $\sigma$ (\Cref{ch:ant}: $\sigma$ has distinct eigenvalues), $a+b\sigma$, of trace $2a-b$; so the trace-zero ones
form the line spanned by $1+2\sigma=\sigma-\sigma^{-1}$. That line is $(\ad V)^{I_q}$. On the tame quotient the arithmetic Frobenius acts by $x\mapsto x^q$, so
$\rho(\Frob_q)\sigma\rho(\Frob_q)^{-1}=\sigma^q=\sigma^{\chi_{-3}(q)}$ ($q\equiv\chi_{-3}(q)\pmod3$). Conjugation sends $\sigma-\sigma^{-1}$ to $\pm(\sigma-\sigma^{-1})$ with the sign $\chi_{-3}(q)$, and a sign
is its own inverse, so the geometric Frobenius acts by the same sign. Hence $L_q(A_f,s)^{-1}=1-\chi_{-3}(q)q^{-s}$ and $P_q(X)=L_q(A_f,s-1)^{-1}|_{X=q^{-s}}=1-\chi_{-3}(q)qX$. The naive factor
is $1$, so $e_\ell(E,q)=v_\ell(1-\chi_{-3}(q)q^{-1})=v_\ell(q-\chi_{-3}(q))$. At $q=2$, $\chi_{-3}(2)=-1$, $P_2=1+2X$, $e_3(E,2)=v_3(3/2)=1$.
\end{proof}

\paragraph{The paragraph after the lemma.} It records that $T_3E$ is free of rank one over $\Z_3[X]/(X^2+X+1)\simeq\Z_3[\zeta_3]$ at every $q\ne3$ where $\rho_{E,3}(I_q)$ is
cyclic of order $3$, $q=2$ included: $T_3E$ is a torsion-free module of $\Z_3$-rank $2$ over that discrete valuation ring, hence free of rank one (\Cref{ch:ant}). Lemma
lem:tamadd uses this, also at $q=2$. Only the identification of the order $3$ with Kraus's semistability defect needs $q\ge5$. At $q=3$ the types $\kod{IV}$, $\kod{IV}^*$
are wildly ramified (the image of inertia has order divisible by $3=q$), and the lemma does not apply.

\begin{remark}[the paper's Remark rem:types]\label{rem:paper:types}
At $q\ge5$, $\ell=3$, the paper lists the image of inertia mod $3$ for the other additive types: $\{\pm1\}$ ($\kod{I}_0^*$), cyclic of order $4$ ($\kod{III}$, $\kod{III}^*$), minus a
unipotent of order $6$ ($\kod{II}$, $\kod{II}^*$), a ramified quadratic character times a unipotent ($\kod{I}_n^*$). Each has $E[3]^{I_q}=0$: $-1$, an element $\tau$ with $\tau^2=-1$,
and $-u$ with $u$ unipotent all fix no nonzero vector (for $-u$: $-uv=v$ gives $uv=-v$, but $u$ has the single eigenvalue $1$). So the Serre exponent is $2$. The twist by
$\chi_{q^*}$ multiplies the generator by $-1$; it lowers the Serre exponent for $\kod{I}_0^*$ (to $0$), $\kod{I}_n^*$ (to $1$, or $0$ when $3\mid n$), $\kod{II}$, $\kod{II}^*$ (to $1$, the twist
being of type $\kod{IV}^*$, $\kod{IV}$), and the conductor of $E$ only for $\kod{I}_0^*$ and $\kod{I}_n^*$ (for $\kod{II}\leftrightarrow\kod{IV}^*$ the exponent stays $2$). These are the facts the
remark states (\Cref{ch:galrep}). They explain why Theorem thm:qual treats exactly the types $\kod{IV}$, $\kod{IV}^*$ (no twist), $\kod{II}$, $\kod{II}^*$ (after the
twist) and $\kod{I}_n^*$ with $n\ge1$ and $\ell\mid n$: those are the rows where, after a twist or none, the Serre exponent drops below the conductor exponent.
\end{remark}

\section[Paper 2: twists and Euler factors]{Section 2 of the paper, part 2: twists and the omitted Euler factors}\label{sec:paper:twists-euler}

This section reads Lemma lem:twist, Lemma lem:euler (the table), Proposition prop:shapes (the four shapes of $P_q$ at every bad prime) and Remark
rem:wild. These give the values of $e_\ell(E,q)$ everywhere. The method is the one of Step 5 of Lemma lem:IV, applied type by type: find the image of inertia, compute the
inertia invariants of $\ad V$ (the trace-zero endomorphisms commuting with that image), and read off how Frobenius acts on them. The quadratic twist enters because it does
not change $\ad V$, so a type that twists to a good or multiplicative one has the Euler factor of the twisted curve.

\begin{lemma}[the paper's Lemma lem:twist]\label{lem:paper:twist}
Let $q$ be odd, $\chi=\chi_{q^*}$, $E^{(q)}$ the twist of $E$ by $\chi$.
\begin{enumerate}[label=(\roman*)]
\item At every $p\ne q$, $E$ and $E^{(q)}$ have the same Kodaira type, the same $\Phig_p$ and the same conductor exponent.
\item If $E$ has type $\kod{I}_n$ at $q$ ($n\ge0$), $E^{(q)}$ has type $\kod{I}_n^*$, and conversely. If $E$ has type $\kod{II}$, $\kod{III}$, $\kod{IV}$ at $q$, then $E^{(q)}$ has type $\kod{IV}^*$,
$\kod{III}^*$, $\kod{II}^*$, and conversely. For the exchanges
$\kod{I}_n\leftrightarrow\kod{I}_n^*$, $\kod{I}_0\leftrightarrow\kod{I}_0^*$ at every odd $q$ and for all the exchanges at $q\ge5$, the conductor exponent of $E^{(q)}$ at $q$ is $2$ when its reduction is
additive and $1$ when it is of type $\kod{I}_n$, $n\ge1$. At $q=3$ the exchanges $\kod{II}\leftrightarrow\kod{IV}^*$, $\kod{III}\leftrightarrow\kod{III}^*$, $\kod{IV}\leftrightarrow\kod{II}^*$ keep $f_3$
unchanged, and $f_3$ can exceed $2$ (wild ramification).
\item (Atkin and Li) If $g$ is a newform of weight $2$ on $\Gamma_0(M)$ with $q^2\nmid M$, then $g\otimes\chi$ is a newform on $\Gamma_0(\lcm(M,q^2))$.
\end{enumerate}
\end{lemma}

\noindent\emph{Where proved.} (i), (ii) for $q\ge5$ and the $\kod{I}_n\leftrightarrow\kod{I}_n^*$, $\kod{I}_0\leftrightarrow\kod{I}_0^*$ exchanges at every odd $q$: \Cref{ch:reduction}; the exchanges at $q=3$:
\Cref{ch:reduction}, given the formula of Ogg and Saito; (iii) is stated with a reference \cite[Theorem 3.1]{AtkinLi78} and sketched in \Cref{ch:hecke}.

\begin{proof}[Expanded proof of (i) and (ii)]
(i) $\chi$ is unramified at $p$, so $E\simeq E^{(q)}$ over an unramified extension of $\Q_p$; Kodaira type, $\Phig_p$ and conductor exponent are invariant under unramified base
change (\Cref{ch:reduction}). (ii) A minimal model with invariants $c_4,c_6,\Delta$ twists to a model with $(q^*)^2c_4$, $(q^*)^3c_6$, $(q^*)^6\Delta$ and the same $j$. If $v_q(j)<0$, the
type is $\kod{I}_n$ or $\kod{I}_n^*$ with $n=-v_q(j)$ according as the reduction is multiplicative or not, and a ramified quadratic twist exchanges the two (over $\Q_q$ the curve is a
Tate curve or its ramified twist, \Cref{ch:reduction}). For $\kod{I}_0\leftrightarrow\kod{I}_0^*$: a model $y^2=x^3+a_2x^2+a_4x+a_6$ with good reduction twists to
$y^2=x^3+q^*a_2x^2+(q^*)^2a_4x+(q^*)^3a_6$; Tate's algorithm, after dividing out the powers of $q$ (substitute $x=qT$ and divide the cubic by $q^3$, using $q^*=\epsilon q$ and $\epsilon^2=1$), looks at $T^3+\epsilon a_2T^2+a_4T+\epsilon a_6$, $\epsilon=q^*/q=\pm1$, whose roots are $\epsilon$
times those of $x^3+a_2x^2+a_4x+a_6$, distinct modulo $q$; distinct roots is the criterion for $\kod{I}_0^*$; conversely a type $\kod{I}_0^*$ model at odd $q$ has this shape and
twists to a good one. For $q\ge5$ and potentially good reduction the type is read off $v_q(\Delta_{\min})\in\{0,2,3,4,6,8,9,10\}$; twisting adds $6$, and subtracting $12$ when the
result is at least $12$ gives the minimal model: $2\leftrightarrow8$, $3\leftrightarrow9$, $4\leftrightarrow10$. The conductor exponents are the tame ones, $\dim V-\dim V^{I_q}$, since $V_\ell E^{(q)}=V_\ell E\otimes\chi$ with $\chi$
tame and $V_\ell E$ good, multiplicative or (for $q\ge5$) tame. At $q=3$, for the types $\kod{II},\kod{III},\kod{IV}$ and their duals: $\chi_{-3}$ is tame, so the Swan conductor does not change, and neither $V_\ell E$ nor its twist has inertia invariants (the
image of inertia is not contained in $\{\pm1\}$, so it contains an element of order $3$, $4$ or $6$, which has no eigenvalue $\pm1$), so $f_3$ is unchanged; $v_3(\Delta_{\min})$ changes
by $6$ modulo $12$; by Ogg's formula $v_3(\Delta_{\min})=f_3+m-1$ with $m$ the number of components of the minimal regular model, $1\le m\le9$ for the potentially good types
other than $\kod{I}_0^*$, which forces the exchange $m\mapsto m\pm6$, that is the three exchanges listed (\Cref{ch:reduction}).
\end{proof}

\noindent\emph{Reading guidance.} Two points in (ii) are easy to misread. The sentence on conductor exponents applies to the exchanges it names, not to the
exchanges at $q=3$: there $f_3$ is unchanged and can be $3$ or more, by wild ramification; for example \curve{54a1}, of type $\kod{IV}^*$ at $3$, and its twist by
$\chi_{-3}$, \curve{54b1}, of type $\kod{II}$ at $3$, both have $f_3=3$ (recomputed). And the exchanges at $q=3$ are where the proof leaves Tate's algorithm: it combines the
invariance of the Swan conductor under a tame twist with Ogg's formula, and the bound $1\le m\le9$ on the number of components then forces the shift $m\mapsto m\pm6$. The
appendix of the paper reports a numerical check on $5{,}122$ pairs below conductor $10{,}000$, which the book reproduced (\Cref{ch:reduction}).

\begin{lemma}[the paper's Lemma lem:euler]\label{lem:paper:euler}
Let $q\ge5$ be a bad prime, $q\ne\ell$. Then $L_q(A_f,s)^{-1}$ and $e_\ell(E,q)$ are as in the table of Section~\ref{sec:paper:local-terms}. The rows $\kod{I}_n^*$ and $\kod{I}_0^*$ hold at every
odd $q\ne\ell$, and the row $\kod{I}_n$ at every $q\ne\ell$.
\end{lemma}

\noindent\emph{Where proved.} \Cref{ch:lfunctions} (in its own language); the inertia images: \Cref{ch:galrep} (Kraus's theorem \cite{Kraus90} for $q\ge5$); the $\kod{I}_0^*$ identity:
\Cref{ch:rings}.

\begin{proof}[Expanded proof]
In every case $L_q(A_f,s)^{-1}=\prod(1-\lambda q^{-s})$ over the eigenvalues $\lambda$ of $\Frob_q^{-1}$ on $(\ad V)^{I_q}$ (\Cref{conv:paper:frobenius}).

\emph{Type $\kod{I}_n$, any $q\ne\ell$.} In a basis $e_1,e_2$ of $T$, $\rho|_{G_{\Q_q}}=\sm{\chi_\ell\eps}{\eps\kappa}0{\eps}$ (\Cref{lem:paper:tamfactor}), and a generator $t$ of tame $\ell$-inertia acts by
$te_1=e_1$, $te_2=e_2+n'e_1$, $n'\ne0$. In the basis $e=e_1\otimes e_2^\vee$, $h=e_1\otimes e_1^\vee-e_2\otimes e_2^\vee$, $f=e_2\otimes e_1^\vee$ of $\ad V$ (as matrices, $e=\sm0100$, $h=\sm100{-1}$,
$f=\sm0010$), $t$ acts by conjugation $X\mapsto tXt^{-1}$ with $t=\sm1{n'}01$, $t^{-1}=\sm1{-n'}01$. Multiplying out, $tet^{-1}=e$, $tht^{-1}=\sm1{-2n'}0{-1}=h-2n'e$ and
$tft^{-1}=\sm{n'}{-n'^2}1{-n'}=f+n'h-n'^2e$; so $(t-1)e=0$, $(t-1)h=-2n'e$, $(t-1)f=n'h-n'^2e$. A vector $xe+yh+zf$ is fixed iff $-2n'y-n'^2z=0$ and $n'z=0$, that is $y=z=0$:
$(\ad V)^{I_q}=\Ql e$. The arithmetic Frobenius sends $e$ to $(\eps q)\eps^{-1}e=qe$, so the geometric one acts by $q^{-1}$ and $L_q(A_f,s)^{-1}=1-q^{-1-s}$.
This is also the naive factor ($\delta_q=1$), so $e_\ell=0$.

\emph{Type $\kod{I}_n^*$, $q$ odd.} By \Cref{lem:paper:twist}(ii), $V=V'\otimes\chi_{q^*}$ with $V'=V_\ell E^{(q)}$ multiplicative; $\ad V=\ad V'$, so the true factor is $1-q^{-1-s}$; the
naive factor is $1$, and $e_\ell=v_\ell(1-q^{-2})=v_\ell(q^2-1)$.

\emph{Types $\kod{II},\kod{II}^*,\kod{IV},\kod{IV}^*$, $q\ge5$.} By Kraus's theorem the image of $I_q$ on $V$ is cyclic of order $3$ or $6$, generated by $\sigma$ or $-\sigma$ with $\sigma$ of order $3$.
On $\ad V$ the sign disappears; the computation is Step 5 of \Cref{lem:paper:IV}: the invariant line is spanned by $2\sigma+1$, of square $-3$, inside $\End_{I_q}(T)=\Zl[\sigma]\simeq\Zl\otimes\Z[\zeta_3]$
(not a field when $\ell\equiv1\pmod3$), and either Frobenius acts on it by $\chi_{-3}(q)$; $e_\ell=v_\ell(q-\chi_{-3}(q))$.

\emph{Types $\kod{III},\kod{III}^*$, $q\ge5$.} The image is cyclic of order $4$, generated by $\tau$ with $\tau^2=-1$ (determinant $1$, eigenvalues $\pm i$). The commutant is $\Ql[\tau]$, and its
trace-zero part is the line $\Ql\tau$ ($\tau$ has trace $0$; it is diagonalisable only over $\Ql(i)$, so no diagonal basis exists over $\Ql$ when $\ell\equiv3\pmod4$). Frobenius sends $\tau$
to $\tau^q=\tau^{\chi_{-4}(q)}$, so it acts on the line by $\chi_{-4}(q)$ (either Frobenius), and $e_\ell=v_\ell(q-\chi_{-4}(q))$.

\emph{Type $\kod{I}_0^*$, $q$ odd.} $V=V'\otimes\chi_{q^*}$ with $V'=V_\ell E'$ unramified, arithmetic Frobenius eigenvalues $\alpha,\beta$; $\ad V=\ad V'$ unramified with eigenvalues
$\alpha/\beta,1,\beta/\alpha$ for either Frobenius. With $\alpha\beta=q$ and $\alpha+\beta=a=a_q(E')$,
\[
\Bigl(1-\frac{\alpha}{\beta q}\Bigr)\Bigl(1-\frac1q\Bigr)\Bigl(1-\frac{\beta}{\alpha q}\Bigr)=\frac{q-1}{q}\Bigl(1+\frac1{q^2}-\frac{\alpha^2+\beta^2}{q^2}\Bigr)=\frac{q-1}{q}\cdot\frac{(q+1)^2-a^2}{q^2},
\]
since $\alpha/\beta+\beta/\alpha=(\alpha^2+\beta^2)/q=(a^2-2q)/q$ and $q^2+1-a^2+2q=(q+1)^2-a^2$. So $e_\ell=v_\ell(q-1)+v_\ell((q+1)^2-a^2)$. Only the exchanges $\kod{I}_n\leftrightarrow\kod{I}_n^*$ and
$\kod{I}_0\leftrightarrow\kod{I}_0^*$ were used in the last two cases, valid at every odd $q$; at $q=3$ the values are $L_3(A_f,1)^{-1}=\frac89$ ($\kod{I}_n^*$) and $\frac23\cdot\frac{16-a_3(E')^2}9$ ($\kod{I}_0^*$).
\end{proof}

\noindent\emph{Reading guidance.} Two details of the paper's proof deserve attention. The commutant of the image of inertia is taken in $\End_{I_q}(T)=\Zl[\sigma]\simeq\Zl\otimes\Z[\zeta_3]$,
which is not a field when $\ell\equiv1\pmod3$; only the trace-zero line inside it matters. And for the types $\kod{III}$, $\kod{III}^*$ the generator $\tau$ is itself the trace-zero
invariant, with no diagonalisation over $\Ql$ needed. The sign by which Frobenius acts is the same for the arithmetic and the geometric Frobenius, as a sign is its
own inverse. The closing sentence of the paper's proof, that the values agree with PARI's \texttt{lfunsympow} on every curve of the tables, was
reproduced by the book on all $67{,}745$ pairs (optimal curve, bad $q\ge5$) below conductor $10{,}000$ (\Cref{ch:computation}).

\begin{proposition}[the paper's Proposition prop:shapes]\label{prop:paper:shapes}
Let $q$ be a bad prime and $\ell\ne q$ odd. Then $P_q(X)$ is (i) $1-X$ if $v_q(j)<0$, and then $e_\ell(E,q)=0$ or $v_\ell(q^2-1)$ according as the reduction is multiplicative or
additive; (ii) $(1-qX)(1-(a^2-2q)X+q^2X^2)$ if the reduction is potentially good and $\rho(I_q)=\{\pm1\}$, with $a=a_q(E')$ for a quadratic twist $E'$ with good reduction at
$q$, and then $e_\ell(E,q)=v_\ell(q-1)+v_\ell((q+1)^2-a^2)$; (iii) $1-cqX$ with $c=\pm1$ if the reduction is potentially good and $\rho(I_q)$ is cyclic of order $3$, $4$ or $6$, and
then $e_\ell(E,q)=v_\ell(q-c)$; (iv) $1$ if the reduction is potentially good and $\rho(I_q)$ is not abelian, which happens only for $q=2,3$, and then $e_\ell(E,q)=0$.
\end{proposition}

\noindent\emph{Where proved.} \Cref{ch:lfunctions}; finite subgroups of $\SL_2$: \Cref{ch:groups,ch:reps}; potentially good reduction and finite inertia: \Cref{ch:galrep}.

\begin{proof}[Expanded proof]
The cases are exhaustive: if $v_q(j)<0$ we are in (i); otherwise $\rho(I_q)$ is a finite subgroup of $\SL_2(\Zl)$ ($\det=\chi_\ell$ is trivial on $I_q$) whose elements have
characteristic polynomials in $\Z[X]$ independent of $\ell$ \cite[Section 2]{SerreTate68}, so their orders are in $\{1,2,3,4,6\}$. (An element $g$ of finite order $m$ and determinant $1$ is diagonalisable over
$\overline\Q_\ell$ with eigenvalues $\zeta,\zeta^{-1}$, and $m$ is the order of the root of unity $\zeta$. The cyclotomic polynomial $\Phi_m$, the minimal polynomial of $\zeta$ over $\Q$, divides the
characteristic polynomial of $g$, of degree $2$ in $\Z[X]$; so $\deg\Phi_m\le2$, which holds exactly for $m\in\{1,2,3,4,6\}$.) A finite abelian subgroup of $\SL_2$ over a field of
characteristic $0$ is cyclic (simultaneously diagonalisable over the algebraic closure, and $\mathrm{diag}(x,x^{-1})\mapsto x$ embeds it in the multiplicative group). Order $1$ is good
reduction; order $2$ is $\{\pm1\}$, case (ii); orders $3,4,6$ are case (iii); otherwise (iv).

(i) Over $\Q_q$, $\rho\simeq\rho'\otimes\chi$ with $\rho'$ a Tate curve's representation up to an unramified twist and $\chi$ quadratic; $\ad V=\ad V'$, and the $\kod{I}_n$ computation of
\Cref{lem:paper:euler} gives $L_q(A_f,s)^{-1}=1-q^{-1-s}$, so $P_q(X)=1-X$. The naive factor is $1-q^{-2}$ or $1$.

(ii) $\psi=\rho|_{I_q}$ is a quadratic character of $I_q$, invariant under conjugation by $G_{\Q_q}$ (its values are $\pm1$, central). Choose a Frobenius lift $\phi$ topologically
generating a complement to $I_q$; $\chi(i\phi^k)=\psi(i)$ defines a quadratic character of $G_{\Q_q}$ extending $\psi$. Every quadratic character of $G_{\Q_q}$ is the local component of a
quadratic Dirichlet character $\chi'$ (local class field theory: $\Q_q^\times/(\Q_q^\times)^2$ has order $4$ for $q$ odd and $8$ for $q=2$, and the local components of $\chi_{q^*}$ and of a character
unramified at $q$, respectively of $\chi_{-4}$, $\chi_8$, $\chi_5$, give that many distinct characters), and $E'=E\otimes\chi'$ has $\rho_{E'}|_{I_q}=\psi^2=1$: good reduction. Then $\ad V=\ad V_\ell E'$ has eigenvalues
$\alpha/\beta,1,\beta/\alpha$ and $P_q(X)=(1-\alpha^2X)(1-qX)(1-\beta^2X)$ (the geometric eigenvalues times $q$, as $L_q(A_f,s-1)$), which is the stated polynomial as
$\alpha^2+\beta^2=a^2-2q$ and $\alpha^2\beta^2=q^2$. Only $a^2$ enters, so the sign ambiguity in the choice of $E'$ does not matter.

(iii) Let $g$ generate $\rho(I_q)$, of order $m\in\{3,4,6\}$, eigenvalues $\zeta^{\pm1}$, $\zeta\ne\zeta^{-1}$. Trace-zero endomorphisms commuting with $g$ form the line spanned by $g-g^{-1}$
(commutant $\Ql[g]$, two-dimensional). Frobenius normalises $\rho(I_q)$, so it conjugates $g$ to a generator $g^k$ with the same eigenvalue set $\{\zeta,\zeta^{-1}\}$, so to $g^{\pm1}$; it acts on
the line by a sign $c$, the same for both Frobenii. Hence $P_q(X)=1-cqX$ and $e_\ell=v_\ell(1-c/q)=v_\ell(q-c)$.

(iv) A faithful two-dimensional representation of a nonabelian finite group in characteristic $0$ is absolutely irreducible (a reducible one would be, after
semisimplification, a sum of two characters, hence abelian image), so its commutant is the scalars and $(\ad V)^{I_q}=0$. For $q\ge5$ the image is a quotient of tame
inertia (its order divides $24$, prime to $q$), hence cyclic; so (iv) needs $q\in\{2,3\}$.
\end{proof}

\noindent\emph{Reading guidance.} The proposition is what lets the paper use PARI's Euler factor at $q=2,3$ without a Kodaira-type table: whatever the wild inertia does, the
factor has one of four shapes, and $e_\ell(E,q)$ is read off it. The only cubic shape is (ii).

\begin{remark}[the paper's Remark rem:wild]\label{rem:paper:wild}
The remark says that at $q=2,3$ the Kodaira type does not determine $P_q$, gives \curve{336f1} (type $\kod{II}$ at $2$, $P_2=1+2X$, $e_3=1$, $\degphi=24$), \curve{1232l1} (type
$\kod{II}^*$ at $2$, the twist of \curve{77a1} by $\chi_{-4}$, $P_2=(1-2X)(1+2X)^2$, $e_3(E,2)=2$, $\degphi=288$: the bound of Theorem thm:main at $\ell=3$ attained by the prime $2$ alone) and
\curve{592c1} (type $\kod{II}^*$ at $2$, the twist of \curve{37a1} by $\chi_{-4}$, $P_2=(1-2X)(1+4X^2)$, $e_5(E,2)=1$, $\degphi=80$). All these values were recomputed for the book
(Section~\ref{sec:paper:verifications}). The two counts of the remark name their populations: the cubic factors at $2$ ($1{,}143$) and at $3$ ($2{,}283$) are counted
over all curves below conductor $10{,}000$, and $534$ and $1{,}308$ over the optimal non-CM curves (\Cref{ch:lfunctions}); since (ii) is the only cubic shape, the
empirical content is that each cubic equals the formula of (ii) for a twist of good reduction found independently. The share of potentially good reduction among
the types $\kod{I}_n^*$ at $2$ is $58\%$ ($2{,}340$ of $4{,}037$) among optimal curves and $61\%$ ($5{,}935$ of $9{,}667$) among all curves (\Cref{ch:reduction}). The book reproduced all
four counts.
\end{remark}

\section[Paper 2: Tamagawa factors]{Section 2 of the paper, part 3: the Tamagawa factors of the adjoint}\label{sec:paper:tamagawa}

This section reads Lemma lem:tamfactor (potentially multiplicative primes) and Lemma lem:tamadd (potentially good additive primes). With them every term
of (eq:identity) other than the Selmer length is explicit (Corollary cor:explicit). Both proofs compute $H^1(I_q,\ad T)$ from the image of inertia found in the previous two
sections: at a potentially multiplicative prime inertia acts through a unipotent, and its first cohomology is the coinvariants, which have torsion; at a potentially good prime it
acts through a finite group of order dividing $24$, the torsion of the cohomology comes from that finite group, so it is killed by the group's order; it vanishes at $\ell\ge5$ for that
reason, and at $\ell=3$ because $\ad T$ is a free module over the group ring of a Sylow $3$-subgroup.

\begin{lemma}[the paper's Lemma lem:tamfactor]\label{lem:paper:tamfactor}
Let $q$ be a prime of multiplicative reduction, $n=n_q$, $\ell$ odd, $\ell\ne q$, $M=\ad T$. Let $q_E$ be the Tate parameter, $u=q_Eq^{-n}\in\Z_q^\times$, and $c(u)$ the class of the
reduction $\bar u$ in $\F_q^\times/(\F_q^\times)^{\ell^v}\simeq\Z/\ell^v$, $v=v_\ell(q-1)$ (so $v_\ell(c(u))\ge k$ iff $\bar u$ is an $\ell^k$-th power, for $k\le v$). Then $H^1(I_q,M)_{\mathrm{tor}}\simeq(\Z/\ell^{v_\ell(n)})^2$ and
\[
v_\ell\bigl(\Tam_q(\ad T)\bigr)=v_\ell(n)+\min\bigl(v_\ell(n),v_\ell(q-1),v_\ell(c(u))\bigr).
\]
In particular $v_\ell(\Tam_q(\ad T))\ge t_\ell(E,q)$, with equality unless $q\equiv1\pmod\ell$, $\ell\mid n$ and $\bar u$ is an $\ell$-th power. The same formula holds at an additive potentially
multiplicative prime $q\ne\ell$, $q=2$ included, with $n=-v_q(j_E)$ and $q_E$ the Tate parameter with $j(q_E)=j_E$.
\end{lemma}

\noindent\emph{Where proved.} \Cref{ch:blochkato} (the lemma), \Cref{ch:cohomology} (the cohomology of $\langle t\rangle$ and the conjugation action), \Cref{ch:galcoh} (the
Kummer class of $q_E$ and the identification of $c'$), \Cref{ch:modules} (Fitting ideals of finite modules).

\begin{proof}[Expanded proof]
(1) \emph{The Tate matrices.} Over the unramified quadratic extension $E$ is a Tate curve; $T$ has a basis $e_1,e_2$ in which $g\in G=G_{\Q_q}$ acts by $\sm{\chi_\ell\eps(g)}{\eps\kappa(g)}0{\eps(g)}$,
$\eps$ unramified of order at most $2$ with $\eps(\Frob)=a_q(E)$, $\kappa$ the Kummer cocycle of $q_E$ in $H^1(G,\Zl(1))=\varprojlim\Q_q^\times/(\Q_q^\times)^{\ell^k}$. The corner is $\eps\kappa$: the
nonsplit curve is the Tate curve twisted by $\eps$, and $\sm{\chi_\ell}{\kappa}01\otimes\eps=\sm{\chi_\ell\eps}{\eps\kappa}0{\eps}$; a corner $c$ must satisfy $c(gh)=\chi_\ell\eps(g)c(h)+c(g)\eps(h)$, which
$\eps\kappa$ does and $\kappa$ does not when $\eps\ne1$.

(2) \emph{Inertia.} $I=I_q$ acts through its maximal pro-$\ell$ quotient, $\Zl(1)$, generated by $t$; $\eps|_I=1$ and $\kappa|_I=n$ times the tame character, so $te_1=e_1$,
$te_2=e_2+n'e_1$ with $v_\ell(n')=v_\ell(n)$. The kernel of $I\to\langle t\rangle$ has no pro-$\ell$ quotient, so inflation-restriction gives $H^1(I,M')=H^1(\langle t\rangle,M')=M'/(t-1)M'$
via $c\mapsto c(t)$, for every finitely generated $\Zl$-module $M'$ with this action (\Cref{ch:cohomology}).

(3) \emph{The torsion.} In the basis $e,h,f$ of $M$: $(t-1)e=0$, $(t-1)h=-2n'e$, $(t-1)f=n'h-n'^2e$. As $2$ is a unit, $(t-1)M=n'(\Zl e\oplus\Zl h)$, so
$H^1(I,M)=(\Zl/n')\bar e\oplus(\Zl/n')\bar h\oplus\Zl\bar f$, with torsion $(\Z/\ell^{v_\ell(n)})^2$.

(4) \emph{Frobenius.} $G/I$ acts by $(\sigma\cdot c)(i)=\sigma c(\sigma^{-1}i\sigma)$. With $\Frob^{-1}t\,\Frob=t^{q^{-1}}$ and $c(t^k)\equiv kc(t)$ modulo $(t-1)M$, the arithmetic Frobenius acts on
$M/(t-1)M$ by $q^{-1}\Ad(\rho(\Frob))$, conjugation by $\rho(\Frob)$ followed by multiplication by $q^{-1}$, as the paper says. The factor $q^{-1}$ comes from the twist
$\Frob^{-1}t\Frob=t^{q^{-1}}$ of the tame generator, and it is what makes the count depend on $q-1$.

(5) \emph{The fixed points.} Write $\rho(\Frob)=\sm{\eps q}{c'}0{\eps}$, $\eps=\eps(\Frob)=\pm1$, $c'=\eps\kappa(\Frob)$. Conjugation sends $e\mapsto qe$ and $h\mapsto h-2c'\eps e$; so on the torsion
$\bar e\mapsto\bar e$ and $\bar h\mapsto q^{-1}(\bar h-2c'\eps\bar e)$. The element $x\bar e+y\bar h$ is fixed iff $(q-1)y\equiv0$ and $c'y\equiv0$ modulo $\ell^{v_\ell(n)}$ (compare the $\bar h$ and $\bar e$
coefficients, using that $q$, $2$, $\eps$ are units: the image of $x\bar e+y\bar h$ is $(x-2c'\eps q^{-1}y)\bar e+q^{-1}y\bar h$). There are $\ell^{v_\ell(n)}$ choices of $x$ and $\ell^{\min(v_\ell(n),v_\ell(q-1),v_\ell(c'))}$ of $y$.
(Count in $\Z/\ell^m$, $m=v_\ell(n)$: for $a=v_\ell(q-1)$ and $b=v_\ell(c')$, the conditions $\ell^ay\equiv0$ and $\ell^by\equiv0$ modulo $\ell^m$ say $v_\ell(y)\ge m-\min(a,b)$, and the number of such $y$ is
$\ell^{\min(m,a,b)}$.) The Fitting ideal of a finite $\Zl$-module of order
$\ell^m$ is $\ell^m\Zl$ (\Cref{ch:modules}), and $\Tam_q(M)$ is that Fitting ideal \cite[I.4.2.2]{FPR94}, recalled in \cite[p.~708]{DFG04}.

(6) \emph{The value of $c'$.} $q$ being a uniformiser, $H^1(\Q_q,\Zl(1))=\Zl\kappa_q\oplus H^1_{\ur}$, and $\kappa=n\kappa_q+\kappa_u$ with $\kappa_u$ inflated from the class of $u$ in
$H^1(G/I,\Zl(1))=\Zl(1)/(q-1)\simeq\F_q^\times\otimes\Zl$. For the Frobenius lift fixing a compatible system of $\ell$-power roots of $q$ (which exists, $\Q_q(q^{1/\ell^\infty})$ being totally
ramified), $\kappa_q(\Frob)=0$ and $c'=\eps\kappa_u(\Frob)$, the class of $\bar u$ up to the sign $\eps$. Another lift $\Frob t^k$ changes $c'$ by $\eps qkn'$; another choice of roots changes $\kappa$ by
a coboundary, hence $c'$ by a multiple of $q-1$. So $c'$ is well defined modulo $\ell^{\min(v_\ell(n),v_\ell(q-1))}$, which does not change the minimum, and $v_\ell(c')$ may be replaced by
$v_\ell(c(u))$ (\Cref{ch:galcoh}).

(7) \emph{Check for $M=T$.} The same computation gives $H^1(I,T)_{\mathrm{tor}}=(\Z/\ell^{v_\ell(n)})\bar e_1$ with $\Frob$ acting by $q^{-1}\cdot\eps q=\eps$: all invariant if split, none if nonsplit
($\ell$ odd), which is the $\ell$-part of $c_q(E)$, as it should be.

(8) \emph{The ``in particular''.} The extra term is positive iff $v_\ell(n)\ge1$, $v_\ell(q-1)\ge1$ and $v_\ell(c(u))\ge1$. At $q=2$, $v_\ell(q-1)=0$, so the extra term vanishes.

(9) \emph{Additive potentially multiplicative $q$.} Over $\Q_q$, $E$ is the twist of the Tate curve $E_{q_E}$ by a ramified quadratic character; $\ad T$ is unchanged as a $G_{\Q_q}$-module, and the
computation used only that module, with $n=v_q(q_E)=-v_q(j_E)$ (not the index of the Kodaira symbol at $q=2$).
\end{proof}

\noindent\emph{Reading guidance.} Three points are easy to get wrong: the corner entry is $\eps\kappa$, not $\kappa$ (step (1) and \Cref{ex:paper:corner}); the value $c'$ is
determined only modulo $\ell^{\min(v_\ell(n),v_\ell(q-1))}$, which is harmless (step (6)); and at an additive potentially multiplicative prime the argument is local, so it
covers $q=2$ (step (9)). Worked values, recomputed for the book with PARI's $q$-adic
Tate parameter: \curve{791a1} at $7$ ($\kod{I}_3$, $\bar u=6$, a cube): $v_3(\Tam_7(\ad T))=2$, $t_3=1$; \curve{539c1} at $7$ ($\kod{I}_3^*$, the same $q_E$ as \curve{77c1}, $\bar u=5$, not a
cube): $1$; \curve{1577a1} at $19$ ($\bar u=8$, a cube but not a ninth power): extra term $1$; \curve{11a1} at $11$, $\ell=5$ ($\bar u=10$, a fifth power): $v_5(\Tam)=2$, although $E[5]$ is
reducible, so only the lemma, not Theorem thm:main, applies there.

\begin{lemma}[the paper's Lemma lem:tamadd]\label{lem:paper:tamadd}
Let $\ell$ be odd and $q\ne\ell$ a prime of additive potentially good reduction. Then $H^1(I_q,\ad T)_{\mathrm{tor}}=0$; in particular $v_\ell(\Tam_q(\ad T))=0$.
\end{lemma}

\noindent\emph{Where proved.} \Cref{ch:blochkato}; inflation-restriction and restriction-corestriction: \Cref{ch:cohomology}; $T_3E$ over $\Z_3[\zeta_3]$: \Cref{ch:ant}.

\begin{proof}[Expanded proof]
Let $M=\ad T$, $K\subseteq I=I_q$ the kernel of the action on $M$, and $\Delta=I/K$, the image of the finite group $\rho(I)$ modulo $\rho(I)\cap\{\pm1\}$ (conjugation kills $\pm1$).
Inflation-restriction: $0\to H^1(\Delta,M)\to H^1(I,M)\to\Hom(K,M)^\Delta$. $K$ is the inertia group of a finite extension of $\Q_q^{\ur}$; a continuous homomorphism $K\to M$ (a pro-$\ell$
group) factors through the maximal pro-$\ell$ abelian quotient of $K$, which is $\Zl(1)$ as $q\ne\ell$; so $\Hom(K,M)\simeq M(-1)$ is torsion-free, and the torsion of $H^1(I,M)$ lies in
$H^1(\Delta,M)$, which is killed by $\#\Delta$. Orders of elements of $\rho(I)$ lie in $\{1,2,3,4,6\}$ and $\#\rho(I)\mid24$, so $\#\Delta$ is a product of $2$s and $3$s: $H^1(\Delta,M)=0$ for
$\ell\ge5$. Let $\ell=3$. Restriction to a Sylow $3$-subgroup $P$ of $\Delta$ followed by corestriction is multiplication by the index, prime to $3$, so the $3$-part of
$H^1(\Delta,M)$ injects into $H^1(P,M)$. As $\#\rho(I)$ divides $24$ ($\rho(I)$ is cyclic of order $1,2,3,4$ or $6$, the quaternion group or $\SL_2(\F_3)$, \Cref{ch:galrep}), a Sylow $3$-subgroup of $\Delta$ has
order $1$ or $3$: $P$ is trivial or generated by the image of an element $\sigma\in\rho(I)$ of order $3$. Then $T$ is free of rank one over $\Z_3[X]/(X^2+X+1)$
with $X$ acting as $\sigma$, so in a suitable basis $\sigma=\sm0{-1}1{-1}$. With $e=\sm0100$: $e$, $\sigma e\sigma^{-1}=\sm0{0}{-1}0$ and $\sigma^2e\sigma^{-2}=\sm{-1}1{-1}1$ have coordinates $(1,0,0)$,
$(0,0,-1)$, $(1,-1,-1)$ in the basis $e,h,f$ (a matrix $\sm b a c{-b}$ has coordinates $(a,b,c)$); the determinant of these three vectors is $-1$, a unit. So $\ad T\simeq\Z_3[P]$ is free over
$\Z_3[P]$ and $H^1(P,\ad T)=0$. Hence $H^1(I,M)_{\mathrm{tor}}=0$, and so are its $G_{\Q_q}$-invariants and their Fitting ideal's valuation.
\end{proof}

\noindent The decisive step at $\ell=3$ is that $\ad T$ is a free $\Z_3[P]$-module: a regular representation has no cohomology. The book checked the three
coordinate vectors and their determinant in PARI. The sentence after the lemma, that $\Tam_q(T)$ can be divisible by $3$ at a type $\kod{IV}$ prime because
$H^1(\langle\sigma\rangle,T)=T/(1-\zeta_3)T\simeq\F_3$, is right: on $T$ the norm $1+\sigma+\sigma^2$ is $0$, so $H^1=\ker N/(\sigma-1)T=T/(1-\zeta_3)T$.

\section[Paper 3: Theorem thm:qual]{Section 3 of the paper: the proof of Theorem thm:qual}\label{sec:paper:qual}

This section reads the paper's Section 3: the congruence number, the congruence criterion (Lemma lem:crit), the oldform lemma (Lemma lem:old), the
optimal-level theorem, the three cases of Theorem thm:qual, the Example and the Remark after it. The method is qualitative: produce a congruent newform
of lower level, and conclude that $\ell$ divides the congruence number, hence the degree.

\begin{definition}[the congruence number; the opening of Section 3]\label{def:paper:congruence-number}
For $E$ optimal with newform $f$ of level $N$, $r_E$ is the largest integer $r$ such that $f\equiv g\pmod r$ coefficientwise for some $g\in S_2(\Gamma_0(N);\Z)$ orthogonal to $f$ for
the Petersson product. Agashe, Ribet and Stein \cite[Theorem 2.2]{ARS12}: $\degphi\mid r_E$, and $v_p(\degphi)=v_p(r_E)$ for every prime $p$ with $p^2\nmid N$.
\end{definition}

\noindent\emph{Where proved.} The theorem of Agashe, Ribet and Stein is stated in \Cref{ch:moddeg} (checked against the source there); the divisibility is reproved
there from the projection formula, and the equality at $\ell\nmid2N$ with $E[\ell]$ irreducible follows from \Cref{prop:paper:pullback} below together with Lemma lem:crit.
The examples where $r_E\ne\degphi$ are computed in \Cref{ch:moddeg,ch:hecke}: \curve{99a1} ($\degphi=4$, $r_E=12$), \curve{726a1} ($144$ and $1584=11\cdot144$), \curve{539d1} ($72$ and
$504=7\cdot72$); in each the gap is at a prime whose square divides $N$.

\begin{lemma}[the paper's Lemma lem:crit]\label{lem:paper:crit}
Let $I=\ker\phi_f$ and $J=\Ann_{\T}(f^\perp)$. Then $\T/(I+J)$ is cyclic of order $r_E$; equivalently $r_E$ generates $\phi_f(J)=\phi_f(\Ann_\T I)$. Consequently, if $g\in S_2(\Gamma_0(N))\otimes\Qbar$
is a normalised eigenform for all of $\T$, orthogonal to $f$, with coefficients in the ring of integers $\mc O$ of a number field, and $a_n(g)\equiv a_n(f)\pmod\lambda$ for all $n$ and a prime
$\lambda\mid\ell$, then $\ell\mid r_E$. The same holds over $\Zl$ with the $\ell$-part of $r_E$.
\end{lemma}

\noindent\emph{Where proved.} In full in \Cref{ch:modgalrep}; restated as a statement about congruence modules in \Cref{ch:congruence}; the perfect pairing
$S_2(\Gamma_0(N);\Z)\times\T\to\Z$ and the Hecke adjoints in \Cref{ch:hecke}; the lattice duality in \Cref{ch:modules}.

\emph{Walk through the paper's proof.} (1) $\T\otimes\Q=\Q\times\T'$ with $\T'$ faithful on $f^\perp$: the line of $f$ and $f^\perp$ are $\T$-stable. For $(n,N)=1$, $T_n$ is self-adjoint
for the Petersson product; for $p\mid N$, $U_p^\dagger=w_NU_pw_N^{-1}$ and $f$ is an eigenvector of $w_N$; so $f$ is an eigenvector of every $T_n^\dagger$, and the orthogonal
complement of a line stable under all $T_n$ and all $T_n^\dagger$ is $\T$-stable (the paper's ``because $f$ is an eigenvector of the adjoints of the $T_n$ as well''). Hence
$I=\T\cap(0\times\T')$, $J=\T\cap(\Q\times0)$, $I\cap J=0$, and $\T/(I+J)\simeq\Z/\phi_f(J)$, of order $r'$ say ($r'$ is finite because $\phi_f(J)\ne0$: $J\otimes\Q=\Q\times0$). (2) Under the perfect
pairing $(g,t)\mapsto a_1(tg)$ the forms killed by $I$ are $\Z f$ and those killed by $J$ are $L=f^\perp\cap S_2(\Gamma_0(N);\Z)$. (3) Dualising $\T\into\T/I\times\T/J$, with cokernel
$\T/(I+J)$, gives $[S_2(\Gamma_0(N);\Z):\Z f\oplus L]=r'$. In detail: the sequence $0\to\T\to\T/I\times\T/J\to\T/(I+J)\to0$, $t\mapsto(t,t)$, $(a,b)\mapsto a-b$, is exact because $I\cap J=0$.
Apply $\Hom(-,\Z)$. The groups $\T/I\simeq\Z$ and $\T/J$ are free ($I$ and $J$ are saturated in $\T$, being intersections of $\T$ with subspaces of $\T\otimes\Q$), and $\Hom(\T/(I+J),\Z)=0$, so
\[
0\to\Hom(\T/I,\Z)\oplus\Hom(\T/J,\Z)\to\Hom(\T,\Z)\to\Ext^1(\T/(I+J),\Z)\to0,
\]
and $\Ext^1(\Z/r',\Z)\simeq\Z/r'$. Under the perfect pairing, $\Hom(\T,\Z)$ is $S_2(\Gamma_0(N);\Z)$ and $\Hom(\T/I,\Z)$, $\Hom(\T/J,\Z)$ are the forms killed by $I$ and by $J$, that is $\Z f$ and $L$ by (2). (4) The quotient embeds in $\Q f/\Z f$ by projection along $f^\perp$, so it is cyclic, generated by a form $h$ with projection
$f/r'$; $r'h=f+g_0$ with $g_0\in L$ gives a congruence modulo $r'$, so $r'\le r_E$; conversely $f\equiv g_1\pmod{r_E}$ with $g_1\in L$ makes $(f-g_1)/r_E$ integral, of order $r_E$ in the
quotient, so $r_E\mid r'$. (5) $\phi_g$ kills $J$ ($g\in f^\perp$) and $\phi_g\equiv\phi_f\pmod\lambda$, so $I+J$ lies in the kernel of $\T\to\mc O/\lambda$, a maximal ideal of residue characteristic
$\ell$, and $\ell\mid r_E$. (6) Every step holds over $\Zl$, the pairing being perfect after $\otimes\Zl$. The hard step is (3): the index of $\Z f\oplus L$ is computed by dualising,
and the cyclicity of $\T/(I+J)$ is what turns the congruence number into a single generator of an ideal.

\begin{lemma}[the paper's Lemma lem:old]\label{lem:paper:old}
Let $h\ne f$ be a newform of weight $2$ on $\Gamma_0(M)$, $M\mid N$, with coefficients in $\mc O$, and $\lambda\mid\ell$ with $a_p(h)\equiv a_p(f)\pmod\lambda$ for all primes $p\nmid\ell N$. Assume $p\nmid M$
for every $p\,\|\,N$ at which $\bar\rho$ is unramified. Then, after enlarging $\mc O$ and $\lambda$, the span of the $h(dz)$, $d\mid N/M$, contains an eigenform $g$ as in Lemma lem:crit; so
$\ell\mid r_E$.
\end{lemma}

\noindent\emph{Where proved.} \Cref{ch:modgalrep}; the $U_p$-structure of old spaces in \Cref{ch:hecke}; Chebotarev and Brauer-Nesbitt in \Cref{ch:galois,ch:reps}; the local
description at $\ell$ (Deligne and Fontaine, as in \cite[Section 2]{Edixhoven92}) and at a multiplicative prime \cite[Theorem 3.1(e)]{DDT} in \Cref{ch:modgalrep,ch:galrep}.

\emph{Walk through.} The span $V$ of the $h(dz)$ is $\T$-stable and orthogonal to $f$. On it $T_n$, $(n,N)=1$, acts by $a_n(h)$, and for $p\mid N/M$ with $e_p=v_p(N/M)$, $U_p$ acts on
the span of $h(p^jz)$, $0\le j\le e_p$, by $U_ph(p^jz)=h(p^{j-1}z)$ ($j\ge1$) and $U_ph=a_p(h)h-ph(pz)$ ($p\nmid M$) or $a_p(h)h$ ($p\mid M$). Its eigenvalues are the roots $\alpha_p,\beta_p$ of
$X^2-a_p(h)X+p$, with $0$ if $e_p\ge2$ ($p\nmid M$), or $a_p(h)$ with $0$ if $e_p\ge1$ ($p\mid M$); each has an eigenvector (the $p$-stabilisations and the depleted forms), though
$U_p$ is not semisimple when $0$ is repeated. As the local spaces tensor together and $\T$ is commutative, any choice of one eigenvalue $\mu_p$ per $p\mid N/M$ gives a
common eigenform $g$; one checks $a_p(g)\equiv a_p(f)$ prime by prime, and then for all $n$ by multiplicativity and the Hecke recursions. The five cases: $p\nmid\ell N$
(hypothesis); $p=\ell$ (both forms have level prime to $\ell$ and weight $2$; $a_\ell\bmod\lambda$ is read off $\bar\rho|_{D_\ell}$, $0$ if irreducible, the eigenvalue of Frobenius on the unramified quotient otherwise, a quotient that is well defined because the characters of the semisimplification restrict to $I_\ell$ as $\bar\chi_\ell$ and $1$, and
$\bar\chi_\ell|_{I_\ell}\ne1$ as $\ell$ is odd, so that it is the module of $I_\ell$-coinvariants);
$p^2\mid N$ ($a_p(f)=0$; $a_p(h)=0$ if $v_p(M)=v_p(N)$ \cite[Theorem 3]{AtkinLehner70}, else $0$ is an available eigenvalue); $p\,\|\,N$ with $\bar\rho$ unramified ($p\nmid M$ by hypothesis,
$\bar\rho(\Frob_p)$ has eigenvalues $a_p(f)p$ and $a_p(f)$, so one of $\alpha_p,\beta_p$ is congruent to $a_p(f)$); $p\,\|\,N$ with $\bar\rho$ ramified ($p\mid M$, and two non-split extensions
of characters match their unique stable lines, so $a_p(h)\equiv a_p(f)$). In the first case the paper adds: ``$\bar\rho_h\simeq\bar\rho$
over a common finite extension of $\F_\ell$, by Chebotarev and Brauer-Nesbitt, $\bar\rho$ being irreducible, hence absolutely irreducible'' (complex conjugation has the
distinct eigenvalues $\pm1$, so an irreducible $\bar\rho$ is absolutely irreducible, and Brauer-Nesbitt then gives an isomorphism, not only of semisimplifications). The last case is the subtle one: the two local representations are matched through their
unique stable lines, which needs the extensions to be non-split, that is, $\bar\rho$ ramified at $p$.

\begin{theorem}[the optimal-level theorem, as the paper uses it]\label{thm:paper:optimal-level}
Let $\ell$ be odd and $\bar\sigma\colon\Gal(\Qbar/\Q)\to\GL_2(\Fbar_\ell)$ irreducible and modular, absolutely irreducible on $\Gal(\Qbar/\Q(\sqrt{-3}))$ if $\ell=3$. Then $\bar\sigma$ arises from a
newform of weight $2$ and level $N(\bar\sigma)\ell^{\delta(\bar\sigma)}$ with nebentypus of order prime to $\ell$, $N(\bar\sigma)$ the Serre conductor and $\delta(\bar\sigma)=0$ when $\bar\sigma|_{D_\ell}$ is good
\cite[Theorem 3.15]{DDT} (Ribet, Carayol, Diamond \cite{Ribet90,Carayol89,Diamond95}). For $\bar\sigma=\bar\rho_{E',\ell}$ with $E'$ of conductor prime to $\ell$ it arises from a newform on
$\Gamma_0(N(\bar\sigma))$; and $N(\bar\rho_{E',\ell})\mid N(E')$.
\end{theorem}

\noindent\emph{Where.} Stated with references and sketched in \Cref{ch:modgalrep}; we do not prove it. The deduction of trivial character is the paper's: $\det\bar\sigma=\bar\chi_\ell$ forces
the nebentypus $\epsilon$ to reduce to $1$; reduction is injective on roots of unity of order prime to $\ell$, so $\epsilon=1$. The theorem number is as cited by the paper.

\begin{theorem}[the paper's Theorem thm:qual]\label{thm:paper:qual}
Let $E/\Q$ have conductor $N$ and $\ell$ an odd prime, $\ell\nmid N$, with $\bar\rho_{E,\ell}$ irreducible and, if $\ell=3$, absolutely irreducible on $\Gal(\Qbar/\Q(\sqrt{-3}))$. Let $q\ne\ell$ be a
prime of additive reduction of one of the kinds: (a) type $\kod{IV}$ or $\kod{IV}^*$, $q\ge5$ or $q=2$, $\ell=3$; (b) type $\kod{II}$ or $\kod{II}^*$, $q\ge5$, $\ell=3$; (c) type $\kod{I}_n^*$ with $n\ge1$ and $\ell\mid n$,
$q$ odd. Then $\ell$ divides the modular degree of the optimal curve of the class of $E$.
\end{theorem}

\noindent\emph{Where proved.} \Cref{ch:modgalrep}, with the inputs \Cref{lem:paper:IV,lem:paper:twist,lem:paper:tamagawa-exp,lem:paper:old}, \Cref{thm:paper:optimal-level} and
\cite[Theorem 2.2]{ARS12}.

\begin{proof}[Expanded proof]
Let $E_0$ be the optimal curve and $r=r_{E_0}$. As $\ell\nmid N$, $v_\ell(\deg\varphi_{E_0})=v_\ell(r)$, so it suffices to find $h$ as in \Cref{lem:paper:old}. In all three cases $q^2\,\|\,N$: the
conductor exponent at $q$ is $2$ (\Cref{lem:paper:IV} in case (a), \Cref{lem:paper:twist}(ii) in cases (b) and (c)). Notation: $h^\chi$ is the newform twist and $h\otimes\chi$ the naive
twist; when $q\,\|\,\mathrm{level}(h)$ or $q\nmid\mathrm{level}(h)$, $h\otimes\chi$ is the newform $h^\chi$ (\Cref{lem:paper:twist}(iii)); for $h=f_E$ with $q^2\,\|\,N$ and $E^{(q)}$ also additive at $q$ (case (b)), $f\otimes\chi$ is the
newform of $E^{(q)}$, as both have $a_q=0$; and $(g\otimes\chi)\otimes\chi=g^{\{q\}}$ is the $q$-depletion of $g$.

\emph{Case (a).} By \Cref{lem:paper:IV}, the Serre conductor $M=N(\bar\rho)$ has $v_q(M)=1$; as $M\mid N$ (Carayol), $M\mid N/q$. \Cref{thm:paper:optimal-level} gives a newform $g$ on $\Gamma_0(M)$ with
$\bar\rho_g\simeq\bar\rho$, so $a_p(g)\equiv a_p(f)\pmod\lambda$ for $p\nmid3N$. Take $h=g$: $h\ne f$ since the levels differ, and $p\nmid M$ for every $p$ at which $\bar\rho$ is unramified. At $p=q$ the
eigenvalue $0$ of $U_q$ is available ($v_q(M)=1<2$). \Cref{lem:paper:old} gives $3\mid r$.

\emph{Case (b).} Let $\chi=\chi_{q^*}$, $E'=E^{(q)}$: type $\kod{IV}^*$ or $\kod{IV}$ at $q$, the same types elsewhere, conductor $N$. Its newform is $f'=f\otimes\chi$: both $f'$ and $f\otimes\chi$ have
$a_q=0$ (additive reduction) and $a_p(f')=\chi(p)a_p(f)$ for $p\ne q$; this is not Atkin-Li, since $q^2\mid N$. The representation $\bar\rho\otimes\chi=\bar\rho_{E',3}$ satisfies the hypotheses
of \Cref{thm:paper:optimal-level} (irreducibility and absolute irreducibility over $\Q(\sqrt{-3})$ survive a twist). By \Cref{lem:paper:IV} for $E'$, $M=N(\bar\rho\otimes\chi)$ has $v_q(M)=1$,
$M\mid N/q$; there is a newform $g$ on $\Gamma_0(M)$ with $a_p(g)\equiv\chi(p)a_p(f)\pmod\lambda$ for $p\nmid3N$. As $q\,\|\,M$, $h=g\otimes\chi$ is a newform on $\Gamma_0(Mq)$, $Mq\mid N$, with
$a_p(h)=\chi(p)a_p(g)\equiv\chi(p)^2a_p(f)=a_p(f)$. If $h=f$, then $a_p(g)=\chi(p)a_p(h)=a_p(f')$ for all $p\nmid qN$, so $g=f'$ by strong multiplicity one (\Cref{ch:hecke}), impossible as
their levels are $M<N$ and $N$. (Here $h\otimes\chi$ is the $q$-depletion of $g$, not $g$, since $a_q(g)=\pm1$.) If $p\,\|\,N$ and $\bar\rho$ is unramified at $p$, then $p\ne q$, so $\bar\rho\otimes\chi$ is
unramified at $p$ and $p\nmid Mq$. \Cref{lem:paper:old} gives $3\mid r$.

\emph{Case (c).} $E'=E^{(q)}$ has type $\kod{I}_n$ at $q$ and conductor $N/q$; $f'=f_{E'}$ has level $N/q$, $q\,\|\,N/q$, $a_q(f')=\pm1$ and $a_p(f')=\chi(p)a_p(f)$ for $p\ne q$, so $f\otimes\chi$ is
the $q$-depletion of $f'$, not $f'$ (for example \curve{539c1} twisted by $\chi_{-7}$ lies in \curve{77c}, and $a_7(\curve{77c1})=-1$). By \Cref{lem:paper:tamagawa-exp} for $E'$, $\ell\mid n$ makes $\bar\rho\otimes\chi$
unramified at $q$, so $M=N(\bar\rho\otimes\chi)$ is prime to $q$ and $M\mid N/q^2$; Ribet's theorem (the case of \Cref{thm:paper:optimal-level} with $q$ exactly dividing the level of $E'$ and
$\bar\rho\otimes\chi$ unramified at $q$) gives a newform $g$ on $\Gamma_0(M)$ with $a_p(g)\equiv\chi(p)a_p(f)$ for $p\nmid\ell N$. As $q\nmid M$, $h=g\otimes\chi$ is a newform on $\Gamma_0(Mq^2)$, $Mq^2\mid N$,
congruent to $f$ at $p\nmid\ell N$. If $h=f$, then $g$ and $f'$ have equal $a_p$ for $p\nmid qN$, so $g=f'$, impossible since their levels $M$ and $N/q$ differ. The hypothesis of
\Cref{lem:paper:old} holds as in case (b), and $\ell\mid r$.
\end{proof}

\noindent\emph{Reading guidance.} The step to watch in cases (b) and (c) is $h\ne f$: it is proved by comparing $g$ with $f'$ through strong multiplicity one and their
levels, and the distinction between the naive twist $f\otimes\chi$ and the newform of $E'$ matters there (in case (c), $f\otimes\chi$ is the $q$-depletion of $f'$). Why $\kod{I}_0^*$, and $\kod{I}_n^*$ with $\ell\nmid n$, are not in the
theorem: after the twist the lowered form may be $f_{E'}$ itself, so $h=f$; Theorem thm:quant reaches these terms by raising the level instead.

\begin{example}[the paper's Example ex:qual, verified]\label{ex:paper:qual}
The paper's congruences, each between two forms in the same $S_2(\Gamma_0(N);\Z)$, so that Sturm's bound $\frac16N\prod_{p\mid N}(1+\frac1p)$ \cite{Sturm87} (\Cref{ch:modforms}) proves them
for all $n$. (a) at $q=2$: \curve{52a1}, type $\kod{IV}^*$ at $2$, Serre conductor $26$ at $\ell=3$; $f_{52a}\equiv f_{26b}(z)-f_{26b}(2z)\pmod3$ (Sturm bound $14$). Here $2\,\|\,26$ and $a_2(\curve{26b1})=1$, so
$f_{26b}(z)-a_2f_{26b}(2z)$ is the $U_2$-eigenvector of eigenvalue $0$ in Lemma lem:old; $\deg\varphi_{\curve{52a1}}=3$ and $e_3(E,2)=1$. (b) \curve{121a1}, type $\kod{II}$ at $11$ ($v_{11}(\Delta)=2$); its twist
by $\chi_{-11}$ is \curve{121c2}, of type $\kod{IV}^*$ ($v_{11}(\Delta)=8$ on the minimal model; the optimal curve \curve{121c1} of that class has type $\kod{IV}$), whose mod $3$ representation has
Serre conductor $11$ and comes from $f_{11a}$; twisting back, $f_{11a}\otimes\chi_{-11}=f_{121d}$, and $f_{121a}\equiv f_{121d}\pmod3$ (Sturm bound $22$); $\degphi=6$, $e_3(E,11)=v_3(12)=1$. (c) at $\ell=3$:
\curve{539c1}, type $\kod{I}_3^*$ at $7$; its twist by $\chi_{-7}$ is \curve{77c1} (type $\kod{I}_3$, conductor $77$); $f_{539c}\equiv f_{11a}\otimes\chi_{-7}=f_{539d}\pmod3$ (Sturm bound $112$); $\degphi=2^5\cdot3^2$,
and $e_3(E,7)=v_3(48)=1$, $t_3(E,7)=1$ are both attained. (c) at $\ell=5$: \curve{980i1}, type $\kod{I}_5^*$ at $7$; its twist by $\chi_{-7}$ is \curve{140b1}; $f_{980i}\equiv f_{20a}\otimes\chi_{-7}=f_{980g}\pmod5$
(Sturm bound $336$); $\degphi=2^6\cdot3^2\cdot5$. Since $5\mid980$, the last illustrates the congruence, not the theorem, as the paper says. The book verified each congruence to
$n\le400$ or $n\le1000$, each twist identification, each degree, and that PARI's \texttt{mfsturm} returns $14$, $22$, $112$, $336$ (Section~\ref{sec:paper:verifications}). The level-raising
conditions are visible: $a_7(\curve{11a1})=-2$ and $(7+1)^2-4=60\equiv0\pmod3$; $a_7(\curve{20a1})=2$ and $60\equiv0\pmod5$.
\end{example}

\noindent Sturm's bound is what makes the congruences proofs rather than observations: both forms lie in the same space $S_2(\Gamma_0(N))$ with integer coefficients. The class \curve{980g}, the twist of \curve{20a}, must not be confused with \curve{980i}.

\paragraph{The Remark after the Example.} It says that $\ell\nmid N$ is used in the optimal-level theorem, in Lemma lem:old at $p=\ell$, and to pass from $r_E$ to $\degphi$, the last
needing only $\ell^2\nmid N$, and gives \curve{99a1} ($N=9\cdot11$, $\degphi=4$, $r_E=12$; \cite[Table 1]{ARS12}), \curve{726a1} and \curve{539d1} (above). All three were recomputed in
\Cref{ch:moddeg,ch:hecke}; at \curve{99a1} and $\ell=3$, moreover, $H^1(X_0(99),\Z_3)_{\mf m}$ is not free over $\T_{\mf m}$ (\Cref{ch:moddeg}). The last sentence, that the data show the conclusion
``at $q=3$ in cases (a) and (b), where $q=\ell$ and the theorem does not apply'', names the obstruction: at $q=\ell$ the optimal-level theorem and Lemma lem:old at $p=\ell$ are not available.

\section[Paper 4.1: congruence ideals]{Section 4.1 of the paper: congruence ideals and the modular degree}\label{sec:paper:congruence-ideals}

This section reads the paper's Section 4.1: the objects of Diamond, Flach and Guo for weight $2$, the comparison of their integral structure with the
cohomology of $X_0(M^\Sigma)$ (Proposition prop:integral), their formula for the variation of the congruence ideal with $\Sigma$ (Theorem thm:dfg14), the identification
of the congruence ideal at level $N$ with the modular degree (Proposition prop:pullback), the commutative algebra lemma (Lemma lem:ddt417), Lemma lem:hecke and Remark
rem:free. This is Step 1 of the proof of Theorem thm:main, and the place where no multiplicity-one hypothesis is needed. Section 3 used congruences one at a time, to show that $\ell$ divides $r_E$; from here on the paper measures
the whole congruence module, as an ideal $\eta$ of $\Zl$, because that is the quantity that Diamond, Flach and Guo relate to a Selmer group.

\begin{definition}[the objects of Diamond, Flach and Guo in weight $2$; Section 4.1]\label{def:paper:dfg-objects}
Let $f$ be a newform of weight $2$ on $\Gamma_0(M)$ with rational coefficients, $\lambda=\ell\nmid2M$, $\Sigma$ a finite set of primes, $M^\Sigma=M\prod_{p\in\Sigma}p^{\delta_p}$ and
$f^\Sigma=\sum_{(n,\Sigma)=1}a_n(f)q^n$, an eigenform of the full Hecke algebra of level $M^\Sigma$ with $U_pf^\Sigma=0$ for $p\in\Sigma$. The integral structure $M(M^\Sigma,1)_{!,\lambda}$ is the image of
the $U_0(M^\Sigma)$-invariants of the compactly supported cohomology of a modular curve of full level in its cohomology \cite[Sections 1.2, 1.4]{DFG04}. The kernel lattice
$M_{f^\Sigma,\lambda}$ is the subgroup killed by the kernel of the eigencharacter of $f^\Sigma$ on the full Hecke algebra. The pairing $\hat\delta_!$ is the cup product composed with
$[UwU]_\omega$, which in weight $2$ with trivial character is the Atkin-Lehner involution $w_{M^\Sigma}$; on classes pulled back from $X_0(M^\Sigma)$ it is $\frac u2\langle x,w^*_{M^\Sigma}y\rangle$ for an $\ell$-adic unit $u$.
The naive $\Sigma$-finite congruence ideal $\eta^\Sigma_{f,\lambda}\subseteq\Zl$ is the ideal of values of $\hat\delta_!$ on $M_{f^\Sigma,\lambda}$, the image of $\wedge^2M_{f^\Sigma,\lambda}$; equivalently the Gram
determinant on a basis is $(\eta^\Sigma_{f,\lambda})^2$ up to a unit.
\end{definition}

\noindent\emph{Where proved.} \Cref{ch:congruence}, where the definitions are compared with the published text of \cite{DFG04} (Sections 1.5.3, 1.6.1, 1.6.2, 1.7.3, 1.8.2,
displays (13), (14)); the cup product, $w_N$ and the Hecke adjoints in \Cref{ch:curves,ch:moddeg}; a rank-two alternating form has Gram matrix $\sm0a{-a}0$, so its Gram
determinant is the square of a generator of its ideal of values (\Cref{ch:modules}).

\begin{proposition}[the paper's Proposition prop:integral]\label{prop:paper:integral}
Let $\ell$ be a prime. Pull-back along the covering of $X_0(M^\Sigma)$ by the modular curve of full level identifies $H^1(X_0(M^\Sigma),\Z)$ with the Betti lattice
$M(M^\Sigma,1)_{!,B}$ of \cite[Section 1.4.2]{DFG04}, compatibly with the Hecke operators, $U_p$ included; in particular $M(M^\Sigma,1)_{!,\lambda}=H^1(X_0(M^\Sigma),\Zl)$, with no
localisation.
\end{proposition}

\noindent\emph{Where proved.} In the paper, in four steps; the book's expanded proof is the comparison proposition of \Cref{ch:congruence}, which
also explains equivariant cohomology with compact support, the Hochschild-Serre spectral sequence and the equivariant Leray spectral sequence that the paper uses
without comment. Riemann surfaces, their compactly supported cohomology and the cup product: \Cref{ch:curves}; the Eichler-Shimura isomorphism and the Hecke
action on $H^1(X_0(N),\Z)$: \Cref{ch:moddeg}; group cohomology of finite groups: \Cref{ch:cohomology}.

\begin{proof}[The paper's proof, read]
Write $M$ for $M^\Sigma$ and $U=U_0(M)$. In weight $2$ with trivial character the coefficient sheaf of \cite[Section 1.2.2]{DFG04} is the constant sheaf $\Z$. The paper
chooses an auxiliary $N\ge3$ with $M\mid N$ and every prime factor of $N$ in the set $S$ of \cite[Section 1.4.1]{DFG04}, for instance $N=4M$. Then $U_N\subseteq U$, and by
\cite[Section 1.4.1]{DFG04} $M(M,1)_{?,B}=H^1_?(Y(N),\Z)^G$ for $?\in\{c,\mathrm{tf}\}$, where $Y(N)$ is the analytic modular curve of full level $N$, a disjoint union of copies of
$\Gamma(N)\backslash\mathcal H$ indexed by $(\Z/N\Z)^\times$, and $G=U/U_N$, a finite group acting through its image in $\GL_2(\Z/N\Z)$; the lattice $M(M,1)_{!,B}$ is the image of the
compactly supported invariants in the torsion-free invariants, and nothing depends on $N$ \cite[Lemma 1.1]{DFG04}. Since $\det U=\widehat\Z^\times$, $G$ permutes the
components transitively, the stabiliser of one component acts on it through the image of $\Gamma_0(M)$, and the quotient map $\pi\colon Y(N)\to Y(N)/G=Y_0(M)$ is a finite map
onto the coarse curve; point stabilisers are finite.

(1) \emph{Compact support and invariants.} $\pi^*\colon H^1_c(Y_0(M),\Z)\to H^1_c(Y(N),\Z)^G$ is an isomorphism. The paper passes through the $G$-equivariant cohomology with
compact support $H^\bullet_{G,c}(Y(N),\Z)$ and two spectral sequences converging to it. In the Hochschild-Serre spectral sequence
$H^p(G,H^q_c(Y(N),\Z))\Rightarrow H^{p+q}_{G,c}(Y(N),\Z)$ the row $q=0$ is zero, because each component is a connected non-compact surface, so $H^0_c(Y(N),\Z)=0$; in degree $1$
only $E_2^{0,1}=H^1_c(Y(N),\Z)^G$ survives. In the equivariant Leray spectral sequence of $\pi$, $H^p_c(Y_0(M),\mathcal R^q)\Rightarrow H^{p+q}_{G,c}(Y(N),\Z)$, the stalk of $\mathcal R^q$ at $y$ is
$H^q(G_x,\Z)$ for a point $x$ over $y$ (Shapiro's lemma); so $\mathcal R^0=\Z$ and $\mathcal R^1=0$, because $H^1(C,\Z)=\Hom(C,\Z)=0$ for a finite group $C$. In degree $1$ only
$H^1_c(Y_0(M),\Z)$ survives. The composite of the two edge isomorphisms is $\pi^*$.

(2) \emph{The image.} Under these identifications DFG's map from compactly supported to torsion-free invariants is the forgetful map $H^1_c(Y(N),\Z)^G\to H^1(Y(N),\Z)^G$; as $\pi$ is
proper, its image is $\pi^*$ of the image of $H^1_c(Y_0(M),\Z)\to H^1(Y_0(M),\Z)$. Since $Y_0(M)$ is the compact surface $X_0(M)$ minus finitely many cusps, $H^1_c(Y_0(M),\Z)\to H^1(X_0(M),\Z)$ is
onto and $H^1(X_0(M),\Z)\to H^1(Y_0(M),\Z)$ is injective, so that image is $H^1(X_0(M),\Z)$. Finally $\pi^*$ is injective on the torsion-free group $H^1(Y_0(M),\Z)$, since $\pi_*\pi^*$ is
multiplication by $\deg\pi$.

(3) \emph{Hecke operators.} Both lattices are torsion-free and $\pi^*$ is injective, so equivariance may be checked after $\otimes\C$. There DFG's operators act on
$\operatorname{Fil}^1M(M,1)_{!,\mathrm{dR}}\otimes\C=S_2(\Gamma_0(M))$ as the classical $T_p$, and as $U_p$ for $p\mid M$ \cite[Proposition 1.3]{DFG04}, commuting with complex conjugation; $\pi^*$ is
compatible with DFG's comparison isomorphism and with the Eichler-Shimura isomorphism, both being pull-back of $f(\tau)\,d\tau$.

(4) \emph{The $\lambda$-adic structure.} $\Zl\otimes M_{?,B}\simeq M_{?,\ell}$ for every prime $\ell$ \cite[Section 1.2.3]{DFG04}, and invariants under a finite group and images commute
with the flat base change $\Z\to\Zl$; so $M(M,1)_{!,\lambda}=\Zl\otimes M(M,1)_{!,B}=H^1(X_0(M),\Zl)$.
\end{proof}

\noindent\emph{Reading guidance.} Step (1) is the one to master, and its point is compact support. Both vanishing statements are elementary: $H^0_c$ of a connected
non-compact surface is $0$ (a compactly supported locally constant function on it is $0$), and a finite group has no nonzero homomorphism to $\Z$. Together they say that
taking $G$-invariants of $H^1_c$ loses nothing and that the finite stabilisers (the elliptic points) contribute nothing in degree $1$. Step (2) is the familiar
description of $H^1(X_0(M),\Z)$ as the interior cohomology of $Y_0(M)$ (\Cref{ch:moddeg}). Steps (3) and (4) are bookkeeping: equivariance between torsion-free lattices can
be tested over $\C$, and Betti-to-$\ell$-adic comparison is a base change. The statement holds for every prime $\ell$, $\ell=2$ and $\ell\mid M^\Sigma$ included; the
hypothesis $\ell\nmid2M^\Sigma$ of the surrounding text is needed for the pairing, not for the comparison.

\emph{Reading remark} (the paper's Remark after the proposition). ``Compact support is what makes the comparison integral'': with ordinary cohomology the analogue
of step (1) fails, because $H^0(Y(N),\Z)\ne0$. The paper identifies the cokernel: $H^1(Y(N),\Zl)^G/\pi^*H^1(Y_0(M),\Zl)$ is the group of characters $\Gamma_0(M)\to\Qlmod$ of $\ell$-power order that
factor through $G$, namely $\gamma\mapsto\psi(d\bmod p)$ for primes $p\mid M$ with $\ell\mid p-1$; ``for $M=11$ and $\ell=5$ it is the unramified covering $X_1(11)\to X_0(11)$''. These
characters cut out abelian unramified subcoverings such as the Shimura covering, on which the Hecke operators act through an Eisenstein ideal
\cite{LingOesterle,Ribet88}; the proof does not pass through them, since the compactly supported lattice never sees this cokernel. Proposition prop:pullback (last
sentence), Lemma lem:hecke, Step 1 of the proof of Theorem thm:main and Step 2 of Theorem thm:quant use the proposition.

\paragraph{The pairing.} The operator in $\hat\delta_!$ is $[UwU]_\omega$, attached to $w=\sm0{-1}{M^\Sigma}0$ \cite[Section 1.5.3, (13), (14)]{DFG04}; in weight $2$ with
trivial character it is the Atkin-Lehner involution with no scalar factor. On classes pulled back from $X_0(M^\Sigma)$, $\hat\delta_!$ is $\frac u2$ times the cup product
on $X_0(M^\Sigma)$ composed with $w_{M^\Sigma}$, for an $\ell$-adic unit $u$, ``so for odd $\ell$ only a unit separates the two''. The paper's justification: for a sufficiently small
level $U'\subseteq U_0(M^\Sigma)$ with $\det U'=\widehat\Z^\times$, such as $U_0(M^\Sigma)\cap U_1(4)$, the normalised pairing $\bar\delta$ of \cite[Section 1.5.2]{DFG04} is Poincar\'e duality on $X_{U'}$
up to an $\ell$-adic unit when $\ell\nmid2M^\Sigma$; it is normalised by dividing by the index, so on $U_0(M^\Sigma)$-invariants it is $[U_0(M^\Sigma):U']^{-1}$ times its value on
$U'$-invariants; and the covering $X_{U'}\to X_0(M^\Sigma)$ has degree $[U_0(M^\Sigma):U']/2$, because $-1$ acts trivially. By the projection formula the two factors combine to $\frac12$,
up to the unit $u$. Only the ideal of values of the pairing enters the congruence ideal, so the unit $u$ never matters. Self-adjointness of the Hecke operators, $U_p$ included, for $\hat\delta_!$ comes from
$w_{M^\Sigma}T_nw_{M^\Sigma}^{-1}$ being the adjoint of $T_n$ for the cup product (\cite[Section 1.8.2, (25)]{DFG04}; \Cref{ch:moddeg}). This is the reason for composing with $w$: the
plain cup product is alternating and perfect, but the $U_p$ are not self-adjoint for it.

\begin{theorem}[Diamond, Flach and Guo, Proposition 1.4(c); the paper's Theorem thm:dfg14]\label{thm:paper:dfg14}
Let $\ell\nmid2M$ and $\ell\notin\Sigma$. If $\bar\rho_f$ is irreducible, then $\eta^\Sigma_{f,\lambda}=\eta^\emptyset_{f,\lambda}\prod_{p\in\Sigma}L^{\nv}_p(A_f,1)^{-1}$.
\end{theorem}

\noindent\emph{Where.} Stated with the reference \cite[Proposition 1.4(c)]{DFG04} (verified), and sketched in \Cref{ch:congruence} from parts (a) and (b), given the integral
comparison isomorphism of \cite{DFG04}, which is cited. We do not prove it here. The three factors (\Cref{ch:congruence}): an additive $p$ contributes $1$ (neither level nor
ideal changes); a good $p$ contributes $(1-\frac\alpha\beta p^{-1})(1-p^{-1})(1-\frac\beta\alpha p^{-1})=\frac{p-1}p\cdot\frac{(p+1)^2-a_p^2}{p^2}$, Wiles's level-raising factor; a multiplicative $p$ contributes
$1-p^{-2}$.

\noindent\emph{Reading guidance.} The hypothesis $\ell\notin\Sigma$ cannot be dropped, and the paper's paragraph after the theorem gives the example: for $f=f_{11a}$, $\ell=3$,
$\Sigma=\{3\}$, $a_3=-1$ and $L^{\nv}_3(A_f,1)^{-1}=\frac23\cdot\frac{16-1}9=\frac{10}9$, of $3$-adic valuation $-2$, so the right side would be $3^{-2}\Zl$ while $\eta^\Sigma_{f,\lambda}$ is an
ideal of $\Zl$ (\Cref{ex:paper:dfg14-ell} gives a second instance). In \cite[Section 1.7.3]{DFG04} $\ell\notin\Sigma$ is a standing assumption, and every use in the paper has it
($\Sigma$ is a set of primes dividing the level, and $\ell$ is prime to the level).

\begin{proposition}[the paper's Proposition prop:pullback]\label{prop:paper:pullback}
Let $E$ be $X_0(N)$-optimal with newform $f$, $\varphi=\varphi_E$, $I_f=\ker\phi_f$, and $\ell$ any prime. Then $\varphi^*\colon H^1(E,\Zl)\to H^1(X_0(N),\Zl)$ is injective with image the kernel lattice
$H^1(X_0(N),\Zl)[I_f]$, and $\langle\varphi^*a,w_N^*\varphi^*b\rangle_{X_0(N)}=\varepsilon_N(f)\,\degphi\,\langle a,b\rangle_E$, with $\varepsilon_N(f)=\pm1$ the eigenvalue of $w_N$ on $f$. For $\ell$ odd the ideal of values of
$\frac12\langle x,w_N^*y\rangle$ on $H^1(X_0(N),\Zl)[I_f]$ is $\degphi\,\Zl$ and its Gram determinant is $(\degphi)^2$ up to a unit. If $\ell\nmid2N$, then
$M_{f,\lambda}=\varphi^*H^1(E,\Zl)$ and $v_\lambda(\eta^\emptyset_{f,\lambda})=v_\ell(\degphi)$.
\end{proposition}

\noindent\emph{Where proved.} \Cref{ch:moddeg} (the image and the projection formula), \Cref{ch:curves} (cup product, pullback and degree), \Cref{ch:congruence} (the
identification with $\eta^\emptyset_{f,\lambda}$, given the comparison of \Cref{prop:paper:integral}).

\begin{proof}[Expanded proof]
(1) \emph{Injectivity and saturation.} By optimality the kernel $B$ of $J_0(N)\to E$ is connected; the long exact sequence of $0\to B\to J_0(N)\to E\to0$ in homology ends with
$H_1(J_0(N),\Z)\to H_1(E,\Z)\to\pi_0(B)=0$, so $\varphi_*$ is onto on $H_1$, and its dual $\varphi^*\colon H^1(E,\Z)\to H^1(X_0(N),\Z)$ is injective with torsion-free cokernel; tensor with $\Zl$.
(2) \emph{The image is the kernel lattice.} $\varphi^*\omega_E=c_E\cdot2\pi i\,f(z)\,dz$, so $\varphi^*H^1(E,\C)=\C f\,dz\oplus\C\overline{f\,dz}$. This plane is $H^1(X_0(N),\C)[I_f]$, because $f$ is new: the
$\phi_f$-eigenspace in $S_2(\Gamma_0(N))$ is $\C f$. The plane is defined over $\Q$: $H^1(X_0(N),\Q)[I_f]$ is a two-dimensional $\Q$-form of it, and by flatness
$H^1(X_0(N),\Zl)[I_f]=H^1(X_0(N),\Z)[I_f]\otimes\Zl$ is a saturated lattice of rank $2$ in $H^1(X_0(N),\Ql)[I_f]$. The image of $\varphi^*$ is a saturated rank-two lattice in the same
$\Ql$-plane, so the two coincide. (Two saturated lattices of the same rank in the same $\Ql$-plane are equal; this is the point where optimality,
through the torsion-free cokernel of step (1), enters.) (3) \emph{Projection formula.} $\varphi^*(a\cup b)=\varphi^*a\cup\varphi^*b$ and $\varphi^*$ multiplies the fundamental
class of $H^2$ by $\degphi$, so $\langle\varphi^*a,\varphi^*b\rangle_{X_0(N)}=\degphi\langle a,b\rangle_E$. (4) \emph{The involution.} $w_N$ preserves the plane: by multiplicity one for newforms the
plane is also the eigenspace of the $T_n$ with $(n,N)=1$ for the eigenvalues $a_n(f)$, and $w_N$ commutes with those $T_n$; it does not commute with the $U_p$, $p\mid N$, which is why the
argument passes through the $T_n$ with $(n,N)=1$, as the paper does. On the plane $w_N$ acts by
$\varepsilon_N(f)=\pm1$, which gives the displayed formula. (5) \emph{Values.} The cup product on $H^1(E,\Z)$ is unimodular, so the ideal of values of $\frac12\langle x,w_N^*y\rangle$ on the
image is $\degphi\Zl$ ($\frac12$ and $\varepsilon_N$ are units), and the Gram determinant is its square. (6) \emph{DFG's ideal.} If $\ell\nmid2N$, DFG's integral structure is $H^1(X_0(N),\Zl)$, Hecke equivariantly with the $U_p$ included
(\Cref{prop:paper:integral}), so $M_{f,\lambda}=H^1(X_0(N),\Zl)[I_f]=\varphi^*H^1(E,\Zl)$ by (2); and $\hat\delta_!$ is $\frac u2\langle x,w_N^*y\rangle$ with $u$ an $\ell$-adic unit (the paragraph on the
pairing above). By (5), with the unit $u$ changing nothing, $\eta^\emptyset_{f,\lambda}=\degphi\Zl$.
\end{proof}

\noindent\emph{Reading guidance.} The proposition is the heart of the identification of $\degphi$ with a congruence ideal, and it uses no multiplicity-one
hypothesis, and no irreducibility of $\bar\rho$: the paper notes that ``for the first part no irreducibility'' is used, and the last sentence needs none either,
because \Cref{prop:paper:integral} holds without localisation. Only the last sentence uses Proposition prop:integral. The consequence the paper draws after the proof,
$M_{f,\lambda}=H^1(E,\Zl)=\Hom(T_\ell E,\Zl)\simeq T_\ell E(-1)$ by Galois equivariance of $\varphi^*$ in \'etale cohomology, so that its adjoint is $\ad T_\ell E$, is what lets
the Selmer groups of \cite{DFG04} be read as Selmer groups of $\ad T_\ell E$: $M_{f,\lambda}$ lives in the dual representation, of determinant $\chi_\ell^{-1}$, but the adjoint
does not see duals or twists (\Cref{conv:paper:duals}).

\begin{lemma}[the paper's Lemma lem:ddt417, the inequality half of {\cite[Lemma 4.17]{DDT}}]\label{lem:paper:ddt417}
Let $\mc O$ be a discrete valuation ring with finite residue field, $\T$ a commutative $\mc O$-algebra, finite and free as an $\mc O$-module, acting on an $\mc O$-module $L$ free of finite rank
with a perfect $\mc O$-bilinear pairing for which $\T$ is self-adjoint; $\pi\colon\T\to\mc O$ a surjective homomorphism, $\wp=\ker\pi$, $I=\Ann_\T(\wp)$, $\eta=\pi(I)$; and suppose $L[\wp]$ free of rank $d$
with Gram determinant $\Delta\ne0$. Then $\eta^d\subseteq\Delta\mc O$, that is $d\,v(\eta)\ge v(\Delta)$, with equality if $L$ is free over $\T$.
\end{lemma}

\noindent\emph{Where proved.} \Cref{ch:congruence} (in length form, over any discrete valuation ring); saturation, the dual map and Gram determinants: \Cref{ch:modules}.

\begin{proof}[Expanded proof]
If $\eta=0$ there is nothing to prove. Otherwise take $t\in I$ with $\pi(t)$ a generator of $\eta$.
(1) $\wp\cap I=0$: for $x\in\wp\cap I$, $tx=0$ (as $t$ kills $\wp$) and $(\pi(t)-t)x=0$ (as $\pi(t)-t\in\wp$ and $x$ kills $\wp$), so $\pi(t)x=0$ and $x=0$, $\T$ being torsion-free. Reducedness of $\T$
is not the reason: the full Hecke algebra can have nilpotents ($U_2$ at level $88$ on the span of the $f_{11a}(dz)$, $d\mid8$; \Cref{ch:hecke}).
(2) $I=\mc Ot$: $\pi|_I$ is injective by (1), with image $\eta=\pi(t)\mc O$.
(3) $L[\wp]\cap L[I]=0$: an element killed by $\wp$ and by $I$ is killed by $\pi(t)=t+(\pi(t)-t)\in I+\wp$, a nonzero scalar, so it is $0$, $L$ being torsion-free. This is the paper's argument; note that the
torsion-freeness of $L$ alone would not suffice, the nonzero scalar $\pi(t)$ in $I+\wp$ is what is used.
(4) $Q=L/(L[\wp]\oplus L[I])$ is killed by $\eta$: for $x\in L$, $tx\in L[\wp]$ (as $\wp I=0$) and $(\pi(t)-t)x\in L[I]$ (for $s\in I$, $s(\pi(t)-t)\in I\cap\wp=0$), and their sum is $\pi(t)x$.
(5) $L[\wp]^\perp=L[I]$: if $y\in L[I]$, $x\in L[\wp]$, then $\pi(t)\langle x,y\rangle=\langle tx,y\rangle=\langle x,ty\rangle=0$. Conversely, if $y\perp L[\wp]$, then for every $z$, $\langle ty,z\rangle=\langle y,tz\rangle=0$ as
$tz\in L[\wp]$; so $ty=0$, and by (2) $Iy=0$.
(6) $L[\wp]$ is saturated, so the pairing maps $L$ onto $\Hom(L[\wp],\mc O)$ with kernel $L[\wp]^\perp=L[I]$; hence $Q$ is the cokernel of $L[\wp]\to\Hom(L[\wp],\mc O)$, $x\mapsto\langle x,\cdot\rangle$. In a basis of $L[\wp]$ and the dual basis this map has the Gram matrix as its matrix, and the
cokernel of a square matrix over a discrete valuation ring with nonzero determinant has length the valuation of the determinant (Smith normal form, \Cref{ch:modules}): length $v(\Delta)$. $Q$ is
generated by $d$ elements (a quotient of $\Hom(L[\wp],\mc O)\simeq\mc O^d$) and killed by $\eta$, so its length is at most $d\,v(\eta)$.
(7) If $L\simeq\T^d$, then $Q\simeq(\T/(\wp+I))^d\simeq(\mc O/\eta)^d$, of length $d\,v(\eta)$.
\end{proof}

\noindent\emph{Reading guidance.} Steps (1) and (6) are the ones to master. Step (1) shows $\wp\cap I=0$ from torsion-freeness of $\T$, not from reducedness, which fails for
the full Hecke algebra. Step (6) identifies the cokernel $Q$ with the discriminant module of the pairing on $L[\wp]$, whose length is $v(\Delta)$; the inequality is then
the trivial bound on a module generated by $d$ elements and killed by $\eta$. The finite residue field is used only to count orders; in length form the lemma holds over any
discrete valuation ring (\Cref{ch:congruence}).

\begin{lemma}[the paper's Lemma lem:hecke]\label{lem:paper:hecke}
Let $E$ be $X_0(N)$-optimal and $\ell$ odd with $\ell\nmid2N$ and $\bar\rho$ irreducible. Then $v_\ell(\degphi)=v_\ell(r_E)=v_\lambda(\eta^\emptyset_{f,\lambda})$.
\end{lemma}

\begin{proof}[Proof]
$v_\ell(\degphi)=v_\ell(r_E)$ is \cite[Theorem 2.2]{ARS12} ($\ell^2\nmid N$), and $v_\ell(\degphi)=v_\lambda(\eta^\emptyset_{f,\lambda})$ is \Cref{prop:paper:pullback}.
\end{proof}

\noindent The paragraph after the lemma describes a second route, Lemma lem:ddt417 applied to $L=H^1(X_0(N),\Zl)_{\mf m}$, which without the projection formula gives
only $v_\ell(r_E)\ge v_\lambda(\eta^\emptyset_{f,\lambda})$, with equality when $L$ is free over $\T_{\mf m}$; it is the argument of Step 2 of Theorem thm:quant, for a depleted form not
attached to an optimal quotient. Two facts make it work: the other local factors of $\T\otimes\Zl$ lie in the kernel $\wp$, so the congruence ideal is computed in $\T_{\mf m}$; and
$e_{\mf m}$ is self-adjoint, so the pairing is perfect on each local component. The paper's example: at \curve{99a1} and $\ell=3$ the Hecke ideal has valuation $v_3(12)=1$ and the cohomological ideal $v_3(4)=0$, so the inequality of Lemma lem:ddt417 is
strict ($2>0$) and $H^1(X_0(99),\Z_3)_{\mf m}$ is not free (\Cref{ch:moddeg}).

\begin{remark}[the paper's Remark rem:free]\label{rem:paper:free}
For $\ell\nmid2N$ and $\bar\rho$ irreducible, $J_0(N)[\mf m]$ is two-dimensional (\cite[Theorem 3.5]{RibetStein} for $J_1(N)$, as cited; the exceptional Serre weight $\ell$ does not occur
in weight $2$ with $\ell\nmid N$; the kernel of the pullback $J_0(N)\to J_1(N)$ is the Shimura subgroup, which is Eisenstein; the references \cite{LingOesterle,Ribet88} for
it are the book's); then $\T_{\mf m}$ is Gorenstein and the Tate module
is free of rank two \cite{Tilouine97}. The identification of $T_\ell J_0(N)$ with $H^1(X_0(N),\Zl)$ uses Poincar\'e duality composed with $w_N$, not Galois equivariant, but freeness
transfers. After \Cref{prop:paper:pullback} this is not used by Theorems thm:main or thm:quant; it gives a second proof of $v_\ell(r_E)=v_\ell(\degphi)$ in this range and explains
why the Hecke and the cohomological congruence ideals agree there (\Cref{ch:moddeg}).
\end{remark}

\section[Paper 4.2, 4.3: Theorem thm:main]{Sections 4.2 and 4.3 of the paper: the adjoint Selmer group and the proof of Theorem thm:main}\label{sec:paper:selmer}

This section reads the set-up of the theorem of Diamond, Flach and Guo, Proposition prop:Sf, the three imported results, Lemmas lem:dfg21, lem:fpr and
lem:h0, the proof of Theorem thm:main, Corollary cor:explicit, Remark rem:degree and Remark rem:ko.

\begin{definition}[the Selmer set-up; Section 4.2]\label{def:paper:selmer-setup}
Let $f=f_E$, $\Sigma$ the set of primes dividing $N$, $\lambda=\ell\nmid2N$, $A=\ad T$ (by \Cref{prop:paper:pullback}, $\ad M_{f,\lambda}=\ad T$), $W$, $T^D\simeq A(1)$, $V^D$, $W^D$ as in
\Cref{conv:paper:duals}. $H^1_\Sigma(\Q,W)\supseteq H^1_f(\Q,W)$ has no condition at $p\in\Sigma$, the Bloch-Kato condition at $\ell$ and the unramified condition elsewhere \cite[Section 2.1]{DFG04}.
$S_f$ is the set of $\lambda$ with $\lambda\mid2N$ or $\bar\rho|_{\Gal(\Qbar/\Q(\sqrt{\ell^*}))}$ not absolutely irreducible \cite[(31)]{DFG04} (verified; it is the weight-$2$ case of $\lambda\mid Nk!$).
\end{definition}

\begin{proposition}[the paper's Proposition prop:Sf]\label{prop:paper:Sf}
Let $\ell$ be odd, $\ell\nmid N$, $E[\ell]$ irreducible. If $\ell\ge5$ then $\ell\notin S_f$. If $\ell=3$ then $3\in S_f$ iff the image of $\bar\rho_{E,3}$ is conjugate to $N_s(\F_3)$, and then $E$ is
ordinary at $3$. In all cases $\ell\notin S_f$ iff $h_\ell(E)=0$.
\end{proposition}

\noindent\emph{Where proved.} \Cref{ch:galrep} (with the description of $\bar\rho|_{I_\ell}$ for good reduction, \cite[Section 1]{Serre72}), Clifford theory in \Cref{ch:reps}.

\begin{proof}[Expanded proof]
Let $K=\Q(\sqrt{\ell^*})$. $\bar\rho$ is absolutely irreducible (complex conjugation has eigenvalues $1\ne-1$). For an absolutely irreducible $\bar\rho$ and an index-two subgroup $G_K$,
$\bar\rho|_{G_K}$ fails to be absolutely irreducible iff $\bar\rho\simeq\bar\rho\otimes\chi_{\ell^*}$ (Clifford theory, \Cref{ch:reps}; this needs $\bar\rho$ irreducible). Suppose so and restrict to $I_\ell$,
where $\chi_{\ell^*}=\omega^{(\ell-1)/2}$, $\omega$ the fundamental character of level one. Good reduction at $\ell$ gives: if $E$ is ordinary, $\bar\rho|_{I_\ell}^{\mathrm{ss}}=\omega\oplus1$; then
$\{\omega,1\}=\{\omega^{(\ell+1)/2},\omega^{(\ell-1)/2}\}$, and $\omega^{(\ell-1)/2}=-1\ne1$ forces $\omega^{(\ell+1)/2}=1$, i.e. $(\ell-1)\mid(\ell+1)/2$, so $\ell=3$. If $E$ is supersingular, $\bar\rho|_{I_\ell}\otimes\F_{\ell^2}=\omega_2\oplus\omega_2^\ell$
with $\omega_2$ of order $\ell^2-1$ and $\omega_2^{\ell+1}=\omega$, so $\omega^{(\ell-1)/2}=\omega_2^{(\ell^2-1)/2}=-1$; then $\{\omega_2,\omega_2^\ell\}=\{-\omega_2,-\omega_2^\ell\}$ needs $\omega_2^\ell=-\omega_2$, i.e.
$\omega_2^{\ell-1}=\omega_2^{(\ell^2-1)/2}$, i.e. $(\ell^2-1)\mid(\ell-1)^2/2$, i.e. $(\ell+1)\mid(\ell-1)/2$: impossible. So $\ell\ge5$ never occurs, and at $\ell=3$ the curve is ordinary; the image is
then $N_s(\F_3)$ by \Cref{prop:paper:ns}. The last statement is \Cref{lem:paper:h0}.
\end{proof}

\noindent\emph{Reading guidance.} The whole proof is a comparison of characters of inertia at $\ell$: twisting by $\chi_{\ell^*}=\omega^{(\ell-1)/2}$ must permute the characters of
$\bar\rho|_{I_\ell}$, and the fundamental characters leave no room for this except in the ordinary case at $\ell=3$. So for $\ell\nmid N$ the residual representations induced from
$\Q(\sqrt{\ell^*})$ with $\ell\ge5$ do not occur; Section 6.6 of the paper notes that such curves all have $\ell\mid N$.

\begin{theorem}[Diamond, Flach and Guo, Theorem 3.7; the paper's Theorem thm:dfg37]\label{thm:paper:dfg37}
If $\ell\notin S_f$, then $H^1_\Sigma(\Q,W)$ is finite of length $v_\lambda(\eta^\Sigma_{f,\lambda})$.
\end{theorem}

\begin{theorem}[Diamond, Flach and Guo, Theorem 2.7; the paper's Theorem thm:dfg27]\label{thm:paper:dfg27}
If $\ell\notin S_f$, then $H^0(\Q,V^D)=H^1_f(\Q,V^D)=0$, and the same for $A\otimes\Ql$.
\end{theorem}

\noindent\emph{Where.} Both are stated with references in \Cref{ch:congruence} (Taylor-Wiles method in the axiomatic form of \cite{DFG04}, sketched there); we do not prove
them. Theorem 3.7 is stated in \cite{DFG04} for $\Sigma\not\ni\ell$ with $M_{f,\lambda}$ minimally ramified outside $\Sigma$, which holds for the primes dividing $N$; this is their Theorem 0.2. It
covers $\ell=3$ in weight $2$: $S_f$ excludes only $\lambda\mid2N$ and the induced case, and the proof treats $\ell=3$ explicitly (both checked against the published text, \Cref{ch:congruence}).
For Theorem 0.1 at $\ell=3$ with $a_3\not\equiv0\pmod3$ Diamond, Flach and Guo also use Perrin-Riou's and Colmez's results \cite[Proposition 2.16]{DFG04}; Theorems 2.7 and 3.7 do not need them.
The imported theorems are the deep input of the paper, and the only place where the Taylor-Wiles method enters.

\begin{lemma}[Diamond, Flach and Guo, Lemma 2.1; the paper's Lemma lem:dfg21]\label{lem:paper:dfg21}
Suppose $H^0(\Q,V^D)=H^1_f(\Q,V^D)=0$, and let $\Sigma$ be a nonempty finite set of primes, $\ell\notin\Sigma$, containing every prime $p\ne\ell$ at which $W$ is ramified. With
$H^1_f(\Q_p,T^D)=\iota^{-1}H^1_f(\Q_p,V^D)$ for $\iota\colon H^1(\Q_p,T^D)\to H^1(\Q_p,V^D)$, there is an exact sequence
\[
0\to H^1_f(\Q,W)\to H^1_\Sigma(\Q,W)\to\bigoplus_{p\in\Sigma\cup\{\infty\}}H^1_f(\Q_p,T^D)^*\to H^0(\Q,W^D)^*\to0.
\]
\end{lemma}

\noindent\emph{Where.} Stated with the reference \cite[Lemma 2.1]{DFG04} (verified), and derived in \Cref{ch:galcoh,ch:blochkato} from Poitou-Tate duality (the comparison of
Selmer groups for two sets of local conditions). The paper's condition on $\Sigma$ implies the one of \cite{DFG04} ($\Sigma$ contains every prime where $H^1_f(\Q_p,W)\ne H^1_{\ur}(\Q_p,W)$), as
the paper says.

\emph{The paragraph after the lemma, expanded.} For $p\in\Sigma$, $p\ne\ell$: $H^1_f(\Q_p,V^D)=H^1_{\ur}(\Q_p,V^D)=(V^D)^{I_p}/(\Frob_p-1)$, and $\Frob_p-1$ is injective on $(V^D)^{I_p}$ because
$\det(1-\Frob_p^{-1}\mid(V^D)^{I_p})=L_p(A_f(1),0)^{-1}=L_p(A_f,1)^{-1}\ne0$ (geometric Frobenius; with the arithmetic one the determinant differs by a unit at the primes in question, and
injectivity is the same statement for $\Frob$ and $\Frob^{-1}$); so the quotient is $0$, and $H^0(\Q_p,V^D)=\ker(\Frob_p-1)=0$, which makes the hypothesis $H^0(\Q,V^D)=0$ automatic.
Hence $H^1_f(\Q_p,T^D)$ is the torsion of $H^1(\Q_p,T^D)$, a finite group (the word ``torsion'' matters: $H^1(\Q_p,T^D)$ itself can be infinite). At $p=\ell$ the computation fails
($\dim H^1_f(\Q_\ell,V^D)=2$), which is why $\ell\notin\Sigma$ is part of the lemma, as the paper says; it is harmless, as $\ell\nmid N$. At $\infty$, $H^1(\R,T^D)$
is killed by $2$, so $H^1_f(\R,T^D)=0$ for $\ell$ odd.

\begin{lemma}[the paper's Lemma lem:fpr]\label{lem:paper:fpr}
For $p\in\Sigma$: $\Fitt_{\Zl}H^1_f(\Q_p,T^D)=L_p(A_f,1)^{-1}\Tam_p(T^D)$ and $\Tam_p(T^D)=\Tam_p(A)$; and $L_p(A_f,1)^{-1}=L_p(A_f(1),0)^{-1}$ is a nonzero $\ell$-integral rational number, of
valuation $e_\ell(E,p)$ at an additive prime and $v_\ell(1-p^{-2})$ at a multiplicative one.
\end{lemma}

\noindent\emph{Where proved.} The first identity is \cite[I.4.2.2]{FPR94} as recalled in the proof of \cite[Theorem 2.15]{DFG04}; it is proved in \Cref{ch:blochkato} (the
$L$-part in \Cref{ch:lfunctions}). The second is \cite[(57)]{DFG04} in general (verified), and for the elliptic adjoint it follows from the explicit computations: both sides are
trivial at potentially good primes (\Cref{lem:paper:tamadd} and the same argument for $A(1)$), and at a potentially multiplicative prime the count of \Cref{lem:paper:tamfactor}, with
the Frobenius action twisted by $\chi_\ell(\Frob)=q$, gives the same number (\Cref{ch:blochkato}; \Cref{ex:paper:tamfactor-coordinates}).

\emph{Expansion of the first identity.} Inflation-restriction with $G/I\simeq\widehat\Z$ gives $0\to X^I/(\Frob-1)X^I\to H^1(G,X)\to H^1(I,X)^{G/I}\to0$ for $X=T^D$ (\Cref{ch:cohomology}). $\Frob-1$ is
injective on $X^I$ (previous paragraph), and $X^I/(\Frob-1)X^I$ is finite of length $v_\ell(\det(\Frob-1\mid X^I))=v_\ell(L_p(A_f,1)^{-1})$. Taking torsion: the torsion of $H^1(G,X)$ is an
extension of $H^1(I,X)^{G/I}_{\mathrm{tor}}=(H^1(I,X)_{\mathrm{tor}})^{G/I}$ by the finite group $X^I/(\Frob-1)X^I$; lengths add, and $\Fitt=\ell^{\len}$ for finite $\Zl$-modules. Integrality: $P_p$ has
integral coefficients and $p^{-2}$ is an $\ell$-adic unit.

\begin{lemma}[the paper's Lemma lem:h0]\label{lem:paper:h0}
Let $\ell$ be odd with $\bar\rho$ irreducible. Then $H^0(\Q,W^D)=H^0(\Q,\ad T(1)\otimes\Qlmod)$ is finite cyclic of length $h_\ell(E)$, the largest $k$ with $\rho\bmod\ell^k\simeq(\rho\otimes\chi_\ell)\bmod\ell^k$.
Moreover $h_\ell(E)=0$ for $\ell\ge5$, and $h_3(E)\le1$, with equality iff $\bar\rho_{E,3}\simeq\bar\rho_{E,3}\otimes\chi_{-3}$, iff $\bar\rho_{E,3}$ is not absolutely irreducible on $\Gal(\Qbar/\Q(\sqrt{-3}))$, iff the image
of $\bar\rho_{E,3}$ is conjugate to $N_s(\F_3)$. This holds with or without complex multiplication.
\end{lemma}

\noindent\emph{Where proved.} \Cref{ch:galrep}; Schur's lemma, Clifford theory and Frobenius reciprocity in \Cref{ch:reps}; $\chi_\ell(G_\Q)=\Zl^\times$ in \Cref{ch:ant}.

\begin{proof}[Expanded proof]
(1) $\bar\rho$ is absolutely irreducible. (2) $H^0(\Q,W^D)[\ell^k]=H^0(\Q,\ad(T/\ell^k)(1))$, the trace-zero $G_\Q$-maps $T/\ell^k\to(T/\ell^k)(1)$; every $G_\Q$-map $\phi$ has trace $0$, since
$\tr\phi=\chi_\ell(g)\tr\phi$ for all $g$ (equivariance says $\phi\rho(g)=\chi_\ell(g)\rho(g)\phi$, so $\phi=\chi_\ell(g)\rho(g)\phi\rho(g)^{-1}$, and the trace is invariant under conjugation) and
$1-\chi_\ell(c)=2$ is a unit for complex conjugation $c$. (3) For $k=1$ these are the maps $\bar\rho\to\bar\rho\otimes\bar\chi_\ell$, at most a line by Schur.
So the $\ell$-torsion of the submodule $H^0(\Q,W^D)$ of $(\Qlmod)^3$ has rank at most one, and $H^0(\Q,W^D)$ is finite cyclic or isomorphic to $\Qlmod$. (4) An element of exact
order $\ell^k$ is a map $\phi\colon T/\ell^k\to(T/\ell^k)(1)$ nonzero mod $\ell$, hence an isomorphism mod $\ell$ (Schur), hence an isomorphism (Nakayama); conversely. (5) Determinants:
taking determinants in $\phi\rho(g)\phi^{-1}=\chi_\ell(g)\rho(g)$ modulo $\ell^k$ gives $\det\rho(g)=\chi_\ell(g)^2\det\rho(g)$, and $\det\rho(g)=\chi_\ell(g)$ is a unit, so
$\chi_\ell(g)^2\equiv1\pmod{\ell^k}$ for all $g$. Since $\chi_\ell(G_\Q)=\Zl^\times\ni2$ and $4\not\equiv1$ modulo $\ell$ for $\ell\ge5$ and modulo $9$, $k=0$ for $\ell\ge5$ and $k\le1$ for $\ell=3$; in
particular the $\Qlmod$ alternative of (3) is excluded and the group is finite. (6) At $\ell=3$, $\bar\chi_3^{\mathrm{cyc}}=\chi_{-3}$. If $\phi\colon\bar\rho\to\bar\rho\otimes\chi_{-3}$ is an isomorphism, it
commutes with $\bar\rho(G_K)$, $K=\Q(\sqrt{-3})$, and is not scalar, so $\bar\rho|_{G_K}$ is not absolutely irreducible; conversely a $G_K$-stable line over $\Fbar_3$ with character $\psi$ gives, by
Frobenius reciprocity, a nonzero map $\Ind\psi\to\bar\rho\otimes\Fbar_3$, an isomorphism by absolute irreducibility, and $\Ind\psi\otimes\chi_{-3}\simeq\Ind\psi$; the isomorphism descends to $\F_3$,
Hom spaces commuting with extension of scalars. (7) The last equivalence is \Cref{prop:paper:ns}.
\end{proof}

\noindent\emph{Reading guidance.} Two steps carry the proof. Step (5), the determinant argument, is what makes $H^0(\Q,W^D)$ finite and bounds $k$: before it, the module
could a priori be $\Qlmod$, which is why the paper says ``finite cyclic or isomorphic to $\Q_\ell/\Z_\ell$'' first and excludes the second case afterwards. Step (6) is Clifford
theory for the index-two subgroup $G_{\Q(\sqrt{-3})}$. The last equivalence, with the image $N_s(\F_3)$, is the first statement of Proposition prop:ns, whose proof uses this lemma
only for the value $h_3(E)=1$; there is no circularity. The clause ``in particular $H^0(\Q,W^D)=0$ if $\ell\notin S_f$'' combines the lemma with \Cref{prop:paper:Sf}.

\paragraph{The proof of Theorem thm:main (Section 4.3).} The proof uses, besides the lemmas above, \Cref{thm:paper:dfg14,thm:paper:dfg37,thm:paper:dfg27,lem:paper:dfg21},
\cite[Theorem 2.2]{ARS12} and the comparison of integral structures (\Cref{prop:paper:integral}). Each step is listed with the hypothesis it uses.

\begin{proof}[Proof of \Cref{thm:paper:main}, expanded]
Let $\Sigma=\{q\mid N\}$; then $\ell\notin\Sigma$.
\emph{Step 1} (optimality, $\ell\nmid2N$, $\bar\rho$ irreducible): $v_\ell(\degphi)=v_\lambda(\eta^\emptyset_{f,\lambda})$, \Cref{lem:paper:hecke}.
\emph{Step 2} ($\bar\rho$ irreducible, $\ell\notin\Sigma$): $v_\lambda(\eta^\emptyset_{f,\lambda})=v_\lambda(\eta^\Sigma_{f,\lambda})-\sum_{p\in\Sigma}v_\ell(L^{\nv}_p(A_f,1)^{-1})$, \Cref{thm:paper:dfg14}.
\emph{Step 3} ($\ell\notin S_f$): $v_\lambda(\eta^\Sigma_{f,\lambda})=\len H^1_\Sigma(\Q,W)$, \Cref{thm:paper:dfg37}.
\emph{Step 4} ($\ell\notin S_f$, $\ell$ odd): by \Cref{thm:paper:dfg27} only $H^1_f(\Q,V^D)=0$ is needed in \Cref{lem:paper:dfg21}, $H^0(\Q,V^D)=0$ being automatic; $H^0(\Q,W^D)=0$ by
\Cref{lem:paper:h0} (as $h_\ell(E)=0$ for $\ell\notin S_f$, \Cref{prop:paper:Sf}); the term at $\infty$ is $0$. Lengths add along the exact sequence (Pontryagin duality preserves length):
$\len H^1_\Sigma(\Q,W)=\len H^1_f(\Q,W)+\sum_{p\in\Sigma}\len H^1_f(\Q_p,T^D)$.
\emph{Step 5} ($p\ne\ell$): $\len H^1_f(\Q_p,T^D)=v_\ell(L_p(A_f,1)^{-1})+v_\ell(\Tam_p(A))$, \Cref{lem:paper:fpr}.
\emph{Step 6} (assembly): combining, $v_\ell(\degphi)=\len H^1_f(\Q,W)+\sum_pv_\ell(\Tam_p(\ad T))+\sum_pv_\ell\bigl(L_p(A_f,1)^{-1}/L^{\nv}_p(A_f,1)^{-1}\bigr)$, and the last sum is $\sum_{p\mid N}e_\ell(E,p)$ by
\Cref{def:paper:terms}: the naive factors at multiplicative primes cancel against the true ones, and at additive primes the naive factor is $1$. This is (eq:identity).
\emph{Step 7} (eq:thm): $\len H^1_f\ge0$; $v_\ell(\Tam_p(\ad T))\ge t_\ell(E,p)$ at every potentially multiplicative $p$, multiplicative or additive, $p=2$ included (\Cref{lem:paper:tamfactor}, step
(9)); $t_\ell(E,p)=0$ at the other primes.
\end{proof}

\noindent\emph{Reading guidance.} The bookkeeping in Step 6 is the point of the paper: the naive factors that Theorem thm:dfg14 removes and the true factors that Lemma lem:fpr puts
back agree at the multiplicative primes and differ exactly at the additive ones, and the difference is $\sum e_\ell$. Step 7 is local at each prime, so it covers $p=2$, where
$t_\ell(E,2)$ can be positive although no quadratic twist of conductor dividing $N$ need make the reduction multiplicative.

\begin{corollary}[the paper's Corollary cor:explicit]\label{cor:paper:explicit}
Under the hypothesis of \Cref{thm:paper:main},
\begin{multline*}
v_\ell(\degphi)=\len H^1_f(\Q,\ad T\otimes\Qlmod)+\sum_{q\mid N}e_\ell(E,q)\\
+\sum_{q:\,v_q(j_E)<0}\Bigl(t_\ell(E,q)+\min\bigl(t_\ell(E,q),v_\ell(q-1),v_\ell(c(u_q))\bigr)\Bigr).
\end{multline*}
\end{corollary}

\begin{proof}
(eq:identity) with $\Tam_q(\ad T)$ trivial at good primes (unramified $\ad T$: $H^1(I_q,\ad T)=\Hom(I_q,\ad T)$ is torsion-free) and at potentially good additive primes
(\Cref{lem:paper:tamadd}), and given by \Cref{lem:paper:tamfactor} at the potentially multiplicative ones, where $v_\ell(n)=t_\ell(E,q)$.
\end{proof}

\noindent So the excess of $v_\ell(\degphi)$ over an explicit local bound is exactly the length of the adjoint Selmer group. The paper's Table tab:excess gives its
distribution on the optimal non-CM curves below conductor $10{,}000$ with $\ell\nmid N$, $E[\ell]$ irreducible and, at $\ell=3$, image not $\Ns$, for $\ell=3,5,7$; at $\ell=3$ the
excess is $0$, $1$, $2$, $3$, at least $4$ for $8{,}688$, $3{,}288$, $942$, $248$, $37$ of $13{,}203$ curves, and it is never negative, as the corollary requires. The book reproduced
the row $\ell=3$ (\Cref{ch:blochkato}). Examples, with the
Tate data recomputed: \curve{121a1} ($e_3(E,11)=1$, $\degphi=6$): Selmer length $0$; \curve{791a1} (explicit bound $1+1$ at $7$, $\degphi=144$): $0$; \curve{1184f1} (explicit bound $2$,
$\degphi=1728$): $1$; \curve{1748e1} (explicit bound $3$, $\degphi=3240$): $1$.

\begin{proposition}[the degree formula; the paper's Remark rem:degree]\label{prop:paper:degree}
For $E$ optimal, $\degphi\cdot\vol(\Lambda_E)=4\pi^2c_E^2\langle f,f\rangle$ (unnormalised Petersson product) and $\langle f,f\rangle=N\,L^{\nv}(\Sym f,2)/(8\pi^3)$ give
\[
\degphi=\frac{c_E^2\,N\,L^{\nv}(\ad f,1)}{2\pi\vol(\Lambda_E)}=B_E\prod_{q^2\mid N}L_q(A_f,1)^{-1},\qquad B_E=\frac{c_E^2\,N\,L(\ad f,1)}{2\pi\vol(\Lambda_E)}\in\Q^\times;
\]
so $v_\ell(\degphi)=v_\ell(B_E)+\sum_{q\mid N}e_\ell(E,q)$ for every $\ell\nmid N$, and (eq:thm) is equivalent to $v_\ell(B_E)\ge\sum_qt_\ell(E,q)$.
\end{proposition}

\noindent\emph{Where proved.} The formula of Zagier and Watkins (with $c_E^2$ in the numerator) in \Cref{ch:moddeg}; the Petersson norm at every level in \Cref{ch:hecke,ch:lfunctions};
\cite[Lemma 2.12]{DFG04} uses the naive value (verified).

\begin{proof}
Divide the first formula by $\vol(\Lambda_E)$ and substitute the second: $\degphi=4\pi^2c_E^2N\,L^{\nv}(\Sym f,2)/(8\pi^3\vol)$, and $L^{\nv}(\Sym f,2)=L^{\nv}(\ad f,1)$. The naive and true $L$-functions
differ exactly at the additive primes ($q^2\mid N$), where the naive factor is $1$: $L^{\nv}(\ad f,1)=L(\ad f,1)\prod_{q^2\mid N}L_q(A_f,1)^{-1}$. $B_E$ is rational because $\degphi$ and the $L_q(A_f,1)^{-1}$
are; taking $v_\ell$ with $\ell\nmid N$ (so each $L_q(A_f,1)^{-1}$ has valuation $e_\ell(E,q)$, and $e_\ell=0$ at multiplicative $q$) gives the rest.
\end{proof}

\noindent The paper checks the first formula to $30$ digits on five curves (with $c_E=1$). The formula isolates the role of the omitted
factors in one line: (eq:thm) is a statement about the true adjoint $L$-value.

\begin{remark}[the paper's Remark rem:ko]\label{rem:paper:ko}
(1) At a potentially multiplicative $q\equiv1\pmod\ell$ with $\ell\mid n$ and $\bar u$ an $\ell$-th power, \Cref{lem:paper:tamfactor} gives $v_\ell(\Tam_q)\ge t_\ell+1$, and (eq:identity) raises the bound by
exactly $\min(v_\ell(n),v_\ell(q-1),v_\ell(c(u)))$ when $\ell\notin S_f$ (\Cref{cor:paper:explicit}). (2) The formula of Kim and Ota \cite[Theorem 1.1, Definition 2.12]{KimOta} (as cited; their
hypotheses as the paper lists them) gives $v_\ell(r_E)\ge\sum_{q\mid N^-}t_\ell(E,q)$, contained in (eq:thm) under its hypothesis; the equality is theirs. Whether their $\eta_f(N)$ is the Hecke or
the cohomological congruence ideal does not matter at $\ell\nmid N$ with $E[\ell]$ irreducible, since the two agree there (\Cref{lem:paper:hecke}, \Cref{ch:congruence}).
\end{remark}

\section[Paper 4.4: Theorem thm:quant]{Section 4.4 of the paper: twist monotonicity and Theorem thm:quant}\label{sec:paper:twisting}

This section reads the paper's Section 4.4: Theorem thm:twist (comparison of congruence numbers along a twist), the data paragraph after it, the set-up
and Lemma lem:step1 (the twist transport), the proof of Theorem thm:quant in four steps with its equality clause, and Remark rem:afterquant. Theorem thm:quant is
the only result of the paper on the $\kod{I}_0^*$ and $\kod{I}_n^*$ terms that does not need the hypothesis over $\Q(\sqrt{\ell^*})$.

\begin{theorem}[the paper's Theorem thm:twist]\label{thm:paper:twist}
Let $E/\Q$ have conductor $N$, $S$ a set of odd primes, $D=\prod_{q\in S}q$, $E_S=E\otimes\chi_S$ of conductor $N_S$. If $q^2\mid N$ for every $q\in S$ and $\lcm(N_S,D^2)\mid N$, then $r_{E_S}\mid r_E$;
if moreover $\lcm(N,D^2)\mid N_S$, then $r_{E_S}=r_E$.
\end{theorem}

\noindent\emph{Where proved.} \Cref{ch:congruence}; the twist map, strong multiplicity one and the $\T$-stable complement in \Cref{ch:hecke}. The twist map is written
$\iota_S$ in the book, since $\iota$ is also the map $H^1(\Q_p,T^D)\to H^1(\Q_p,V^D)$.

\begin{proof}[Expanded proof]
Let $f=f_E$, $g=f_{E_S}$, $\iota_S(h)=h\otimes\chi_S=\sum\chi_S(n)a_n(h)q^n$. (1) $\iota_S$ maps $S_2(\Gamma_0(N_S))$ into $S_2(\Gamma_0(\lcm(N_S,D^2)))\subseteq S_2(\Gamma_0(N))$ (Shimura's
level of a twist, \Cref{ch:hecke}), preserves integrality, and intertwines $T_p$ with $\chi_S(p)T_p$ for $p\nmid DN_S$. (2) $\iota_S(g)=f$: both sides are multiplicative in $n$;
$a_p(f)=\chi_S(p)a_p(g)$ for $p\nmid D$, the twist being unramified there; and $a_{q^k}(f)=0$ for $q\in S$ as $q^2\mid N$. (3) Let $g^\perp$ be the $\T$-stable complement of the $g$-line at
level $N_S$, the sum of the generalised eigenspaces of the other systems of eigenvalues (at level $N_S$, $g$ is new, and this is also the Petersson complement). A
generalised eigenform $h$ of level dividing $N_S$ whose $T_p$-eigenvalues ($p\nmid N_S$) differ from those of $g$ at infinitely many $p$ is sent to one whose eigenvalues differ from
those of $f$ at infinitely many $p$, hence (strong multiplicity one) into $f^\perp$. (4) A congruence $g\equiv h\pmod r$ with $h\in g^\perp$ integral gives $f\equiv\iota_S(h)\pmod r$ with $\iota_S(h)\in f^\perp$
integral, so $r\mid r_E$ (\Cref{lem:paper:crit}); with $r=r_{E_S}$ this is the first claim. (5) Under the second hypothesis $\iota_S$ also maps level $N$ to level $N_S$ and the roles
exchange.
\end{proof}

\noindent The hypotheses hold when $E$ has type $\kod{I}_n^*$ at each $q\in S$ ($N_S=N/D$; case (a)), type $\kod{I}_0^*$ ($N_S=N/D^2$; case (b)), or a potentially good type
$\kod{II},\kod{III},\kod{IV}$ or a dual at each $q\in S$, $q\ge5$ ($N_S=N$, equality; case (c)). The step to watch is (3): the complement of the $g$-line must be the $\T$-stable one, the
sum of the other generalised eigenspaces, and the twist map respects it because it moves systems of Hecke eigenvalues by $\chi_S$.

\paragraph{The data paragraph after the theorem.} The paper reports $157{,}938$ and $292{,}580$ working-set pairs with $v_\ell(\degphi)\ge v_\ell(\deg\varphi_{E^{(q)}_0})$, $E^{(q)}_0$ the optimal
curve of the class of $E^{(q)}$, with no exception, equality in every pair of case (c), and a difference equal to the Euler term in cases (a) and (b), with the examples
\curve{539c1}/\curve{77c1}, \curve{980i1}/\curve{140b1}, \curve{121d1}/\curve{11a1} ($v_3$: $2$ against $1$, $2$ against $1$, $1$ against $0$) and \curve{539d1}/\curve{11a1} ($2$ against $0$, type
$\kod{I}_0^*$, $e_3(E,7)=2$). The examples were recomputed for the book (Section~\ref{sec:paper:verifications}); below $10{,}000$ the difference equals the Euler term in all $25{,}477$ pairs
(curve, $\ell$) with $S\ne\emptyset$, $\ell\nmid N$ and $E[\ell]$ irreducible, as the paper says (\Cref{ch:congruence,ch:computation}). In case (a) at $\ell=3$ the difference is always positive,
since $3\mid q^2-1$ for every $q\ne3$.

\paragraph{The set-up.} For any $E$ of conductor $N$ let $S$ be the set of odd primes at which $E$ has type $\kod{I}_0^*$ or $\kod{I}_n^*$, $D=\prod_{q\in S}q$, $E_S=E\otimes\chi_S$ of
conductor $N_S$, $g=f_{E_S}$ and $\iota_S$ the twist map (the book writes $\iota_S$, since $\iota$ is also the map $H^1(\Q_p,T^D)\to H^1(\Q_p,V^D)$). By \Cref{lem:paper:twist}, valid at
every odd $q$ for these exchanges, $E_S$ is good at the $\kod{I}_0^*$ primes of $S$ and multiplicative at the $\kod{I}_n^*$ primes; so $\delta_q(g)=2$, resp. $1$, and
$N_S\prod_{q\in S}q^{\delta_q}=N$: the depleted form $g^S$ lives at level $N$, and $\iota_S(g^S)=f$, $\iota_S(f)=g^S$ ($\iota_S$ kills the coefficients at multiples of $D$). Let $\pi\colon\T\to\Z$ be
the eigencharacter of $g^S$ on the full Hecke algebra of level $N$ and $\eta^{\mathrm{Hecke}}(g^S)=\pi(\Ann_\T\ker\pi)$. Its generalised eigenspace is the line of $g^S$: the forms
with the $T_p$-eigenvalues of $g$ ($p\nmid N$) span the $g(dz)$, $d\mid N/N_S$, on which $U_q$, $q\in S$, has the eigenvalue $0$ with multiplicity one. So the proof of \Cref{lem:paper:crit}, with
$f^\perp$ replaced by the $\T$-stable complement $C$ of the $g^S$-line (not the Petersson complement, which is not $\T$-stable when the $U_p$ are involved; \Cref{ch:hecke}), shows that
$\ell^{v_\ell(\eta^{\mathrm{Hecke}}(g^S))}$ is the largest $\ell^t$ with $g^S\equiv h\pmod{\ell^t}$ for some $h\in C\cap S_2(\Gamma_0(N);\Zl)$.

\begin{lemma}[the paper's Lemma lem:step1]\label{lem:paper:step1}
In this situation, $v_\ell(r_E)\le v_\ell(\eta^{\mathrm{Hecke}}(g^S))$ for every prime $\ell$, and $v_\ell(r_E)=v_\ell(\eta^{\mathrm{Hecke}}(g^S))$ for every prime $\ell\notin S$. No other hypothesis
on $\ell$ is needed: $\ell=2$, $\ell\mid N$, $E[\ell]$ reducible and $\ell\in S_f$ are allowed, and $E$ need not be optimal.
\end{lemma}

\noindent\emph{Where proved.} \Cref{ch:congruence}; the twist map, strong multiplicity one and the $\T$-stable complement in \Cref{ch:hecke}.

\begin{proof}[Expanded proof]
\emph{The inequality $\le$.} By \Cref{lem:paper:crit} there is an integral $h\in f^\perp$ with $f\equiv h\pmod{r_E}$. As $f$ is new, $f^\perp$ is the sum of the generalised eigenspaces of
the $T_p$, $p\nmid N$, for systems differing from that of $f$ at infinitely many $p$. The map $\iota_S$ preserves integrality, maps $S_2(\Gamma_0(N))$ into itself (the level of a twist by
$\chi_S$ divides $\lcm(N,D^2)=N$), and intertwines $T_p$ with $\chi_S(p)T_p$ for $p\nmid N$; so $\iota_S(h)\in C$ is integral, and $g^S=\iota_S(f)\equiv\iota_S(h)\pmod{r_E}$.

\emph{The inequality $\ge$ for $\ell\notin S$.} Let $t=v_\ell(\eta^{\mathrm{Hecke}}(g^S))$, $\mf m$ the maximal ideal of $\T_\ell=\T\otimes\Zl$ containing $\ker\pi$ and $\ell$; it contains $U_q$ for
$q\in S$. Take $h\in C$ integral with $g^S\equiv h\pmod{\ell^t}$; the idempotent $e_{\mf m}$ fixes $g^S$ and preserves $C$ and integrality, so $h'=e_{\mf m}h$ does as well. The form $h'$ lies in
the generalised eigenspaces of eigencharacters $\phi\ne\pi$ congruent to $\pi$, and such a $\phi$ has $\phi(U_q)\equiv0$. Every eigenvalue of $U_q$ on $S_2(\Gamma_0(N))$ is $0$, $\pm1$
or a root of $X^2-a_q(h_0)X+q$ for a newform $h_0$, hence $0$ or a unit at every prime not above $q$; as $\ell\ne q$, $\phi(U_q)=0$. This is the only place where $\ell\notin S$ is
used. Each such $\phi$ comes from a newform $h_0\ne g$ (if $h_0=g$, $\phi$ would agree with $\pi$ everywhere), so $\iota_S$ maps its generalised eigenspace into that of the system
$\chi_S(p)a_p(h_0)$, which differs from that of $f$ at infinitely many $p$; so $\iota_S(h')\in f^\perp$, and $f=\iota_S(g^S)\equiv\iota_S(h')\pmod{\ell^t}$ gives $t\le v_\ell(r_E)$.
\end{proof}

\noindent\emph{Reading guidance.} The lemma is elementary: no Galois representation enters, only the Hecke action and strong multiplicity one. The step that needs
$\ell\notin S$ is the vanishing $\phi(U_q)=0$ from $\phi(U_q)\equiv0$. The paper's sentence after the lemma says what happens at $\ell\in S$: at level $539$ the congruence modules of
$f_{539c}$, $f_{539d}$ and of the depleted forms $(f_{77c})^{\{7\}}$, $(f_{11a})^{\{7\}}$ have orders $288$, $504$, $288$, $504$, and their $7$-parts, $7\in S$, are not covered by the lemma
(\Cref{ch:congruence}).

\begin{theorem}[the paper's Theorem thm:quant]\label{thm:paper:quant}
Let $\ell$ be an odd prime with $\ell\nmid2N$ and $E[\ell]$ irreducible, $S$ the set of odd primes at which $E$ has type $\kod{I}_0^*$ or $\kod{I}_n^*$, $E_S=E\otimes\chi_S$ and $E'_S$ the optimal
curve of its class.
\begin{enumerate}[label=(\roman*)]
\item $v_\ell(\degphi)\ge v_\ell(r_{E_S})+\sum_{q\in S}e_\ell(E,q)\ge\sum_{q\in S}e_\ell(E,q)$.
\item If moreover $\ell\notin S_f$, the first inequality is an equality, and $v_\ell(r_{E_S})=v_\ell(\deg\varphi_{E'_S})$: $v_\ell(\degphi)=v_\ell(\deg\varphi_{E'_S})+\sum_{q\in S}e_\ell(E,q)$.
\end{enumerate}
\end{theorem}

\noindent\emph{Where proved.} \Cref{ch:congruence}, given \Cref{thm:paper:dfg14} and \Cref{thm:paper:main}. (The paper numbers the two statements as one theorem; the book splits them
as (i) and (ii) for reference.)

\begin{proof}[Expanded proof]
Use the set-up above; $\ell\notin S$ as $\ell\nmid N$, so $\ell\ne3$ when $3\in S$.

\emph{Step 1 (the twist map; elementary).} By \Cref{lem:paper:step1}, $v_\ell(r_E)\ge v_\ell(\eta^{\mathrm{Hecke}}(g^S))$, and the congruence ideal is computed in $\T_{\mf m}$: the other local
factors of $\T_\ell$ lie in $\ker\pi$, and the annihilator decomposes along the idempotents.

\emph{Step 2 (Hecke versus cohomological ideal).} \Cref{lem:paper:ddt417} with $L=H^1(X_0(N),\Zl)_{\mf m}$ (the pairing of \Cref{def:paper:dfg-objects} is perfect on the $\mf m$-component, $e_{\mf m}$ being
self-adjoint) and $\pi$: $L[\wp]=M_{g^S,\lambda}$ has rank $2$ and Gram determinant $(\eta^S_{g,\lambda})^2$ (\Cref{prop:paper:integral}; it lies in the $\mf m$-component, as the paper says). Hence $v_\ell(\eta^{\mathrm{Hecke}}(g^S))\ge v_\lambda(\eta^S_{g,\lambda})$, with no freeness hypothesis.

\emph{Step 3 (Diamond, Flach and Guo).} \Cref{thm:paper:dfg14} for $g$ ($\bar\rho\otimes\chi_S$ irreducible, $\ell\nmid2N_S$, $\ell\notin S$): $v_\lambda(\eta^S_{g,\lambda})=v_\lambda(\eta^\emptyset_{g,\lambda})+\sum_{q\in S}v_\ell(L^{\nv}_q(A_g,1)^{-1})$.
Since $A_g=A_f$ and $\delta_q(g)\ge1$, the naive factor of $g$ at $q$ is the true factor, the one that $f$ (with $a_q(f)=0$) omits: the sum is $\sum_{q\in S}e_\ell(E,q)$ (values from
\Cref{lem:paper:euler}, $q=3$ included).

\emph{Step 4 (level $N_S$).} \Cref{lem:paper:hecke} for $E'_S$ ($\ell\nmid2N_S$, $\bar\rho_g$ irreducible): $v_\lambda(\eta^\emptyset_{g,\lambda})=v_\ell(\deg\varphi_{E'_S})=v_\ell(r_{E_S})$, with no freeness hypothesis. With
$v_\ell(\degphi)=v_\ell(r_E)$ this proves (i); the second inequality of (i) holds since $r_{E_S}$ is an integer.

\emph{The equality (ii).} If $\ell\notin S_f$ then $\ell\notin S_{f_{E_S}}$ (twisting preserves absolute irreducibility over $\Q(\sqrt{\ell^*})$, and $N_S\mid N$). Apply \Cref{cor:paper:explicit} to $E$ and to $E'_S$:
$\ad T_\ell E\simeq\ad T_\ell E'_S$ (twist and prime-to-$\ell$ isogenies do not change the adjoint), so the Selmer groups and the Tamagawa factors coincide; at $q\mid N$, $q\notin S$, the twist is
unramified and the $e_\ell$ agree; at $q\in S$, $e_\ell(E_S,q)=0$ ($E_S$ good or multiplicative). The difference is $\sum_{q\in S}e_\ell(E,q)$, and $v_\ell(\deg\varphi_{E'_S})=v_\ell(r_{E_S})$ as $\ell^2\nmid N_S$.
\end{proof}

\noindent\emph{Reading guidance.} The inequality (i) uses, from the Selmer side, only Theorem thm:dfg14, Lemma lem:ddt417 and Lemma lem:hecke, never Theorem thm:dfg37; that is why it
needs no hypothesis over $\Q(\sqrt{\ell^*})$. The equality (ii) squeezes the chain of Steps 1 to 4 between two applications of Theorem thm:main, so every inequality of the proof is then
an equality, Step 2 included, with no multiplicity one at level $N$ or $N_S$; Step 1 is an equality at every $\ell\notin S$ in any case, by \Cref{lem:paper:step1}.

\emph{A worked instance.} \curve{539c1} has type $\kod{I}_3^*$ at $7$ and $\kod{I}_2$ at $11$, $\degphi=288=2^5\cdot3^2$, and surjective mod $3$ image (Cremona's \texttt{galrep} lists only $2\mathrm B$),
so $3\nmid2N$, $E[3]$ is irreducible and $3\notin S_f$ (\Cref{prop:paper:Sf}). Here $S=\{7\}$, $E_S$ is the twist by $\chi_{-7}$, which lies in the class \curve{77c} (type $\kod{I}_3$ at $7$), and
$\deg\varphi_{\curve{77c1}}=6$. The table gives $e_3(E,7)=v_3(7^2-1)=v_3(48)=1$. Part (ii) predicts $v_3(288)=v_3(6)+1=2$, which is right. The inequality of (i) needs none of the Selmer
theorems, only Theorem thm:dfg14: the depleted form $g^S$ of the newform of \curve{77c1} lives at level $539$, and depleting at the multiplicative prime $7$ multiplies the congruence ideal by
the naive factor $1-7^{-2}$, whose $3$-part is $3$.

\paragraph{Remark rem:afterquant.} (1) Theorem thm:quant is the twist-and-lower argument of Section 3 read the other way: the level is raised from $N_S$ to $N$ at the price of
the level-raising Euler factor; for $\ell=3$ a single $\kod{I}_n^*$ prime (necessarily $q\ge5$, since $3\nmid N$) forces $3\mid\degphi$ because $3\mid q^2-1$, and so does an $\kod{I}_0^*$ prime with
$q\equiv1\pmod3$ (the term $v_3(q-1)$); neither is in Theorem thm:qual, and both need $E[3]$ irreducible, the standing hypothesis of the theorem. (2) The equalities, as in the reading
guidance above. (3) For the potentially good types the twist does not lower the level, $f$ is its own depletion, and Theorem thm:dfg14 gives no factor; these terms come only through
Theorem thm:dfg37, and Section 6.7 of the paper shows the hypothesis cannot be dropped there (the $\Ns$ curves fail through a type $\kod{III}$ term).

\section[Paper 5: the boundary]{Section 5 of the paper: the boundary of the hypothesis}\label{sec:paper:boundary}

This section reads the proof of Proposition prop:ns, Proposition prop:ns-structure and the count after it, and Remark rem:ns. They describe the curves at
which the hypothesis of Theorem thm:main fails at $\ell=3$ while $E[3]$ is irreducible, and the correction term of Conjecture conj:euler.

\begin{proposition}[the paper's Proposition prop:ns]\label{prop:paper:ns}
Let $E/\Q$ have $E[3]$ irreducible. Then $\bar\rho_{E,3}$ is absolutely irreducible on $\Gal(\Qbar/\Q(\sqrt{-3}))$ unless its image is conjugate to $N_s(\F_3)$; in that case $\bar\rho_{E,3}\otimes\F_9$ is
induced from a character $\psi$ of order $4$ of $\Gal(\Qbar/\Q(\sqrt{-3}))$, $\bar\rho_{E,3}\simeq\bar\rho_{E,3}\otimes\chi_{-3}$, $H^0(\Q,\ad\bar\rho_{E,3}(1))\ne0$ and $h_3(E)=1$. The $3$-division polynomial is then a
product of two irreducible quadratics (times the leading coefficient $3$), and otherwise irreducible.
\end{proposition}

\noindent\emph{Where proved.} \Cref{ch:galrep}; the subgroups of $\GL_2(\F_3)$ and the commutator argument in \Cref{ch:groups}; Clifford theory, Frobenius reciprocity and
descent of isomorphisms in \Cref{ch:reps}; the division polynomial in \Cref{ch:ellcurves}; the quartic and its lines in \Cref{ch:galois}.

\begin{proof}[Expanded proof]
(1) $\det\bar\rho_{E,3}=\chi_{-3}$ (Weil pairing), which cuts out $\Q(\sqrt{-3})$; so the image of $G_{\Q(\sqrt{-3})}$ is $H\cap\SL_2(\F_3)$, $H$ the image of $\bar\rho$. (2) $H$ has surjective
determinant, acts irreducibly on $\F_3^2$, and contains $\bar\rho(c)$ of determinant $-1$ and trace $0$. Inspection of the subgroups of $\GL_2(\F_3)$ (\Cref{ch:groups}) leaves $\GL_2(\F_3)$
(order $48$), $N_{ns}(\F_3)$ (order $16$) and $N_s(\F_3)$ (order $8$); the nonsplit Cartan subgroup $\F_9^\times$ is excluded, its elements of determinant $-1$ having order $8$ and nonzero
trace. (Why $\bar\rho(c)$ has trace $0$: $c^2=1$ and $\det\bar\rho(c)=-1$, so $\bar\rho(c)$ has eigenvalues $1$ and $-1$. In $\F_9^\times$, acting on $\F_9\simeq\F_3^2$ by multiplication, the determinant
is the norm $x\mapsto x^{1+3}=x^4$; $x^4=-1$ means $x$ has order $8$; then $x^2$ has order $4$, so $x^2\ne-1$, and the trace $x+x^3=x(1+x^2)$ is nonzero.) Complex multiplication plays no role; the paper cites \cite{Zywina15} for the list and derives it directly.
(3) The determinant-one subgroups are $\SL_2(\F_3)$, $Q_8$ and $C_4=\langle J\rangle$, $J=\sm01{-1}0$. (4) The first two act absolutely irreducibly: a subgroup of $\GL_2(\Fbar_3)$ fixing a line lies
in a Borel subgroup, whose commutator subgroup is unipotent, so abelian of exponent $3$; the commutator subgroups $\{\pm1\}$ of $Q_8$ and $Q_8$ of $\SL_2(\F_3)$ are not. (5) $C_4$ is abelian; $J$ has characteristic polynomial $X^2+1$, irreducible over $\F_3$, so the
restriction is irreducible over $\F_3$ and splits over $\F_9$ as $\psi\oplus\psi^{-1}$, $\psi(J)=i$ of order $4$. (6) Frobenius reciprocity gives a nonzero, hence bijective, map
$\Ind\psi\to\bar\rho\otimes\F_9$; an induced representation from an index-two subgroup is invariant under twisting by that subgroup's character, so $\bar\rho\otimes\F_9\simeq(\bar\rho\otimes\chi_{-3})\otimes\F_9$,
and the isomorphism descends to $\F_3$ (Hom spaces commute with extension of scalars). Then $H^0(\Q,\ad\bar\rho(1))\ne0$ and $h_3=1$ by \Cref{lem:paper:h0}. (7) The four roots of $\psi_3$ are the
$x$-coordinates of the four lines of $E[3]$; $N_s(\F_3)$ has two orbits of size two on the lines (the two Cartan lines and the two others), $N_{ns}(\F_3)$ and $\GL_2(\F_3)$ are transitive;
$\psi_3$ has no rational root ($E[3]$ irreducible); so $\psi_3$ factors as two irreducible quadratics in the first case and is irreducible otherwise.
\end{proof}

\noindent\emph{Reading guidance.} Step (4) is the one to check: being non-abelian and acting faithfully does not exclude a fixed line (a Borel subgroup is both), and it is
the commutator subgroup that does. Example (recomputed): \curve{726a1},
$\psi_3=(x^2-x-74)(3x^2+8x+20)$. Below conductor $10{,}000$ the function $\texttt{ns3}$ of the appendix agrees with Cremona's label $\Ns$ on all non-CM optimal curves (\Cref{ch:computation}).
$N_s(\F_3)$ is not a maximal subgroup of $\GL_2(\F_3)$ (it has index $2$ in $N_{ns}(\F_3)$ in the standard models; \Cref{ch:groups}); the paper does not say otherwise, but a reader coming
from large $\ell$ should not assume it.

\begin{proposition}[the paper's Proposition prop:ns-structure]\label{prop:paper:ns-structure}
Let $E/\Q$, with or without complex multiplication, have mod $3$ image conjugate to $N_s(\F_3)=\{\pm1,\pm J,\pm h,\pm w\}$, with $J=\sm01{-1}0$, $h=\mathrm{diag}(1,-1)$ and $w=hJ$; let
$C_s=\{\pm1,\pm h\}$ be the split Cartan subgroup, $C_4=\langle J\rangle=N_s(\F_3)\cap\SL_2(\F_3)$, and $K_1$, $K_2$ the fixed fields of $\bar\rho^{-1}(C_s)$ and $\bar\rho^{-1}(\{\pm1,\pm w\})$, the
fields of definition of the roots of the two quadratic factors of $\psi_3$. Let $R$ be the set of primes $p\ne3$ with $\bar\rho(I_p)=C_4$.
\begin{enumerate}[label=(\roman*)]
\item For every $p\ne3$, $\bar\rho(I_p)\subseteq C_4$. Consequently every $p\ne3$ with $v_p(j)<0$ has $3\mid v_p(j)$, so $t_3(E,p)\ge1$; no $p\ne3$ has type $\kod{IV}$, $\kod{IV}^*$; no $p\ge5$ has
type $\kod{II}$, $\kod{II}^*$.
\item $R$ is the set of primes $p\ne3$ ramified in $K_1$, and also the set of primes $p\ne3$ ramified in $K_2$; it is nonempty. At $p\in R$ the reduction is potentially good,
$\rho_{E,\ell}(I_p)$ is cyclic of order $4$ for every $\ell\ne p$, and $(\ad V_\ell E)^{I_p}$ is a line on which either Frobenius acts by $\chi_{-3}(p)$; $P_p(X)=1-\chi_{-3}(p)pX$ and
$e_3(E,p)=v_3(p-\chi_{-3}(p))\ge1$; in particular $P_2(X)=1+2X$ and $e_3(E,2)=1$ if $2\in R$.
\item If $p\in R$ and $p\ge5$, then $E$ has type $\kod{III}$ or $\kod{III}^*$ at $p$ and $p\equiv\pm1\pmod{12}$. Conversely every $p\ge5$ of type $\kod{III}$ or $\kod{III}^*$ is in $R$.
\item The splitting field of $\psi_3$ is biquadratic, with quadratic subfields $K_1$, $K_2$ and $\Q(\sqrt{-3})$.
\item $a_p(E)\equiv0\pmod3$ for every $p\nmid3N$ with $p\equiv2\pmod3$.
\item If $3\nmid N$, $E$ is ordinary at $3$.
\item $h_3(E)=1$.
\end{enumerate}
\end{proposition}

\noindent\emph{Where proved.} \Cref{ch:galrep}; Kraus's theorem \cite{Kraus90} for the inertia images at $p\ge5$.

\begin{proof}[Expanded proof]
Write $\bar\rho=\bar\rho_{E,3}$. (i) For $p\ne3$, $\chi_{-3}=\det\bar\rho$ is unramified at $p$, so $\bar\rho(I_p)\subseteq N_s(\F_3)\cap\SL_2(\F_3)=C_4$, which has no element of order $3$ or $6$. If
$v_p(j)<0$, on $I_p$ the representation is $\chi(\sigma)\sm1{n\theta(\sigma)}01$ with $n=-v_p(j)$, $\theta$ the $3$-adic tame character, $\chi$ quadratic (\Cref{lem:paper:tamfactor}); on $\ker\chi$, where $\theta$
takes a unit value, this is a unipotent of order $3$ mod $3$ unless $3\mid n$. Type $\kod{IV}$, $\kod{IV}^*$ at $p\ne3$ gives $\bar\rho(I_p)$ of order $3$ (\Cref{lem:paper:IV}); type $\kod{II}$, $\kod{II}^*$ at
$p\ge5$ gives $\rho(I_p)$ cyclic of order $6$ (Kraus), injecting mod $3$. Neither fits in $C_4$.

(iv) is proved before (ii), as the paper signposts: ``We prove (iv) next, because (ii) uses the fields $K_1$ and $K_2$ it describes.'' The splitting field of $\psi_3$ is cut out by the kernel of the action on lines, $\bar\rho^{-1}(\{\pm1\})$, with group $N_s(\F_3)/\{\pm1\}\simeq(\Z/2)^2$;
its quadratic subfields are the fixed fields of the three index-two subgroups containing $-1$: $C_4$ (fixed field $\Q(\sqrt{-3})$, the determinant), $C_s$ (the stabiliser of the
Cartan lines) and $\{\pm1,\pm w\}$ (the stabiliser of the other two lines).

(ii) For $p\ne3$: $\bar\rho(I_p)\subseteq C_4$ and $C_4\cap C_s=C_4\cap\{\pm1,\pm w\}=\{\pm1\}$, so $\bar\rho(I_p)\not\subseteq C_s$ iff $\bar\rho(I_p)=C_4$ iff $\bar\rho(I_p)\not\subseteq\{\pm1,\pm w\}$; that is, $p$ is
ramified in $K_1$ iff $p\in R$ iff $p$ is ramified in $K_2$. $K_1\ne\Q(\sqrt{-3})$, because $C_s\ne C_4$; every quadratic field other than $\Q(\sqrt{-3})$ is ramified at some prime $p\ne3$; so
$R\ne\emptyset$. (This does not use $3\nmid N$, as the paper notes.) For $p\in R$ the reduction is not potentially multiplicative ($\bar\rho(I_p)$ would consist of $\pm$ unipotents, while $C_4$
has elements of order $4$), so $\rho(I_p)$ is finite, injects into $\GL_2(\F_3)$, and is cyclic of order $4$, generated by $\tau$ with $\tau^2=-1$; the same holds on $T_\ell E$ for every
$\ell\ne p$ (Serre and Tate). The trace-zero commutant of $\tau$ is $\Ql\tau$. Frobenius normalises $\langle\tau\rangle$ and conjugates $\tau$ to $\tau^{\pm1}$; the sign can be read mod $3$, where
$\tau$ reduces to $J^{\pm1}$. The elements of $N_s(\F_3)$ outside $C_4$ (the coset $hC_4$, of determinant $-1$) conjugate $J$ to $J^{-1}$, those in $C_4$ commute with it; and
$\det\bar\rho(\Frob_p)=\chi_{-3}(p)$. So Frobenius acts on the line by $\chi_{-3}(p)$; the Euler factor and $e_3$ follow as in \Cref{prop:paper:shapes}(iii), $e_3\ge1$ because
$p\equiv\chi_{-3}(p)\pmod3$. At $p=2$, $\chi_{-3}(2)=-1$.

(iii) For $p\ge5$ the image is tame and $\rho(\Frob_p)\tau\rho(\Frob_p)^{-1}=\tau^p$, which by (ii) equals $\tau^{\chi_{-3}(p)}$; so $p\equiv\chi_{-3}(p)\pmod4$, i.e. $\chi_{-4}(p)=\chi_{-3}(p)$, i.e.
$p\equiv\pm1\pmod{12}$. A cyclic image of order $4$ means semistability defect $4$, $v_p(\Delta)\in\{3,9\}$, types $\kod{III}$, $\kod{III}^*$ (Kraus). Conversely, type $\kod{III}$ or $\kod{III}^*$ at $p\ge5$ gives
$\rho(I_p)$ cyclic of order $4$, injecting mod $3$ onto a subgroup of order $4$ of $C_4$, so $p\in R$.

(v) For $p\equiv2\pmod3$, $p\nmid3N$, $\det\bar\rho(\Frob_p)=-1$, so $\bar\rho(\Frob_p)\in hC_4=\{\pm h,\pm w\}$, all of trace $0$; and $a_p\equiv\tr\bar\rho(\Frob_p)$. (vi) is in \Cref{prop:paper:Sf}.
(vii) is \Cref{prop:paper:ns}.
\end{proof}

\noindent\emph{Reading guidance.} Everything follows from two facts: the determinant $\chi_{-3}$ is unramified away from $3$, which confines inertia to the cyclic group $C_4$; and
the three index-two subgroups of $N_s(\F_3)$ containing $-1$ correspond to the three quadratic subfields of the splitting field of $\psi_3$. The prime $p=2$ is the delicate case of (iii):
Kraus's theorem is not available there, and the paper leaves open whether the type at $2$ is always $\kod{III}$ or $\kod{III}^*$ when $R=\{2\}$.

\paragraph{The count after the proposition.} Among the $90$ optimal non-CM curves with image $\Ns$ below conductor $10{,}000$ ($44$ of them with $3\mid N$), $74$ have a prime $p\ge5$ in $R$,
of type $\kod{III}$ or $\kod{III}^*$ with $p\equiv1$ or $11\pmod{12}$, and $16$ have $R=\{2\}$, of type $\kod{III}$ or $\kod{III}^*$ at $2$ with $P_2=1+2X$. The proposition also covers the $13$ optimal CM curves
with image $N_s(\F_3)$ in that range ($j=-32768$ or $8000$: \curve{121b1}, \curve{256a1}, \curve{256d1}, \curve{1089e1}, \curve{1936f1}, \curve{2304c1}, \curve{2304k1}, \curve{3025a1}, \curve{5929a1},
\curve{6400f1}, \curve{6400r1}, \curve{7744a1}, \curve{7744s1}), $7$ with $R=\{11\}$ and $6$ with $R=\{2\}$, of type $\kod{III}$ at $2$. All these counts were recomputed for the book (PARI 2.17.2,
Cremona's files). That the type at $2$ is always $\kod{III}$ or $\kod{III}^*$ when $R=\{2\}$ is not proved, as the paper says (its last open question).

\begin{proposition}[the paper's Remark rem:ns, as a conditional statement]\label{prop:paper:ns-conditional}
Let $\ell=3$ and the image of $\bar\rho_{E,3}$ be $N_s(\F_3)$, $3\nmid N$, $\Sigma=\{q\mid N\}$. Suppose $H^1_f(\Q,V^D)=0$ (equivalently $H^1_f(\Q,\ad V_3E)=0$). Then
\begin{multline*}
\len H^1_\Sigma(\Q,W)=\len H^1_f(\Q,W)\\
+\sum_{p\in\Sigma}\bigl(v_3(L_p(A_f,1)^{-1})+v_3(\Tam_p(A))\bigr)-h_3(E),\qquad h_3(E)=1.
\end{multline*}
If moreover $v_\lambda(\eta^\Sigma_{f,\lambda})=\len H^1_\Sigma(\Q,W)$, then (eq:identity) holds with $h_3(E)$ subtracted on the right, which is the shape of Conjecture conj:euler.
\end{proposition}

\begin{proof}
\Cref{lem:paper:dfg21} needs only $H^0(\Q,V^D)=0$, automatic as above, and the assumed $H^1_f(\Q,V^D)=0$ (it is a statement of Poitou-Tate duality, valid for any $\ell\notin\Sigma$). Lengths add
along the exact sequence; the last term is $H^0(\Q,W^D)^*$, of length $h_3(E)=1$ (\Cref{lem:paper:h0}); the local terms are given by \Cref{lem:paper:fpr}. For the second statement, Steps 1 and 2 of
the proof of Theorem thm:main need only $\bar\rho$ irreducible, and Step 3 is the added assumption.
\end{proof}

\noindent\emph{Reading remark.} The remark is a conditional statement, and the paper says so: ``The vanishing of $H^1_f(\Q,V^D)$, equivalently of $H^1_f(\Q,\ad V_3E)$, is not known
for these curves: we assume it here''. Neither assumption is proved: \cite[Theorem 2.7]{DFG04} gives the vanishing only for $\ell\notin S_f$, and the $\Ns$ curves at $\ell=3$ are exactly in
$S_f$; the equality of Step 3 would need a Taylor-Wiles argument for a residually dihedral representation induced from $\Q(\sqrt{-3})$ (\Cref{ch:congruence}).

\section[Paper 6: the computations]{Section 6 of the paper: the computations}\label{sec:paper:data}

This section reads the paper's Section 6, subsection by subsection. Its six data tables (Tables tab:uniform, tab:iso, tab:counts, tab:ns, tab:isolated and
tab:oot) are read from their printed rows: for each, what the columns mean, which rows carry the paper's claims, and whether the rows agree with the text; table
entries are copied exactly as printed. Table tab:excess was read with Corollary cor:explicit. Every count below is the paper's unless it is labelled as the book's, and the book's counts (conductor below
$10{,}000$) are checks made with the statements known, not blind tests. How to rerun everything is \Cref{ch:computation}; here we say what each statement means and what it
tests.

\paragraph{6.1 Data and protocol.} The paper uses Cremona's tables \cite{Cremona97} (the \texttt{ecdata} repository at commit \texttt{25cec5e}), which list the curves of conductor
below $500000$ with per-prime local data, modular degrees, isogeny classes, torsion and mod $\ell$ images in the notation of \cite{Sutherland16}; optimality is determined for
every class up to $400000$ and for $359{,}009$ further classes. Four ranges: the working set ($N\le300000$), the hold-out, the extension and the curves outside the tables. The text
gives the order of discovery, from which the status of each statement can be read off:

\begin{center}\small
\begin{tabular}{p{0.32\textwidth}p{0.6\textwidth}}
\toprule
statement & status, as the paper describes it\\
\midrule
(eq:main), Conjecture conj:eis & found on the working set, then tested once on the hold-out and on the extension: a blind test\\
(eq:twist) & replaced an earlier, weaker formulation; formulated on the working set after the counts of (eq:main) on the hold-out and the extension were known, and evaluated on all three ranges: for it the hold-out and the extension are a confirmation; the out-of-table curves are expressly not counted as a blind test of it\\
Conjecture conj:euler (Euler terms) & formulated on the working set; hold-out and extension are a confirmation and not a blind test; the out-of-table curves, generated before the rule was formulated, were evaluated with each version of it, no version failed on them and none was changed because of them: an independent confirmation, not a strictly blind test\\
the weights $w_3$ & fixed afterwards from PARI's factors and the minima at $q=3$; the out-of-table curves are not counted as a blind test of them\\
the correction $h_\ell$ & identified last, from the exceptions; never blind-tested (no out-of-table curve has image $\Ns$)\\
Lemmas lem:euler, lem:twist, Proposition prop:forms, $h_3=1$ & theorems; the data are checks of computations\\
\bottomrule
\end{tabular}
\end{center}

Two further points in the protocol paragraph deserve attention. On the extension range, optimality is determined only for the $359{,}009$ named classes, and for the
others the paper shows that the right side of (eq:rule) at $\ell=3$ is the same for every curve of the class. ``Indeed, when $E[3]$ is irreducible the class has no
$3$-isogeny, so every isogeny in it is a composite of isogenies of prime degree $p\ne3$, and an isogeny of prime degree $p$ preserves the Kodaira type at a prime $q\ne p$ of
potentially good reduction'' \cite[Theorem 5.4(1)]{DokDok15}; so $w_3$, read off the type at $q=3$ when the reduction there is potentially good, is constant on the class
(\Cref{ch:ellcurves}). With a $3$-isogeny, where (eq:rule) does not apply at $\ell=3$, $w_3$ can vary; the paper's count below $10{,}000$ ($w_3$ constant on all $1{,}295$
classes without a $3$-isogeny, ``as it must be'', varying in $828$ of the $1{,}293$ with one, for instance \curve{54a}, of types $\kod{IV}^*$, $\kod{II}^*$, $\kod{II}$ at $3$) was reproduced for
the book. The other terms of (eq:rule) are invariant under isogenies of degree prime to $\ell$ (\Cref{ch:ellcurves}). And the PARI versions: the paper used
PARI/GP 2.15 for the large ranges and 2.17.2 for the checks below $10{,}000$, and it records that 2.17.2 returns $q_E^5$ in place of $q_E$ for \texttt{ellinit(E, O(2\^{}k)).tate} when
$v_2(j_E)=-1$; no count is affected, as the Tate parameter at $2$ enters none (\Cref{ch:reduction,ch:computation}).

\paragraph{6.2 The rule on whole ranges (Table tab:uniform).} The table has twenty rows, one for each of the four ranges (working set, hold-out, extension, out
of table) and each $\ell\in\{3,5,7,11,13\}$. How to read the columns: ``curves'' is the population, the optimal curves of the range with $E[\ell]$ irreducible and
not CM (so it depends on $\ell$); ``(eq:rule) $>0$'' counts those with a positive bound; ``viol.'' the violations of (eq:rule); ``equality'' the share of the
curves with a positive bound at which $v_\ell(\degphi)$ equals the bound; ``viol.\ without $h$'' the violations once the correction $h_\ell$ is dropped; and the
last column gives the same three numbers (positive bounds, violations, equality) for the ``III-free'' rule, with the type $\kod{III}$, $\kod{III}^*$ terms left out and no
correction. The first row reads: working set, $\ell=3$, $1{,}233{,}729$ curves, $867{,}578$ with a positive bound, $0$ violations, equality $57.5\%$, $364$ violations without
$h$, III-free $861{,}074$ / $0$ / $56.8\%$. Two columns carry the paper's claims. The column ``viol.'' is $0$ in all twenty rows. The column ``viol.\ without $h$'' is $0$
at every $\ell\ge5$ and reads $364$, $68$, $43$ and $0$ at $\ell=3$ on the four ranges, the failing curves of Table tab:ns (6.7 below); the III-free column has $0$
violations in every row. The equality column is the sharpness of Section sec:sharp: on the working set $57.5\%$, $75.6\%$, $84.1\%$, $91.6\%$ and $92.4\%$ at
$\ell=3,5,7,11,13$; at $\ell=3$ it reads $57.5\%$, $54.6\%$, $54.0\%$ and $52.5\%$ on the four ranges, the last for the $828$ out-of-table curves with a positive bound. The
III-free rule bounds fewer curves ($861{,}074$ against $867{,}578$ on the working set at $\ell=3$) and is attained slightly less often: leaving out the type $\kod{III}$ terms costs
a little sharpness and removes the need for a correction. These rows agree with the text, which quotes $867{,}578$, the five equality rates and the $52.5\%$ of $828$.
That the III-free rule needs no correction is explained by \Cref{prop:paper:ns-structure}: every $\Ns$ curve has a prime of $R$, whose term $e_3\ge1$ is a type $\kod{III}$ term at
every $p\ge5$. Book's range ($\ell=3$): in the base population of $33{,}851$ optimal non-CM curves below $10{,}000$ with $E[3]$ irreducible, of which $20{,}317$ have a positive bound, there is no
violation of (eq:rule) or of the III-free rule (\Cref{ch:computation}).

\paragraph{6.3 Calibration and sharpness.} The argument: among working-set curves with bound $0$, $53.2\%$ have $3\nmid\degphi$; if $v_3(\degphi)$ were independent of the local data, the
$867{,}578$ curves with a positive bound would give about $605{,}000$ violations; there are none. (A sanity check, the book's arithmetic: the share $53.2\%$ alone, applied to the
$867{,}578$ curves, gives the lower bound $0.532\times867{,}578\approx461{,}551$.) The paper adds the right reading: ``This shows that the statement is strong, far from vacuous, not that it is
valid; the figure depends on the reference population.'' On the book's range the same calculation gives about $14{,}275$ expected violations, with a stratified permutation null of
about $7{,}600$ (\Cref{ch:computation}): the number changes by a factor of two with the reference population, the conclusion not at all. Under the hypothesis of Theorem thm:main the
excess over the rule is the Selmer length plus the Tamagawa excess (\Cref{cor:paper:explicit}), as the paper says. The odd-part statistic ($1.85\%$ on the working set, $17.1\%$ below
$10{,}000$) is range dependent, as the paper notes; below $10{,}000$ the smallest prime factor of the quotient is $3$ in $58.7\%$ of the cases, it divides $N$ in $47.5\%$ of the cases with a
nontrivial quotient, and some prime factor of the quotient divides $N$ in $51.8\%$. The book reproduced these three figures.

\paragraph{6.4 The Tamagawa factors.} The test of the refined bound with the extra term $\min(v_\ell(n),v_\ell(q-1),v_\ell(c(u)))$: at $\ell=3$, $46{,}222$ working-set curves, $12{,}095$ with a positive
extra term, no exception, equality rising from $45.0\%$ to $60.6\%$; at $\ell=5$, $10{,}940$, $1{,}646$, from $67.0\%$ to $78.7\%$; including the $\kod{I}_n^*$ primes below $10{,}000$: $1{,}050$ curves at $\ell=3$
(from $56.8\%$ to $68.0\%$) and $252$ at $\ell=5$. The book reproduced the small-range numbers (\Cref{ch:blochkato,ch:computation}). Reading: under the hypothesis, (eq:identity) is an equality,
so a violation is impossible and the equality rate is the proportion of curves with $H^1_f(\Q,W)=0$; the test checks the Tamagawa term of the identity, as the paper says. At $q=2$ the
extra term vanishes ($v_\ell(2-1)=0$), which is why the Tate parameter at $2$ is never needed.

\paragraph{6.5 Each term in isolation (Table tab:iso).} The table is printed for $\ell=3$ on the working set. Rows are the Kodaira types $\kod{II}$, $\kod{II}^*$, $\kod{IV}$, $\kod{IV}^*$, $\kod{III}$,
$\kod{III}^*$, $\kod{I}_0^*$, $\kod{I}_n^*$ of the one additive prime $q\ge5$; the column pairs are $e=e_3(E,q)=1,2,3,4$, and in each pair ``min'' is the minimum of $v_3(\degphi)$ over the cell and ``$n$''
its number of curves (curves with $E[3]$ irreducible and no multiplicative $3$-contribution). The claim is carried by the ``min'' columns: in every printed cell the
minimum equals $e$, that is $1$, $2$, $3$, $4$ along each row ($\kod{III}$ and $\kod{III}^*$ have no curve at $e=4$, and the cell is printed as a dash). The ``$n$'' columns say how much each cell weighs:
the largest are $\kod{I}_n^*$ at $e=1$ ($30{,}497$ curves), $\kod{II}$ at $e=1$ ($8{,}327$) and $\kod{IV}^*$ at $e=1$ ($8{,}221$); the cells at $e=4$ are small ($\kod{II}$ and $\kod{IV}^*$: $2$;
$\kod{II}^*$ and $\kod{IV}$: $1$; $\kod{I}_0^*$: $247$; $\kod{I}_n^*$: $13$). The rows come in pairs: $\kod{II}$ and $\kod{IV}^*$ have $821$, $59$, $2$ curves at $e=2,3,4$ and $8{,}327$, $8{,}221$ at
$e=1$; $\kod{II}^*$ and $\kod{IV}$ have $296$, $12$, $1$ and $4{,}184$, $4{,}290$; $\kod{III}$ and $\kod{III}^*$ have $1{,}271$, $76$, $9$ in both. These are the pairs exchanged by the twist at $q$ in
Lemma lem:twist(ii), which keeps the conductor exponent at $q\ge5$ (a reading of the book; the paper does not comment on it). What the table does not print, the
text supplies: the cells with $e=0$ ($\kod{III}$: $5{,}534$ curves, $\kod{I}_0^*$: $5{,}394$), the $\kod{I}_0^*$ cells up to $e=6$, the values at $\ell=5,7$ (``The same holds at $\ell=5$ and $7$'', in the
caption) and, for $\kod{I}_n^*$, that the excess over $e_\ell$ has minimum $v_\ell(n)$; the printed cells agree with the text's ``In every cell, for $\ell=3,5,7$, the minimum equals
$e_\ell$''. The book found the same pattern cell
by cell below $10{,}000$ (\Cref{ch:computation}). Multiplicative primes alone show no Euler term, for every value of $v_\ell(q^2-1)$ up to $6$; at $\ell=3$ the value $0$ occurs only for $q=\ell=3$, as
the paper notes, since $3\mid q^2-1$ for every other $q$. This is the data side of the paper's thesis: the naive factor at a multiplicative prime is the true one, so nothing is omitted. The
two-prime observation ($9\mid\degphi$ for all $2{,}502$ curves with two odd primes of type $\kod{II}$ or $\kod{II}^*$; it is the entry $2$ ($2{,}502$) of Table tab:isolated, 6.8 below) is consistent with the additivity of the $e_3$ terms.

\paragraph{6.6 The two corollary forms (Table tab:counts).} The table has eighteen rows: $\ell=3,5,7,11,13$ on the working set ($N\le300000$), the hold-out
($300000<N\le400000$) and the extension ($400000<N\le500000$), and $\ell=3,5,7$ out of table ($500000<N\le4\cdot10^6$). How to read the columns: ``curves'' is the size of the
range, the same for every $\ell$ ($1{,}311{,}066$, $429{,}936$, $359{,}009$, $893$; it is not the population of Table tab:uniform, which keeps only the curves with $E[\ell]$
irreducible and not CM); ``(eq:main) $>0$'' and ``(eq:twist) $>0$'' count the curves with $E[\ell]$ irreducible, not CM, with a positive right-hand side; the next two columns are
the violations; the last column gives the Eisenstein curves with a positive bound in (eq:main) and the violations of Conjecture conj:eis. The claim is carried by the three
violation columns, which are $0$ in all eighteen rows. The first row: $1{,}311{,}066$ curves, $598{,}833$ with (eq:main) $>0$, $705{,}563$ with (eq:twist) $>0$, $0$, $0$, and
$67{,}845$ / $0$. The twisted form bounds more curves than (eq:main) in every row (out of table at $\ell=7$: $1$ against $0$), and the rule bounds more still ($867{,}578$ in
Table tab:uniform). The Eisenstein column on the working set reads $67{,}845$, $4{,}454$, $759$, $0$, $26$ at $\ell=3,\dots,13$, and $0$ / $0$ out of table. The equality rates
of the two forms ($39.4\%$, $71.1\%$, $81.2\%$, $91.5\%$ for (eq:twist) at $\ell=3,5,7,13$, and $36.9\%$, $69.7\%$, $80.4\%$, $91.2\%$ for (eq:main)) are in the text, not in the
table. The rows agree with the text: ``There are none'' for the violations, ``$867{,}578$ curves instead of $705{,}563$ at $\ell=3$'', and the $67{,}845$ Eisenstein curves of
Section sec:eis. With
\Cref{prop:paper:forms} the absence of violations under the hypothesis of Theorem thm:main is a theorem; the data test the forms outside it, chiefly at the curves with $\ell\mid N$. The
paragraph lists the failures of the hypothesis on the working set: the $1{,}170$ non-CM curves with image $\Ns$ at $\ell=3$ and the curves with $\ell\mid N$; by Proposition prop:Sf there are no
others among the non-CM curves with $E[\ell]$ irreducible. Its parenthesis on the CM curves with $j\in\{-32768,8000\}$, which also have image $N_s(\F_3)$ and fail at $\ell=3$ when $3\nmid N$
($20$ below $10{,}000$) and which the \texttt{galrep} file leaves unlabelled, matches \Cref{ch:galrep}.

\paragraph{6.7 The boundary: the $\Ns$ curves (Table tab:ns).} The table has one row per range. How to read the columns: ``$\Ns$'' is the number of optimal
non-CM curves with mod $3$ image $\Ns$ and $E[3]$ irreducible; ``with III term'' those with a positive type $\kod{III}$ term; ``failing'' those violating (eq:thm) with all types
and no correction, with their share; ``others failing'' the curves without image $\Ns$ that violate it; ``$h_3=1$'' the $\Ns$ curves with $h_3(E)=1$; the next column the
residues modulo $12$ of the type $\kod{III}$ primes of the failing curves, or $q=2$; and ``shortfall $1$'' the failing curves that miss the bound by exactly one. The rows:
\begin{center}\footnotesize
\begin{tabular}{lrrrrrrr}
\toprule
set & $\Ns$ & III term & failing & others & $h_3=1$ & $q\equiv1$ / $11\ (12)$ / $q=2$ & shortfall $1$\\
\midrule
working set & 1,170 & 1,170 & 364 (31\%) & 0 & 1,170 of 1,170 & 140 / 160 / 72 & 364 of 364\\
hold-out & 290 & 290 & 68 (23\%) & 0 & 290 of 290 & 34 / 30 / 4 & 68 of 68\\
extension & 237 & 237 & 43 (18\%) & 0 & 237 of 237 & 16 / 25 / 2 & 43 of 43\\
out of table & 0 & 0 & 0 (0\%) & 0 & 0 of 0 & 0 / 0 / 0 & 0 of 0\\
\bottomrule
\end{tabular}
\end{center}
Every column carries a claim: ``III term'' equals ``$\Ns$'' (every $\Ns$ curve has a positive type $\kod{III}$ term); ``others'' is $0$ (only $\Ns$ curves fail, which is what
``viol.\ without $h$'' of Table tab:uniform counts); ``$h_3=1$'' is the whole population (Lemma lem:h0, a theorem); and the shortfall is always $1$, so the correction
$h_3=1$ is exactly enough. The residue column counts primes, not curves: on the working set $140+160+72=372$ exceeds the $364$ failing curves, so some failing
curves have two such primes. The table lists only the classes $1$ and $11$ modulo $12$ at $q\ge5$, and Proposition prop:ns-structure shows that no other class can occur ($p\equiv\pm1\pmod{12}$).
The rows agree with the text. The paper: $364$, $68$, $43$ and $0$ violations of (eq:thm), with all types and no correction, on the four ranges, all at $\ell=3$ and
all at $\Ns$ curves; every $\Ns$ curve has a type $\kod{III}$ prime with positive $e_3$; the inequality fails for $31\%$, $23\%$, $18\%$ of them, always by one; each violation involves a prime
$p\equiv-1\pmod3$ with $t_3(E,p)\ge1$, which the paper marks as ``an observation, not proved''. Proposition prop:ns-structure explains part of it (every potentially multiplicative $p\ne3$ of
an $\Ns$ curve has $t_3\ge1$). For the curves with $3\mid N$ the count depends on the term $t_3(E,3)$: below $10{,}000$, $36$ violations with it, as (eq:thm) is written, and $24$ without, the
twelve curves that differ (\curve{726a1}, \curve{726f1}, \curve{2178b1}, \curve{2178h1}, \curve{3840h1}, \curve{3840k1}, \curve{3840o1}, \curve{3840r1}, \curve{5070h1}, \curve{5070p1}, \curve{5808y1},
\curve{5808z1}) all having $t_3(E,3)=1$ from a type $\kod{I}_n$ or $\kod{I}_n^*$ at $3$; the weights $w_3$ change nothing. The book reproduced these counts. Example: \curve{726a1}, types $\kod{I}_6$ at
$2$, $\kod{I}_3$ at $3$, $\kod{III}$ at $11$, $\degphi=144$: the bound of (eq:thm) as written is $1+1+1=3>2$, by one, through $p=2\equiv-1\pmod3$ with $t_3(E,2)=1$; without $t_3(E,3)$ it is $2$
(\Cref{ex:paper:726a1} treats \curve{3840h1}). The value $h_3(E)=1$ for every $\Ns$ curve is \Cref{lem:paper:h0}, a theorem.

\paragraph{6.8 The wild primes (Table tab:isolated).} The table isolates the additive types at $\ell=3$ on the working set. Each entry is the minimum of $v_3(\degphi)$ with the
number of curves in parentheses. The header reads ``weight 1, 2, 3'' for every row, but the caption gives it two meanings: in the five untwisted rows ($\kod{IV}$, $\kod{IV}^*$,
$\kod{III}$, $\kod{III}^*$, $\kod{I}_0^*$, at any $q$) it is the number of primes of that type, all additive primes having that type; in the four twisted rows it is the total weight of
(eq:twist) from the primes named. The rows:
\begin{itemize}
\item untwisted: $\kod{IV}$: $1$ ($25{,}414$), $2$ ($969$), $3$ ($11$); $\kod{IV}^*$: $1$ ($30{,}168$), $2$ ($2{,}268$), $3$ ($45$); $\kod{III}$: $0$ ($34{,}904$), $0$ ($2{,}208$), $0$ ($34$); $\kod{III}^*$: $0$ ($25{,}574$),
$0$ ($1{,}307$), $0$ ($29$); $\kod{I}_0^*$: $0$ ($42{,}924$), $0$ ($3{,}093$), $0$ ($53$);
\item twisted, odd $q$: $\kod{II}$, $\kod{II}^*$: $1$ ($56{,}509$), $2$ ($2{,}502$), $3$ ($12$); $\kod{I}_n^*$ with weight $v_3(n)$: $1$ ($28{,}175$), $2$ ($3{,}453$), $3$ ($89$);
\item twisted, $q=2$: $\kod{II}$, $\kod{II}^*$: $0$ ($44{,}672$), with no curve of weight $2$ or $3$ (a dash in the table); $\kod{I}_n^*$ with weight $v_3(n)$: $0$ ($21{,}176$), $1$ ($5{,}329$), $1$ ($342$).
\end{itemize}
Reading guidance. Where the minimum equals the column number ($\kod{IV}$, $\kod{IV}^*$ and both odd-$q$ twisted rows), each prime, or each unit of weight, forces one factor $3$ in
$\degphi$, as the terms of (eq:twist) say. The rows with minimum $0$ ($\kod{III}$, $\kod{III}^*$, $\kod{I}_0^*$) carry no forced factor when the curves are grouped by the number of
primes; their term is the Euler term $e_3$, which can vanish, and Table tab:iso shows the minimum $e$ once the curves are grouped by $e$. The paper's claims about the
prime $2$ are in the last two rows: types $\kod{II}$, $\kod{II}^*$ at $2$ force nothing, and a type $\kod{I}_n^*$ at $2$ forces nothing when $3\,\|\,n$ and only one factor $3$ when $9\mid n$
or $27\mid n$, where the odd-$q$ weights would give $2$ and $3$. The rows agree with the text: ``$v_3(\degphi)=0$ for $14{,}182$ of the $44{,}672$ curves whose only candidate is
such a prime'' (the $14{,}182$ is in the text only), ``although with $9\mid n$ the minimum is $1$'', and the $2{,}502$ curves of Section sec:iso. Its explanation is Definition def:terms with PARI's
factor, the factor having one of the four shapes of \Cref{prop:paper:shapes}: a type $\kod{II}$ or $\kod{II}^*$ at $2$ \emph{usually} has $P_2=1$ (among the optimal non-CM curves below
$10{,}000$, $1{,}092$ of the $1{,}537$ with type $\kod{II}$ at $2$ and $514$ of the $779$ with type $\kod{II}^*$; \Cref{ch:computation}); a type $\kod{IV}$ or $\kod{IV}^*$ at $2$ has $e_3=1$ by Lemma lem:IV;
and at $q=2$ the right invariant of a type $\kod{I}_n^*$ is $-v_2(j)$, not the index $n$. The weights $w_3$ at $q=3$ are the observed minima ($1$ for $\kod{II},\kod{IV},\kod{III}^*$; $2$ for
$\kod{II}^*,\kod{IV}^*$; $0$ for $\kod{III},\kod{I}_0^*$), reproduced exactly on the book's range (\Cref{ch:computation}). For $\ell=5,7$ no term is needed at $q=\ell$ or at $q=2,3$ beyond Definition
def:terms, and the term at $2$ can be positive (\curve{592c1}); the paper adds that this is ``true but not sharp'' at $q=\ell$, by Remark rem:oq1.

\paragraph{Remark rem:oq1 (a twist relation at $q=\ell$; data only).} For optimal non-CM curves below $10{,}000$ with $\ell\mid N$, $\ell\in\{3,5,7,11,13\}$, the paper compares $v_\ell(\degphi)$ with
$v_\ell$ of the degree of the optimal curve of the class of the twist $E^{(\ell)}$ by $\chi_{\ell^*}$: (a) for type $\kod{II}^*$, $\kod{III}^*$ or $\kod{IV}^*$ at $\ell$ the difference is exactly $1$ in all
$4{,}412$ cases ($4{,}251$ with $E[\ell]$ irreducible); (b) for type $\kod{I}_n^*$ at $\ell$ it is $0$; (c) for type $\kod{I}_0^*$ at $\ell$ with $E[\ell]$ irreducible it is $v_\ell((\ell+1)^2-a_\ell(E^{(\ell)})^2)$
in all $1{,}915$ cases, failing in $13$ cases with $E[\ell]$ reducible (\curve{126a1} at $\ell=3$, \curve{275b1} at $\ell=5$). How to read it: at $\ell=3$, (a) is the pattern of the weights $w_3$, each
starred type weighing one more than its partner under the exchanges of Lemma lem:twist(ii); in (b) and (c) the difference is the numerator of the table of Section 1.1 evaluated at
$q=\ell$. Nothing here is proved, the observations were made on the whole range with no hold-out, and the paper offers no explanation; the book reproduced the counts
(\Cref{ch:computation}).

\paragraph{6.9 Out of table (Table tab:oot).} $893$ curves with prescribed local types and $500000<N\le4\cdot10^6$, kept when the isogeny class is trivial (then the curve is optimal and $E[\ell]$ is
irreducible for every $\ell$, since a rational $\ell$-subgroup would give an isogeny; \Cref{ch:ellcurves}) and not CM. Conjecture conj:euler and both forms hold for all of them, with equality in
(eq:rule) for $52.5\%$ of the $828$ with a positive bound at $\ell=3$. The table itself gives, per family of prescribed types, the number of curves, the largest conductor
and the median time in seconds to compute one modular degree with \texttt{ellmoddegree}: $\kod{II},\kod{II}$: $98$, $3{,}968{,}800$, $0.4$; $\kod{II},\kod{II},\kod{II}$: $1$, $2{,}832{,}200$, $0.2$;
$\kod{IV}$: $334$, $3{,}988{,}012$, $1.1$; $\kod{IV}^*$: $104$, $3{,}994{,}600$, $4.9$; ``IVmixed'' (the text's ``mixed types''): $266$, $3{,}983{,}504$, $1.3$; random: $90$, $2{,}710{,}716$, $32.1$; all: $893$,
$3{,}994{,}600$, $1.5$. Reading guidance: the table describes the sample, not the test; the results on these curves are the last rows of Tables tab:uniform and tab:counts
(``out of table''), and the last row of Table tab:ns records that none of them has image $\Ns$. The families add up to $893$, the largest conductor is below $4\cdot10^6$, and
the cost column shows why the sample is small: a random curve costs about twenty times the median. These rows agree with the text. For the Euler terms of the rule these curves are an independent confirmation: they ``were generated before the rule was formulated and were evaluated with each version of it; no version failed on them and none was changed because of them''. As the paper says, they are not counted as a blind test of (eq:twist) or of the weights $w_3$.

\paragraph{6.10 The Eisenstein case.} At $\ell=3$ on the working set, $75{,}430$ optimal non-CM curves with a rational $3$-isogeny, $2{,}777$ with $D>0$, $D=2$ for $22$: seventeen with $E[3]\simeq\Z/3\oplus\mu_3$
(image 3Cs.1.1, the subgroup of order $2$ of the split Cartan subgroup that fixes a point) and five with a point of order $9$ (image 3B.1.1). The seventeen have torsion $\Z/3$ or $\Z/6$ (for
instance \curve{14a1}, torsion $\Z/6$, types $\kod{I}_6$ at $2$ and $\kod{I}_3$ at $7$, $\degphi=1$, so $D=v_3(6)+v_3(3)-0=2$); with a rational point of order $3$, Borel image and no point of order $9$, $D\le1$
(\curve{714i1}, torsion $\Z/9$, $\degphi=1080=2^3\cdot3^3\cdot5$, has $D=2$). The book recomputed \curve{14a1}, \curve{714i1}, \curve{20a1} ($D=1=v_3(\#E(\Q)_{\mathrm{tors}})$) and the other examples
(Section~\ref{sec:paper:verifications}). The labels are used as in Cremona's \texttt{galrep} file, with their intrinsic descriptions; the paper says it has not compared them with the tables
of \cite{Sutherland16} themselves. The two sources of a deficit of $2$ never combine, and the paper turns this into a theorem:

\begin{proposition}[the paper's Proposition prop:no9]\label{prop:paper:no9}
Let $E/\Q$ be without complex multiplication, with $E[3]\simeq\Z/3\oplus\mu_3$. Then $E(\Q)$ has no point of order $9$. Consequently, for every non-CM $E/\Q$,
$v_3(\#E(\Q)_{\mathrm{tors}})+[E[3]\simeq\Z/3\oplus\mu_3]\le2$, and $v_\ell(\#E(\Q)_{\mathrm{tors}})\le1$ for $\ell\ge5$.
\end{proposition}

\noindent\emph{Where proved.} \Cref{ch:ellcurves}; the input that a rational cyclic $27$-isogeny forces complex multiplication is Ligozat's determination of the rational points of $X_0(27)$
\cite{Ligozat75}, stated there; Mazur's theorem on torsion \cite{Mazur77} is stated in \Cref{ch:ellcurves}.

\begin{proof}[Expanded proof]
Suppose $P\in E(\Q)$ has order $9$. Then $3P$ is a rational point of order $3$, so $\langle3P\rangle$ is the line $\Z/3$ of $E[3]$ ($\mu_3$ has no nonzero rational point), and $\langle P\rangle\cap E[3]=\langle3P\rangle$
meets $\mu_3$ trivially. Let $\phi\colon E\to E'=E/\mu_3$ (rational, as $\mu_3$ is $G_\Q$-stable) and $\hat\phi$ its dual. From $\phi\circ\hat\phi=[3]$, $\hat\phi(E'[3])\subseteq\ker\phi=\mu_3$, with equality since
$\#\hat\phi(E'[3])=9/\#\ker\hat\phi=3$. The subgroup $C=\hat\phi^{-1}(\langle P\rangle)\subset E'$ is $G_\Q$-stable of order $27$ ($\hat\phi$ has degree $3$ and is onto on points). If $C\supseteq E'[3]$ then
$\mu_3=\hat\phi(E'[3])\subseteq\langle P\rangle$, a contradiction; a subgroup of order $27$ of $E'[27]\simeq(\Z/27)^2$ not containing $E'[3]$ is cyclic. So $E'\to E'/C$ is a rational cyclic $27$-isogeny. By
Ligozat, $X_0(27)$ has genus $1$ and its non-cuspidal rational points have $j=-2^{15}\cdot3\cdot5^3$, CM by the order of discriminant $-27$; so $E'$, hence $E$, has CM, a contradiction. The
consequences: by Mazur's list, $v_3(\#E(\Q)_{\mathrm{tors}})\le2$ with equality only for $\Z/9$, and $v_\ell\le1$ for $\ell\ge5$.
\end{proof}

\noindent\emph{Reading guidance.} The proof is a descent of a cyclic subgroup through the dual isogeny: the point of order $9$ and the subgroup $\mu_3$ together produce a
rational cyclic $27$-isogeny on $E/\mu_3$, and such an isogeny exists only with complex multiplication. The paper cites \cite{Ligozat75} for that fact and \cite[Theorem 1]{Kenku} for
the list of levels $N$ with $Y_0(N)(\Q)\ne\emptyset$. The remaining paragraph of Section 6.10 (deficits at $\ell=5,7$ only with rational $\ell$-torsion; the deficit against (eq:twist) not
controlled by torsion, which is why Conjecture conj:eis is stated for (eq:main)) is data. In the Eisenstein case the additive terms do not contribute reliably: \curve{539a1} (image 3B, type
$\kod{I}_2^*$ at $7$) has $e_3=1$ and $v_3(\degphi)=v_3(320)=0$.

\section[Paper 7, 8: the prime 2, open questions]{Sections 7 and 8 of the paper: the prime $2$ and the open questions}\label{sec:paper:two-open}

This section reads the paper's Section 7 (the prime $2$) and its list of open questions, and says where the book bears on each.

\paragraph{Section 7: the prime $2$.} Every theorem of the paper needs $\ell$ odd: Lemma lem:tamfactor uses that $2$ is a unit, the trace form needs $\ell\ne2$, the real place contributes
$H^1(\R,T^D)$ killed by $2$, and $S_f$ contains the prime $2$. Section 7 is data and pointers: a $2$-adic inequality with the types $\kod{III},\kod{III}^*,\kod{I}_0^*,\kod{I}_n^*$ counted once, holding on
the three ranges, and failing when $\kod{I}_0^*$ is counted twice; the book's range shows the same (\Cref{ch:computation}). The theorem of Calegari and Emerton \cite{CalegariEmerton09} is
summarised in the numbering of its arXiv version, as the paper says: odd modular degree implies at most two odd bad primes, even analytic rank, and either a rational
$2$-torsion point, or prime conductor with supersingular reduction at $2$ and $E(\R)$ connected, or complex multiplication with $N\in\{27,32,49,243\}$. Watkins's conjecture
$2^{\rank E(\Q)}\mid\degphi$ is stated in \Cref{ch:moddeg}.

\paragraph{Open question 1 (the weights at $q=\ell$).} The paper asks to extend Theorem thm:main to $\ell\mid N$ and to identify the weights $w_3$ with a local term at $\ell$. What fails
there: at $\ell\mid N$, Agashe-Ribet-Stein can fail ($\ell^2\mid N$: \curve{99a1}), $\ell\in S_f$ (Fontaine-Laffaille needs $\ell\nmid N$), and the Euler factor at $q=\ell$ is not the relevant
local quantity. The paper's lead is Remark rem:oq1 (read in Section~\ref{sec:paper:data}): it reproduces the pattern of the $w_3$ and suggests a weight $1$ for the types $\kod{II}^*$,
$\kod{III}^*$, $\kod{IV}^*$ at $q=\ell\ge5$; added to (eq:rule), that weight gives no violation below $10{,}000$ at $\ell=5,7,11,13$ and is attained by most of these curves.

\paragraph{Open question 2 (the $\Ns$ curves).} Two inputs are missing, as \Cref{prop:paper:ns-conditional} shows: $H^1_f(\Q,\ad V_3E)=0$ and the equality
$v_\lambda(\eta^\Sigma_{f,\lambda})=\len H^1_\Sigma(\Q,W)$ for a residual representation induced from $\Q(\sqrt{-3})$. The data are consistent: failures by exactly one, and equality in Theorem
thm:quant in the $8$ pairs with image $\Ns$ below $10{,}000$, which lie outside its equality clause.

\paragraph{Open question 3 (the Selmer group).} With Lemmas lem:tamfactor and lem:tamadd, the excess of $v_\ell(\degphi)$ over the explicit local bound is the Selmer length
(\Cref{cor:paper:explicit}, Table tab:excess), and what remains is to compute that Selmer group independently, for instance for \curve{1184f1} and \curve{1748e1} at $\ell=3$
(\Cref{ch:blochkato}).

\paragraph{Open questions 4 to 6.} The Eisenstein deficit and the cuspidal subgroup (\cite{Mazur77}; \Cref{ch:moddeg}); the prime $2$ and Watkins's conjecture; and the smaller point
whether, in Proposition prop:ns-structure, the type at $2$ is always $\kod{III}$ or $\kod{III}^*$ when $R=\{2\}$, true for all $16$ non-CM and all $6$ CM optimal curves below $10{,}000$ where it
occurs (\Cref{ch:galrep}). The book adds nothing beyond the statements and \Cref{prop:paper:no9}.

\section[Paper appendix: \texttt{rule}]{The appendix of the paper: the function \texttt{rule}}\label{sec:paper:appendix}

This section reads the paper's appendix (reproducibility and the PARI functions) and its statement on the use of an AI system. The program itself is read line by line in
\Cref{ch:computation}; here we say what it computes and how to use it.

\begin{lstlisting}[language=GP]
rule(E, l, {iii = 0}) =
{ my(R = ellglobalred(E), P = R[4][,1], K = R[5], S = lfunsympow(E, 2), b = 0, j = E.j);
  for(i = 1, #P,
    my(q = P[i], k = K[i][2], e = 0, t = 0, vj = valuation(j, q));
    if(q != l,
      my(F = 1/lfuneuler(S, q), r = subst(F, x, 1/q^2));
      if(k >= 5, r /= (1 - 1/q^2));
      if(abs(k) != 3 || iii, e = valuation(r, l)),
      if(l == 3, e = if(k == 4 || k == 2, 1,
                     if(k == -4 || k == -2, 2, if(k == -3, 1, 0)))));
    if(vj < 0, t = valuation(-vj, l));
    b += e + t);
  if(iii && l == 3, b -= ns3(E));
  [valuation(ellmoddegree(E), l), b]; }
ns3(E) = my(F = factor(elldivpol(E, 3))); (#F[,1] == 2 && poldegree(F[1,1]) == 2);
\end{lstlisting}

\paragraph{What it computes.} (The program is copied from the paper, comments omitted; the paper's comment line after \texttt{ns3} explains the test: $\psi_3$ has degree $4$ and is squarefree, so
two factors, one of them quadratic, means two quadratics.) $P$ is the list of bad primes and $k$ PARI's Kodaira code at $q$. For $q\ne\ell$, $F=1/\texttt{lfuneuler}(S,q)$
is $P_q(x)$ (for \curve{121a1}, \texttt{lfuneuler(S, 11)} is $1/(11x+1)$), $r=P_q(q^{-2})=L_q(A_f,1)^{-1}$, divided at a multiplicative prime ($k\ge5$) by the naive factor $1-q^{-2}$; so $e=e_\ell(E,q)$,
except that the type $\kod{III}$ terms ($|k|=3$) are left out unless \texttt{iii} is set. At $q=\ell=3$ the code puts $e=w_3(E)$ (codes $4,2,-3$ give $1$; $-4,-2$ give $2$; others $0$), even when
\texttt{iii} is not set. $t=t_\ell(E,q)$ at every bad $q$, $q=\ell$ included. With \texttt{iii} set and $\ell=3$, $\texttt{ns3}(E)$ is subtracted; it is $1$ exactly when $\psi_3$ is a product of two
irreducible quadratics, that is, for $E[3]$ irreducible, exactly for image $N_s(\F_3)$ (\Cref{prop:paper:ns}), and then it equals $h_3(E)$. So \texttt{rule(E, l, 1)} returns $[v_\ell(\degphi),\text{right side
of (eq:rule)}]$ and \texttt{rule(E, l)} the III-free bound, provided $E$ is the optimal curve: \texttt{ellmoddegree} returns the degree of the given model. The paper says that the curve must be
optimal with $E[\ell]$ irreducible, with the examples \curve{990h1} (returns \texttt{[4, 3]}) against the optimal \curve{990h3} (\texttt{[3, 3]}) and \curve{11a1} at $\ell=5$ (\texttt{[0, 1]}); the book
reran them, and \curve{11a3} at $\ell=5$ returns \texttt{[-1, 0]} (its degree is $1/5$ relative to the optimal curve). PARI's Kodaira codes: $n+4$ for $\kod{I}_n$ and $-(n+4)$ for $\kod{I}_n^*$ when $n\ge1$,
$1$ for $\kod{I}_0$ and $-1$ for $\kod{I}_0^*$, $\pm2$, $\pm3$, $\pm4$ for $\kod{II}/\kod{II}^*$, $\kod{III}/\kod{III}^*$, $\kod{IV}/\kod{IV}^*$; the qualifier ``when $n\ge1$'' matters, as $\kod{I}_0$ and $\kod{I}_0^*$ have the codes $\pm1$. The book's runs (PARI
2.17.2) of \texttt{rule} on the paper's and the book's example curves are in Section~\ref{sec:paper:verifications}; on \curve{726a1} both modes return \texttt{[2, 2]}, the type $\kod{III}$ term at $11$ and the
correction cancelling, as the paper says.

\paragraph{Reproducibility and the statement on AI.} The repository named in the appendix holds the scripts it names and the result files from which the six
data tables are generated; the book compared the printed tables with those files, and the pipeline was rerun in full from the pinned data in the book's
repository, reproducing every committed summary and table (\Cref{ch:computation}, in its section on reproducibility). The appendix says so: ``The whole pipeline was rerun from the pinned data with PARI/GP 2.17.2 and reproduced every committed result summary and Tables tab:uniform, tab:iso, tab:counts, tab:ns, tab:isolated and~tab:oot byte for byte.'' A
third stage of that run recomputed the numbers that the paper states without a committed result file; it led to two corrections of wording, in Section
sec:two ($19$ exceptions on the extension) and Section sec:eis (the $18$ curves with a point of order $9$, ``optimal or not''). The
book's own reproduction is in
\Cref{ch:computation} (Lemma lem:twist checked on twists of the book's range, Lemma lem:euler on $67{,}745$ pairs, \texttt{ellmoddegree} against Cremona's \texttt{degphi} on all $38{,}042$ optimal curves below
$10{,}000$, no mismatch). The statement on the use of an AI system is part of the paper; Section~\ref{sec:paper:reading} summarises it.

\section{The map of the paper}\label{sec:paper:map}

This section collects three devices that the section-by-section reading cannot show: which result feeds which (the dependency map), which hypothesis is used
where and what fails without it (the hypothesis table), and the status of every statement imported from the literature.

\paragraph{Dependencies.} An arrow means ``is used in the proof of''. Imported statements are in brackets.
{\small
\begin{longtable}{p{0.40\textwidth}p{0.52\textwidth}}
\toprule
input & used in\\
\midrule
\endhead
{[}Tate curves, Kummer theory] & Lemma lem:tamagawa-exp $\to$ Theorem thm:qual (c), Section 1.5, Remark rem:ko\\
{[}N\'eron models] $\to$ inertia invariants (\Cref{lem:paper:inertia-invariants}) & Lemma lem:IV at every $q\ne3$, $q=2$ included $\to$ Theorem thm:qual (a), (b), Proposition prop:forms, the table\\
{[}Tate's algorithm]; [Ogg-Saito] (at $q=3$) & Lemma lem:twist (i), (ii) $\to$ Lemma lem:euler, Theorem thm:qual (b), (c), Theorem thm:twist, Theorem thm:quant\\
{[}\cite[Theorem 3.1]{AtkinLi78}] & Lemma lem:twist (iii) $\to$ Theorem thm:qual (b), (c)\\
{[}\cite{Kraus90}; \cite{SerreTate68}] & Lemma lem:euler, Proposition prop:shapes $\to$ Definition def:terms values, Theorem thm:quant Step 3, Proposition prop:ns-structure\\
{[}\cite[I.4.2.2]{FPR94}] + group cohomology & Lemma lem:tamfactor, Lemma lem:tamadd $\to$ Theorem thm:main Step 7, Corollary cor:explicit, Remark rem:ko\\
{[}\cite{Ribet83,AU96}: the perfect pairing] & Lemma lem:crit $\to$ Lemma lem:old, Theorem thm:twist, Theorem thm:quant Step 1\\
{[}\cite{Edixhoven92,AtkinLehner70}; \cite[Theorem 3.1(e)]{DDT}; Chebotarev; Brauer-Nesbitt] & Lemma lem:old $\to$ Theorem thm:qual\\
{[}\cite[Theorem 3.15]{DDT}]; [\cite[Theorem 2.2]{ARS12}] & Theorem thm:qual\\
{[}\cite[Sections 1.2, 1.4]{DFG04}; Hochschild-Serre; equivariant Leray] $\to$ Proposition prop:integral; projection formula & Proposition prop:pullback $\to$ Lemma lem:hecke $\to$ Theorem thm:main Step 1, Theorem thm:quant Step 4\\
{[}\cite[Lemma 4.17]{DDT}] = Lemma lem:ddt417 & Theorem thm:quant Step 2 (and the second route to Lemma lem:hecke, Remark rem:free)\\
{[}\cite[Proposition 1.4(c)]{DFG04}] = Theorem thm:dfg14 & Theorem thm:main Step 2, Theorem thm:quant Step 3\\
{[}\cite[Theorem 3.7]{DFG04}] = Theorem thm:dfg37 & Theorem thm:main Step 3\\
{[}\cite[Theorem 2.7]{DFG04}] = Theorem thm:dfg27 & $H^1_f(\Q,V^D)=0$ in Lemma lem:dfg21 ($H^0(\Q,V^D)=0$ automatic)\\
{[}\cite[Lemma 2.1]{DFG04}] = Lemma lem:dfg21; Lemma lem:fpr & Theorem thm:main Steps 4, 5; Remark rem:ns (conditionally)\\
{[}Clifford theory; subgroups of $\GL_2(\F_3)$] & Proposition prop:ns $\leftrightarrow$ Lemma lem:h0 $\to$ Theorem thm:main Step 4, Conjecture conj:euler ($h_\ell$), Section 6.7\\
{[}inertia at $\ell$, Serre] & Proposition prop:Sf: $\ell\notin S_f$ automatic for $\ell\ge5$\\
Theorem thm:main & Proposition prop:forms, Corollary cor:explicit, Theorem thm:quant (ii)\\
Theorem thm:twist; the twist map $\to$ Lemma lem:step1 & Theorem thm:quant Step 1\\
{[}\cite{Ligozat75}: $X_0(27)$]; [\cite{Mazur77}: torsion] & Proposition prop:no9 $\to$ $D\le2$ in Conjecture conj:eis, given its first bound\\
data (Section 6) & Conjecture conj:euler ($w_3$, $h_\ell$), Conjecture conj:eis (first bound), Section 7\\
\bottomrule
\end{longtable}}
Three spines: local (Lemmas lem:tamagawa-exp, lem:IV, lem:twist, lem:euler, lem:tamfactor, lem:tamadd), congruences (Lemmas lem:crit, lem:old, the optimal-level theorem,
Theorem thm:qual) and Selmer (Propositions prop:integral, prop:pullback, Theorems thm:dfg14, thm:dfg37, thm:dfg27, Lemmas lem:dfg21, lem:fpr, lem:h0, Theorem thm:main). Theorem thm:quant uses
the first two and, from the third, only Theorem thm:dfg14, Lemma lem:ddt417 and Lemma lem:hecke; it never uses Theorem thm:dfg37, which is why its inequality needs no
hypothesis over $\Q(\sqrt{\ell^*})$, while its equality uses Theorem thm:main for $E$ and $E'_S$. Multiplicity one (Remark rem:free) feeds nothing.

\paragraph{Hypotheses.} Aligned with the table of \Cref{ch:congruence}.
{\small
\begin{longtable}{p{0.17\textwidth}p{0.36\textwidth}p{0.38\textwidth}}
\toprule
hypothesis & where it is used & what fails without it\\
\midrule
\endhead
$\ell$ odd & Lemma lem:tamfactor ($2$ a unit); trace form; real place; $2\in S_f$; Lemma lem:h0 (complex conjugation); optimal-level theorem; Lemma lem:old at $p=\ell$ &
everything in Sections 2 to 4 at $\ell=2$; Section 7 shows a different $2$-adic pattern\\
$\ell\nmid N$ ($\ell^2\nmid N$) & ARS12 (only $\ell^2\nmid N$); $S_f$ and Fontaine-Laffaille; Lemma lem:old at $p=\ell$; optimal-level theorem; the local terms need $q\ne\ell$ &
\curve{99a1} ($\degphi=4$, $r_E=12$ at $\ell=3$; $H^1(X_0(99),\Z_3)_{\mf m}$ not free); \curve{726a1} and \curve{539d1} at $11$, $7$; at $q=\ell=3$ the empirical weights $w_3$\\
$\ell\notin S_f$ ($\bar\rho$ absolutely irreducible over $\Q(\sqrt{\ell^*})$) & Theorems thm:dfg37, thm:dfg27; Lemma lem:h0 ($h_\ell=0$); Theorem thm:qual at $\ell=3$; the equality in Theorem thm:quant &
automatic for $\ell\ge5$, $\ell\nmid N$, $E[\ell]$ irreducible; at $\ell=3$ fails exactly for image $\Ns$, where (eq:thm) fails for $31\%$, $23\%$, $18\%$ of the curves, by one (\curve{726a1})\\
$E[\ell]$ irreducible & Theorem thm:dfg14; non-Eisenstein $\mf m$; Lemma lem:old; Theorems thm:qual, thm:quant & \curve{539a1} (image 3B, $e_3=1$, $v_3(\degphi)=0$); \curve{539d1} at $\ell=5$ ($e_5=1$, $v_5(\degphi)=0$); \curve{98a1}
($e_3=1$, $\degphi=16$); only (eq:main) survives, with the deficit of Conjecture conj:eis\\
optimality & ARS12; the definition of $\degphi$; Proposition prop:pullback; the function \texttt{rule} & nothing at $\ell$ when $E[\ell]$ is irreducible ($v_\ell$ of the degree is a class invariant); matters for
Conjecture conj:eis and for \texttt{rule} (\curve{990h1}: \texttt{[4, 3]})\\
multiplicity one & nowhere (Proposition prop:pullback and the squeeze in Theorem thm:quant (ii) replace it) & nothing; it holds in the paper's range (Remark rem:free) and
fails for $\ell\mid N$ (\curve{99a1} at $3$)\\
\bottomrule
\end{longtable}}

\paragraph{Imported statements.} ``Verified'' means compared with the published text (for \cite{DFG04}: Ann. Sci. \'Ecole Norm. Sup. (4) 37 (2004); a longer arXiv text, arXiv:2512.02348, numbers its
sections and results differently, and every number in the paper and in this book refers to the published version). ``As cited'' means taken from the paper and not checked for this book.
{\small
\begin{longtable}{p{0.27\textwidth}p{0.45\textwidth}p{0.19\textwidth}}
\toprule
citation & content used & status\\
\midrule
\endhead
\cite[Theorem 0.1]{DFG04} (= 2.15) & Bloch-Kato for $A_f$, $A_f(1)$ at $\lambda\notin S_f$ & verified\\
\cite[Theorem 3.7]{DFG04} (= 0.2) & $\len H^1_\Sigma(\Q,W)=v_\lambda(\eta^\Sigma)$ for $\ell\notin S_f$; covers $\ell=3$ in weight $2$ & verified\\
\cite[Theorem 2.7]{DFG04} & $H^0=H^1_f=0$ for $A_f$, $A_f(1)$ & verified\\
\cite[Proposition 1.4(c)]{DFG04} & $\eta^\Sigma=\eta^\emptyset\prod L^{\nv}_p(A_f,1)^{-1}$, with $\ell\notin\Sigma$ standing & verified\\
\cite[Lemma 2.1, (31), (57)]{DFG04} & the four-term sequence; $S_f$; $\Tam_p(M)=\Tam_p(M^D)$ & verified\\
\cite[Lemma 2.12, p.~708, Sections 1.1.1 to 2.1, (13), (14)]{DFG04} & naive Petersson formula; $\Tam$; Frobenius, $[UwU]_\omega$, kernel lattice, ideal of values & verified\\
\cite[I.4.2.2]{FPR94} & the local Tamagawa formula & as recalled in DFG04 (verified there)\\
\cite[Theorem 2.2]{ARS12} & $\degphi\mid r_E$, equality at $p^2\nmid N$ & checked in \Cref{ch:moddeg}\\
\cite[3.1(e), 3.15, 4.17, Section 4.4]{DDT} & local shape at $q\,\|\,N$; optimal level; numerical criterion; adjoints & as cited\\
\cite[Theorem 3.1]{AtkinLi78}; \cite[Theorem 3]{AtkinLehner70} & level of a twist; $a_p=0$ for $p^2\mid M$ & as cited\\
\cite[Theorem 2.2]{Ribet83}; \cite[Lemma 2.1]{AU96}; \cite[Section 2]{Edixhoven92} & perfect pairing; $a_\ell$ from $\bar\rho|_{D_\ell}$ & as cited (pairing proved in \Cref{ch:hecke})\\
\cite[Theorem 3.5]{RibetStein}; \cite{Tilouine97} & multiplicity one; Gorenstein & as cited\\
\cite[Theorem 1.1, Definition 2.12]{KimOta} & variation of congruence ideals; Tamagawa exponent & as cited\\
\cite[Theorem 6.11]{PollackWeston11} & their identity, under CR & checked (arXiv version)\\
\cite{RibetTakahashi97,Takahashi01} & Shimura-curve degrees, semistable $E$ & checked in \Cref{ch:moddeg}\\
\cite{Ligozat75}; \cite{Kenku} & rational points of $X_0(27)$; the list of levels & Kenku checked; Ligozat as cited\\
\cite{Kraus90}; \cite{SerreTate68} & inertia images at $q\ge5$; independence of $\ell$ & standard\\
\cite{Zywina15}; \cite{Sutherland16} & mod $3$ images; labels & by section; labels via Cremona's files\\
\cite{CalegariEmerton09} & odd modular degree & arXiv version read; published numbering unchecked\\
\bottomrule
\end{longtable}}

\section{Worked verifications}\label{sec:paper:verifications}

This section lists the curve-level numbers the tour quotes, as recomputed for this chapter with PARI/GP 2.17.2 through \texttt{cypari2} and Cremona's files: Kodaira codes from
\texttt{ellglobalred}, $P_q$ from \texttt{lfuneuler(lfunsympow(E,2),q)}, degrees from \texttt{ellmoddegree} (equal to Cremona's \texttt{degphi} on every optimal curve below $10{,}000$),
the outputs of the paper's \texttt{rule} (the code of Section~\ref{sec:paper:appendix}), and Tate parameters from \texttt{ellinit(E, O(q\^{}40)).tate}. Congruence numbers are those computed
in \Cref{ch:moddeg,ch:hecke}.

{\small
\begin{longtable}{p{0.13\textwidth}p{0.47\textwidth}p{0.31\textwidth}}
\toprule
curve & local data and $P_q$ & degree, \texttt{rule}, remarks\\
\midrule
\endhead
\curve{121a1} & $\kod{II}$ at $11$, $P_{11}=1+11X$, $e_3=v_3(12/11)=1$ & $6$; \texttt{[1, 1]} in both modes; $r_E=6$; twist \curve{121c2} ($\kod{IV}^*$)\\
\curve{539c1} & $\kod{I}_3^*$ at $7$, $P_7=1-X$, $e_3=1$, $t_3=1$; $\kod{I}_2$ at $11$; $\bar u=5$ at $7$ & $288=2^53^2$; \texttt{[2, 2]} at $3$, \texttt{[0, 0]} at $5$; twist \curve{77c1} ($\degphi=6$)\\
\curve{980i1} & $\kod{IV}^*$ at $2$ ($P_2=1+2X$), $\kod{I}_1$ at $5$, $\kod{I}_5^*$ at $7$ & $2880=2^63^25$; \texttt{[2, 2]} at $3$, \texttt{[1, 1]} at $5$; twist \curve{140b1} ($60$)\\
\curve{726a1} & $\kod{I}_6$ at $2$, $\kod{I}_3$ at $3$, $\kod{III}$ at $11$ ($P_{11}=1+11X$); $\Ns$, \texttt{ns3} $=1$ & $144$; \texttt{[2, 2]} in both modes; $r_E=1584$\\
\curve{99a1} & $\kod{III}$ at $3$ ($P_3=1+3X$, not used at $q=\ell$), $\kod{I}_1$ at $11$ & $4$; \texttt{[0, 0]}; $r_E=12$\\
\curve{121d1} & $\kod{I}_1^*$ at $11$, $e_3=v_3(120)=1$ & $24$; \texttt{[1, 1]}; twist of \curve{11a3}\\
\curve{121d2} & $\kod{I}_5^*$ at $11$, $e_3=1$ & $120$; twist of \curve{11a1} ($e_3(\curve{11a1},11)=0$)\\
\curve{539d1} & $\kod{I}_0^*$ at $7$, $P_7=1+3X-21X^2-343X^3$, $e_3=2$; $\kod{I}_1$ at $11$ & $72$; \texttt{[2, 2]}; $r_E=504$; at $5$: \texttt{[0, 1]} (5B.4.1)\\
\curve{539a1} & $\kod{I}_2^*$ at $7$, $e_3=1$; image 3B & $320$; \texttt{[0, 1]}: Eisenstein\\
\curve{98a1} & $\kod{I}_2$ at $2$, $\kod{I}_1^*$ at $7$, $e_3=1$; image 3B & $16$; \texttt{[0, 1]}: Eisenstein\\
\curve{52a1} & $\kod{IV}^*$ at $2$ ($P_2=1+2X$), $\kod{I}_2$ at $13$ & $3$; \texttt{[1, 1]}; $f_{52a}\equiv f_{26b}-f_{26b}(2z)\bmod3$\\
\curve{336f1} & $\kod{II}$ at $2$, $P_2=1+2X$, $e_3=1$ & $24$; \texttt{[1, 1]}\\
\curve{1232l1} & $\kod{II}^*$ at $2$, $P_2=1+2X-4X^2-8X^3$, $e_3=2$ & $288$; \texttt{[2, 2]}; $=\curve{77a1}\otimes\chi_{-4}$, $a_2(\curve{77a1})=0$\\
\curve{592c1} & $\kod{II}^*$ at $2$, $P_2=1-2X+4X^2-8X^3$, $e_5=1$ & $80$; \texttt{[1, 1]} at $5$; $=\curve{37a1}\otimes\chi_{-4}$, $a_2=-2$\\
\curve{791a1} & $\kod{I}_3$ at $7$, $\bar u=6$ a cube: Tamagawa excess $1$ & $144$; \texttt{[2, 1]}; explicit bound $2$, Selmer $0$\\
\curve{1184f1} & $\kod{III}^*$ at $2$ ($P_2=1$), $\kod{I}_3$ at $37$, $\bar u=36$ a cube; image 3Nn & $1728$; \texttt{[3, 1]}; explicit bound $2$, Selmer $1$\\
\curve{1748e1} & $\kod{IV}$ at $2$ ($e_3=1$), $\kod{I}_6$ at $19$ ($\bar u=12$ a cube), $\kod{I}_1$ at $23$ & $3240$; \texttt{[4, 2]}; explicit bound $3$, Selmer $1$\\
\curve{11a1} & $\kod{I}_5$ at $11$, torsion $\Z/5$, image 5Cs.1.1, $\bar u=10$ & $1$; \texttt{[0, 1]} at $5$: Eisenstein\\
\curve{14a1} & $\kod{I}_6$ at $2$, $\kod{I}_3$ at $7$, torsion $\Z/6$, 3Cs.1.1 & $1$; $D=2$\\
\curve{714i1} & $\kod{I}_9$ at $2$ and $3$, $\kod{I}_3$ at $7$, $\kod{I}_1$ at $17$, torsion $\Z/9$, 3B.1.1 & $1080$; $D=5-3=2$\\
\curve{20a1} & $\kod{IV}^*$ at $2$, $\kod{I}_2$ at $5$, torsion $\Z/6$, 3B.1.1 & $1$; $D=1$\\
\curve{990h1}, \curve{990h3} & the same class; \curve{990h3} is optimal & \texttt{[4, 3]} against \texttt{[3, 3]}\\
\bottomrule
\end{longtable}}

The twist identifications (\curve{11a1} by $-11$, $-7$ to \curve{121d2}, \curve{539d2}; \curve{20a1} by $-7$ to \curve{980g2}; \curve{539c1} by $-7$ to \curve{77c1}; \curve{121a1} by $-11$ to \curve{121c2}; \curve{980i1}
by $-7$ to \curve{140b1}; \curve{37a1}, \curve{77a1} by $-4$ to \curve{592c1}, \curve{1232l1}) were checked by isomorphism of minimal models; the congruences of \Cref{ex:paper:qual} to $n\le400$ and
$n\le1000$; \texttt{mfsturm} returns $14$, $22$, $112$, $336$. Also recomputed: $j=-2^{15}\cdot3\cdot5^3=-12288000$ for discriminant $-27$; the three coordinate vectors and the determinant $-1$ in the
proof of \Cref{lem:paper:tamadd}; $L^{\nv}_3(A_f,1)^{-1}=10/9$ for $f_{11a}$ (\Cref{thm:paper:dfg14}); the $\Ns$ counts of Section 5 and Section 6.7 of the paper ($90$ non-CM, $44$ with $3\mid N$, $74$
and $16$, $13$ CM; $36$ and $24$ violations).

\section{Exercises}\label{sec:paper:exercises}

Data for the curves are in Section~\ref{sec:paper:verifications}; computations by hand, or with PARI, are allowed, and the hints say which.

\begin{exercise}\label{ex:paper:121a1}
The curve \curve{121a1} is $y^2+xy+y=x^3+x^2-30x-76$, with $\Delta=-11^2$ and $c_4=1441=11\cdot131$. (a) Show that it has additive, potentially good reduction of type $\kod{II}$ at $11$. (b) Compute
$e_3(E,11)$ and $e_5(E,11)$ from the table. (c) Using $\degphi=6$, \Cref{cor:paper:explicit} and \Cref{lem:paper:tamadd}, determine $\len H^1_f(\Q,\ad T_3E\otimes\Q_3/\Z_3)$ (the mod $3$ image is
surjective).
\emph{Hint.} (a) $v_{11}(j)=3v_{11}(c_4)-v_{11}(\Delta)=1\ge0$ and $v_{11}(\Delta)=2$. (b) $\chi_{-3}(11)=-1$: $e_3=v_3(12)=1$, $e_5=0$. (c) $1=\len H^1_f+1+0$, so the Selmer length is $0$.
\end{exercise}

\begin{exercise}\label{ex:paper:twist-e}
(a) The curve \curve{37a1} has type $\kod{I}_1$ at $37$ and $\degphi=2$; its twist by $\chi_{37}$ is \curve{1369d1}, of type $\kod{I}_1^*$ at $37$, alone in its class. Compute $e_3$ and
$e_{19}$ at $37$ for both curves, explain why they differ although $\ad V$ does not change, and use \Cref{thm:paper:quant}(ii) to predict $v_3$ and $v_{19}$ of $\deg\varphi_{\curve{1369d1}}$; check
against $\deg\varphi_{\curve{1369d1}}=2736$. (b) For the class \curve{539d} (type $\kod{I}_0^*$ at $7$, the twist of \curve{11a} by $\chi_{-7}$) compute $e_3$ and $e_5$ at $7$ from $a_7(\curve{11a1})=-2$,
compare with $\degphi=72$, and explain the behaviour at $\ell=5$.
\emph{Hint.} (a) For \curve{37a1} the naive and true factors are both $1-37^{-2}$, so $e_\ell=0$; for \curve{1369d1} the naive factor is $1$ and $e_\ell=v_\ell(37^2-1)=v_\ell(2^3\cdot3^2\cdot19)$: $e_3=2$,
$e_{19}=1$. The hypotheses of (ii) hold for $\ell=3,19$ (the mod $\ell$ images of \curve{37a1} are surjective, $\ell\nmid2\cdot1369$, and $\ell\notin S_f$ by \Cref{prop:paper:Sf}, the image at $3$ not being
$\Ns$), with $S=\{37\}$ and $E'_S=\curve{37a1}$; so $v_3=0+2$ and $v_{19}=0+1$, as $2736=2^4\cdot3^2\cdot19$. (b) $e_\ell=v_\ell(6)+v_\ell(64-4)$: $e_3=2=v_3(72)$, $e_5=1$ but $v_5(72)=0$; the
mod $5$ image of \curve{539d1} is reducible (5B.4.1), outside Theorem thm:quant.
\end{exercise}

\begin{exercise}\label{ex:paper:dfg14-ell}
Let $f=f_{37a}$, with $a_3=-3$, and $\ell=3$. Compute $L^{\nv}_3(A_f,1)^{-1}$ and show that the formula of \Cref{thm:paper:dfg14} with $\Sigma=\{3\}$ cannot hold. Which ingredient of the theory of
\cite{DFG04} needs $\ell\notin\Sigma$?
\emph{Answer.} $\frac23\cdot\frac{16-9}9=\frac{14}{27}$, of $3$-adic valuation $-3$, so the right side is not in $\Z_3$. Adding $\ell$ to the level destroys the condition $\ell\nmid$ level under which the integral
structures and the Fontaine-Laffaille description at $\ell$ are defined (\Cref{ch:congruence}).
\end{exercise}

\begin{exercise}\label{ex:paper:tamfactor-coordinates}
Lemma lem:tamfactor in coordinates. Let $\ell=3$, $n=9$, $q=19$ ($v_3(q-1)=2$), and let $\Frob$ act on $(\Z/9)\bar e\oplus(\Z/9)\bar h$ by $\bar e\mapsto\bar e$, $\bar h\mapsto q^{-1}(\bar h-2c'\eps\bar e)$, $\eps=\pm1$. (a) For
each $c'\in\Z/9$ count the fixed points and check that the count is $9\cdot3^{\min(2,2,v_3(c'))}$. (b) Replace the action by $q\cdot\Frob$ (the twist by $\chi_3(\Frob)=q$, which is the action on $H^1(I,A(1))$)
and show that the count is unchanged, so $\Tam_q(A(1))=\Tam_q(A)$ here. (c) Explain why (b) holds for every $n$, $q$, $\ell$.
\emph{Hint.} (a) Fixed iff $(q-1)y\equiv0$ and $c'y\equiv0$. (b) $x\bar e+y\bar h$ is fixed by $q\Frob$ iff $(q-1)x\equiv2c'\eps y$; here $q-1\equiv0\pmod9$, so the condition is $c'y\equiv0$ with $x$ free. (c)
In general the two conditions are $(q-1)y\equiv0,\ c'y\equiv0$ and $(q-1)x\equiv2c'\eps y$; count the solutions of the second system by first choosing $y$ and compare.
\end{exercise}

\begin{exercise}\label{ex:paper:corner}
Let $G$ be a group, $\chi,\eps\colon G\to\Zl^\times$ characters with $\eps^2=1$, and $c\colon G\to\Zl$ a function. (a) Show that $g\mapsto\sm{\chi\eps(g)}{c(g)}0{\eps(g)}$ is a homomorphism iff
$c(gh)=\chi\eps(g)c(h)+c(g)\eps(h)$. (b) If $\kappa$ satisfies $\kappa(gh)=\chi(g)\kappa(h)+\kappa(g)$, show that $c=\eps\kappa$ satisfies (a), and that $c=\kappa$ does not unless $(\eps(g)-1)\chi(g)\kappa(h)+(\eps(h)-1)\kappa(g)=0$
for all $g,h$. (c) Deduce the description of the nonsplit Tate curve used in \Cref{lem:paper:tamfactor}.
\emph{Hint.} (a) The diagonal entries of the product of the matrices for $g$ and $h$ are $\chi\eps(gh)$ and $\eps(gh)$; the upper right entry is $\chi\eps(g)c(h)+c(g)\eps(h)$. (b) Multiply the
cocycle relation for $\kappa$ by $\eps(gh)=\eps(g)\eps(h)$. For $c=\kappa$, subtract the relation for $\kappa$ from the one required in (a). (c) The nonsplit curve is the Tate curve twisted by the
unramified quadratic character $\eps$; its representation is $\sm{\chi_\ell}{\kappa}01\otimes\eps$.
\end{exercise}

\begin{exercise}\label{ex:paper:ddt-toy}
Let $\mc O=\Zl$, $\T=\{(a,b)\in\Zl^2:a\equiv b\pmod{\ell^r}\}$ with $r\ge1$, and $\pi(a,b)=a$. (a) Compute $\wp$, $I=\Ann_\T(\wp)$ and $\eta$, and find $t$ with $I=\Zl t$. (b) Let $\T$ act on $L'=\Zl^2$ by componentwise
multiplication, with the pairing $\langle(a,b),(c,d)\rangle=ac+bd$. Check the hypotheses of \Cref{lem:paper:ddt417}, compute $L'[\wp]$ and $\Delta$, and show that the inequality is strict; is $L'$ free over $\T$?
(c) Let $L=\T$ with $\langle(a,b),(c,d)\rangle=(ac-bd)/\ell^r$. Show that this pairing is integral, perfect and $\T$-self-adjoint, and that \Cref{lem:paper:ddt417} holds with equality.
\emph{Answer.} (a) $\wp=0\times\ell^r\Zl$, $I=\ell^r\Zl\times0$, $\eta=\ell^r\Zl$, $t=(\ell^r,0)$. (b) $L'[\wp]=\Zl(1,0)$, $\Delta=1$, $v(\eta)=r>0$; $L'$ is not free over $\T$ (it is the normalisation of $\T$). (c) In the
basis $(1,1)$, $(\ell^r,0)$ of $\T$ the Gram matrix is $\sm011{\ell^r}$, of determinant $-1$; $L[\wp]=\Zl(\ell^r,0)$ has $\Delta=\ell^r$, and $v(\eta)=r=v(\Delta)$.
\end{exercise}

\begin{exercise}\label{ex:paper:gram}
(a) Show that an alternating form on $\mc O^2$ has Gram matrix $\sm0a{-a}0$ in every basis, so its Gram determinant is $a^2$ and its ideal of values is $a\mc O$. (b) Deduce from the projection formula
that the Gram determinant of $\frac12\langle x,w_N^*y\rangle$ on $\varphi_E^*H^1(E,\Zl)$ is $(\degphi)^2$ up to a unit, for $\ell$ odd. (c) For \curve{11a1}, where $X_0(11)\to E$ has degree $1$, check this directly.
\emph{Hint.} (a) $\langle x,x\rangle=0$ for all $x$ forces zero diagonal entries and $\langle e_2,e_1\rangle=-\langle e_1,e_2\rangle$; a change of basis by $P$ multiplies $a$ by $\det P$, a unit; and
$\langle x,y\rangle=a\det(x,y)$ takes every value in $a\mc O$. (b) On the image $w_N^*$ acts by $\varepsilon_N(f)=\pm1$, so the form is $\frac{\varepsilon_N(f)}2\degphi$ times the pullback of the
unimodular cup product of $E$; it is alternating, and (a) applies with $a=\pm\frac12\degphi$. (c) The parametrisation is an isomorphism, so $\varphi_E^*$ is an isomorphism of lattices carrying one
unimodular cup product to the other; the Gram determinant is $1=1^2$.
\end{exercise}

\begin{exercise}\label{ex:paper:726a1}
The curve \curve{3840h1} has types $\kod{III}$ at $2$ (with $P_2=1+2X$), $\kod{I}_3$ at $3$ and $\kod{I}_6$ at $5$, mod $3$ image $\Ns$, and $\degphi=1152=2^7\cdot3^2$. (a) Compute the right side of
(eq:thm) at $\ell=3$ as written, and the same without the terms at $q=3$; which prime $p\equiv-1\pmod3$ carries the failure, and by how much? (b) Compute \texttt{rule(E, 3)} and
\texttt{rule(E, 3, 1)} by hand. (c) Identify the order-four prime of \Cref{prop:paper:ns-structure} and check part (ii) there. (d) Why is $w_3(E)=0$?
\emph{Answer.} (a) $e_3(E,2)=v_3(P_2(1/4))=v_3(3/2)=1$, $t_3(E,3)=1$, $t_3(E,5)=v_3(6)=1$: the bound is $3>2$, a failure by one, carried by $p=5$ with $t_3(E,5)=1$; without $t_3(E,3)$ the
bound is $2$. (b) Without \texttt{iii} the type $\kod{III}$ term at $2$ is left out: $0+1+1=2$; with it, $1+1+1-\texttt{ns3}(E)=2$; both return \texttt{[2, 2]}. (c) The prime $2$, of type $\kod{III}$: $5$
is potentially multiplicative, so $\bar\rho(I_5)$ consists of $\pm$ unipotents, while $R\ne\emptyset$; by (ii), $P_2=1-\chi_{-3}(2)2X=1+2X$ and $e_3(E,2)=1$, as computed. (d) The reduction at $3$ is
multiplicative.
\end{exercise}

\begin{exercise}\label{ex:paper:rankin}
Let $p$ be a good prime, $a_{p^k}$ the coefficients of $f$, $\alpha+\beta=a_p$, $\alpha\beta=p$, $X=p^{-s}$. (a) Using $a_{p^k}=(\alpha^{k+1}-\beta^{k+1})/(\alpha-\beta)$, prove
$\sum_ka_{p^k}^2X^k=(1+pX)/((1-\alpha^2X)(1-pX)(1-\beta^2X))$ and $\sum_ka_{p^{2k}}X^k=(1+pX)/((1-\alpha^2X)(1-\beta^2X))$. (b) Deduce $\sum_na_n^2n^{-s}=\zeta^{(N)}(s-1)L^{\nv}(\Sym f,s)/\zeta^{(N)}(2s-2)$, checking the
multiplicative and additive primes separately.
\emph{Hint.} (a) $(\alpha-\beta)^2a_{p^k}^2=\alpha^{2k+2}-2p^{k+1}+\beta^{2k+2}$; summing geometric series,
$(\alpha-\beta)^2\sum_ka_{p^k}^2X^k=\frac{\alpha^2}{1-\alpha^2X}-\frac{2p}{1-pX}+\frac{\beta^2}{1-\beta^2X}$; put everything over the common denominator and use $\alpha\beta=p$. For the second series,
$(\alpha-\beta)a_{p^{2k}}=\alpha^{2k+1}-\beta^{2k+1}$ and $\alpha(1-\beta^2X)-\beta(1-\alpha^2X)=(\alpha-\beta)(1+pX)$. (If $\alpha=\beta$, argue by continuity in the coefficients.) (b) At a good prime the
local factors of $\zeta^{(N)}(s-1)$ and $\zeta^{(N)}(2s-2)$ are $(1-pX)^{-1}$ and $(1-p^2X^2)^{-1}$; at $p\mid N$ they are $1$, and $a_{p^k}=a_p^k$ with $a_p\in\{0,\pm1\}$.
\end{exercise}

\begin{exercise}\label{ex:paper:deficits}
Compute the deficit $D$ of Conjecture conj:eis at $\ell=3$ for the optimal curves \curve{26a1} (types $\kod{I}_3$ at $2$ and $13$, torsion $\Z/3$, image 3Cs.1.1, $\degphi=2$), \curve{30a1}
($\kod{I}_4$ at $2$, $\kod{I}_3$ at $3$, $\kod{I}_1$ at $5$, torsion $\Z/6$, image 3B.1.1, $\degphi=2$) and \curve{6942n1} ($\kod{I}_{18}$ at $2$, $\kod{I}_9$ at $3$, $\kod{I}_3$ at $13$, $\kod{I}_1$ at $89$,
torsion $\Z/9$, image 3B.1.1, $\degphi=65664$), and say which summand of the first bound is responsible in each case. Is \curve{6942n1} consistent with \Cref{prop:paper:no9}?
\emph{Answer.} \curve{26a1}: $D=1+1-0=2=v_3(3)+[E[3]\simeq\Z/3\oplus\mu_3]$, both summands. \curve{30a1}: $D=1-0=1=v_3(6)$, the torsion. \curve{6942n1}: $65664=2^7\cdot3^3\cdot19$, so
$D=(2+2+1+0)-3=2=v_3(9)$, the torsion. Yes: its image 3B.1.1 means $E[3]\not\simeq\Z/3\oplus\mu_3$, and the proposition excludes a point of order $9$ only when $E[3]\simeq\Z/3\oplus\mu_3$.
\end{exercise}

\begin{exercise}\label{ex:paper:ns-ap}
Let $E$ have mod $3$ image $N_s(\F_3)$. (a) Show that $a_p\equiv0\pmod3$ for every prime $p\nmid3N$ with $p\equiv2\pmod3$. (b) Check it on \curve{726a1} for $p=5,17,23,29,41$ ($a_p=0,6,6,-6,-6$). (c) Show
that (a) also holds when the image is contained in the split Cartan subgroup $C_s(\F_3)$ (then $E[3]$ is reducible), but not for the image $N_{ns}(\F_3)$.
\emph{Hint.} (a) $\det\bar\rho(\Frob_p)=\chi_{-3}(p)=-1$ puts $\bar\rho(\Frob_p)$ in the coset of determinant $-1$, whose elements $\pm h$, $\pm w$ have trace $0$. (c) The elements of determinant $-1$
of $C_s(\F_3)$ are $\pm h$; those of the nonsplit Cartan subgroup $\F_9^\times\subset N_{ns}(\F_3)$ have order $8$ and nonzero trace (proof of \Cref{prop:paper:ns}), and by Chebotarev they occur as
$\bar\rho(\Frob_p)$.
\end{exercise}

\begin{exercise}\label{ex:paper:wild-two}
The curves \curve{848e1} and \curve{1072a1} are the twists by $\chi_{-4}$ of \curve{53a1} and \curve{67a1}, with $a_2(\curve{53a1})=-1$ and $a_2(\curve{67a1})=2$. PARI gives \curve{848e1} the symbol
$\kod{I}_4^*$ at $2$, with $v_2(j)=0$, and \curve{1072a1} the type $\kod{II}^*$ at $2$. (a) Using \Cref{prop:paper:shapes}(ii), derive $P_2$ for both curves. (b) Show that $e_\ell(\curve{848e1},2)=0$
for every odd $\ell$ and $t_\ell(\curve{848e1},2)=0$, although the symbol at $2$ is $\kod{I}_4^*$. (c) Compute $e_5(\curve{1072a1},2)$ and compare with $\degphi=200$.
\emph{Answer.} (a) $(1-2X)(1-(a^2-4)X+4X^2)$: for \curve{848e1}, $a^2=1$, $P_2=(1-2X)(1+3X+4X^2)$; for \curve{1072a1}, $a^2=4$, $P_2=(1-2X)(1+4X^2)$. (b) $P_2(1/4)=\frac12(1+\frac34+\frac14)=1$, so
$e_\ell=0$; $v_2(j)\ge0$, so the reduction is potentially good and $t_\ell=0$: at $q=2$ a symbol $\kod{I}_m^*$ does not imply potentially multiplicative reduction. (c) $P_2(1/4)=\frac12\cdot\frac54=\frac58$,
$e_5=1$, while $v_5(200)=2$: the bound is not attained (see \Cref{ex:paper:selmer-examples}).
\end{exercise}

\begin{exercise}\label{ex:paper:selmer-examples}
Using \Cref{cor:paper:explicit}, compute the explicit bound and the Selmer length $\len H^1_f(\Q,\ad T\otimes\Qlmod)$: (a) for \curve{1577a1} at $\ell=3$ and at $\ell=5$ (types $\kod{I}_3$ at $19$ with
$\bar u=8$, and $\kod{I}_2$ at $83$; $\degphi=720$; mod $\ell$ images surjective for every $\ell$); (b) for \curve{1072a1} and \curve{67a1} at $\ell=5$ (\Cref{ex:paper:wild-two}; \curve{67a1} has type
$\kod{I}_1$ at $67$ and $\degphi=5$; both have surjective mod $5$ image). Explain why the two Selmer lengths in (b) must agree.
\emph{Answer.} (a) At $\ell=3$: $t_3(19)=1$, $v_3(18)=2$, and $8$ is a cube but not a ninth power in $\F_{19}^\times$, so the Tamagawa excess is $\min(1,2,1)=1$; $t_3(83)=0$; bound $2=v_3(720)$,
Selmer length $0$. At $\ell=5$: $t_5=0$ at both primes and $v_5(18)=v_5(82)=0$; bound $0$, $v_5(720)=1$, Selmer length $1$. (b) \curve{1072a1}: bound $e_5(2)=1$, $v_5(200)=2$, Selmer length $1$;
\curve{67a1}: bound $0$, $v_5(5)=1$, Selmer length $1$. The modules $\ad T_5E$ of a curve and of its quadratic twist are isomorphic as $G_\Q$-modules, and the Bloch-Kato conditions depend only on
the module, so the Selmer groups coincide; the degrees differ by the omitted factor $e_5(\curve{1072a1},2)$.
\end{exercise}

\section{Notes and references}\label{sec:paper:notes}

\paragraph{What was read.} The paper in full, as it accompanies this book, including the printed rows of its tables; the entries quoted in
Section~\ref{sec:paper:data} are copied from them. The repository named in its appendix was read for the result files behind the six data tables
(\Cref{ch:computation}, in its section on reproducibility); its large-range computations were rerun in full by its own pipeline in the book's
repository, not by the book's scripts.

\paragraph{Sources for the imported statements.} The theorems of Diamond, Flach and Guo are cited from the published version \cite{DFG04}; their numbers were compared with it (\Cref{ch:congruence}),
and a later arXiv version (arXiv:2512.02348) numbers them differently. Agashe-Ribet-Stein \cite{ARS12}, Ribet-Takahashi \cite{RibetTakahashi97} and Takahashi \cite{Takahashi01}, Pollack-Weston
\cite{PollackWeston11} (arXiv version) and Kenku \cite{Kenku} were checked for the book (\Cref{ch:moddeg,ch:modgalrep,ch:ellcurves}); the other theorem numbers are as cited by the paper, and the list in
Section~\ref{sec:paper:map} says which. The standard references behind the local computations are \cite{SilvermanATAEC} (Tate's algorithm, Tate curves), \cite{SerreTate68}, \cite{Kraus90} and \cite{Tate75};
behind the Hecke side \cite{DDT}, \cite{AtkinLehner70}, \cite{AtkinLi78}, \cite{Ribet83}, \cite{AU96}; behind the Selmer side \cite{FPR94} and \cite{BlochKato}; for the $L$-functions \cite{Shimura75,Shimura76,Rankin39}.
For the Eisenstein case and the prime $2$: \cite{Mazur77,Watkins02,CalegariEmerton09}.

\paragraph{Computations.} Every curve-level number printed in this chapter was recomputed with PARI/GP \cite{PARI} 2.17.2 through \texttt{cypari2}, on Cremona's files \cite{ecdata} (commit \texttt{25cec5e},
conductor below $10{,}000$); the scripts are in the style of \Cref{ch:computation}. The population counts attributed to the book (the $\Ns$ counts, the $36$ and $24$ violations, the class-invariance
of $w_3$) were recomputed for this chapter; the others are taken from the chapters cited at each place. None of these computations is a blind test: they were made with the statements known.

\paragraph{The paper and the book.} The paper's statement on the use of an AI system says that several of its results and proofs were worked out while writing a companion textbook on the
note; among them are Propositions prop:pullback, prop:Sf, prop:shapes, prop:ns-structure and prop:no9, Lemmas lem:tamadd and lem:step1, Corollary cor:explicit, the equality in Theorem
thm:quant and the exchanges at $q=3$ in Lemma lem:twist. In this book they are proved in the chapters named at each result; this chapter only reads them in the paper's order.
\end{sloppypar}
